% Generated by scripts/integrate_plethysm.py; edit original branch files.
\section{Conventions, source locators and status}
The source is Mulmuley, \emph{Geometric complexity theory VI:
the flip via positivity}, arXiv:0704.0229v4 (GCT VI).
For a global quasipolynomial $f$, with constituent polynomials
$q_a(X)$ at its least period, $p(f)$ is the least nonnegative
integer $u$ for which every polynomial $q_a(X+u)$ has
nonnegative coefficients.  The index $s(f)$ is the least
nonnegative integer $u$ for which every nonzero constituent
$q_a(X+u)$ is positive at every positive integer $X$.
These are Definitions 1.2.2 and 1.2.4, printed pp.~13--14.
In particular the positivity test is not restricted to the
residue class indexing a constituent.

The positive-form convention is Definition 1.2.7, p.~15;
the plethysm hypotheses are 1.6.5--1.6.7, pp.~24--25.
The general-form hypotheses are 3.3.1--3.3.4, pp.~50--51;
the subgroup specifications and costs are in Sections
3.4--3.4.2, pp.~52--57.  Their rank parameter differs from
$h=\dim G+\operatorname{ht}(\lambda)+\operatorname{ht}(\pi)$
in the narrow plethysm hypotheses.  The canonical Cartan algebra
is defined in Section 4.1.2, equations (4.2)--(4.5), pp.~67--69.

The geometric sections below reproduce the full proofs of
ALPOGE\_GCT P-F19, P-F23, P-F24, P-F26, P-F27, P-F28,
P-F30 and P-F31, followed by the needed P-F56 Fourier and
recurrence proofs.  These are inherited results, not newly
claimed discoveries.  Each proof needed for this chain has
been read in full in the present branch, including the inputs
previously carried only at statement level.  Exact provenance
and hashes are in \texttt{source\_manifest.json}.
Classical background theorems (finite generation of reductive
invariants, the affine quotient and Luna slice theorem,
smooth-hypersurface stability, and global quasipolynomiality)
are invoked with precise primary locators rather than reproved.
The obstructing support, noncancellation, residue, wave and
quantitative arguments themselves are included in full.

The new numerical estimates PNEW-002--005 have received two
independent whole-proof readings, first at the conservative
threshold and then at $79,000$.  This is mathematical review,
not formal proof-assistant verification.  The exact finite
certificates use integer and rational arithmetic only.

\section{The canonical cubic ring and the fixed-locus proof chain}


% Source: branches/plethysm/sections/positivity_cubic_counterexample.tex
% Reproduced inherited proof; finding P-F19.
\subsection{An exponential canonical modular index inside ordinary plethysm}

\begin{proposition}[P-F19: chain cubics obstruct the canonical clause of PH3]
\label{prop:cubiccanonicalcounterexample}
Let $n\ge2$, put
\[
 g=\gcd(n,3),\qquad d_0=\frac ng,\qquad
 W=\operatorname{Sym}^3\mathbb C^n,\qquad
 m=\dim W=\binom{n+2}{3},
\]
and consider the ordinary plethysm ray with
\[
 \mu=(3),\qquad \lambda=(d_0),\qquad
 \pi=((3/g)^n).
 \tag{P-F19a}
\]
Its canonical multiplicity algebra is
\[
 T^{(n)}=\bigoplus_{s\ge0}
  \operatorname{Hom}_{GL_n}
  \left((\det)^{3s/g},\operatorname{Sym}^{d_0s}W\right),
 \tag{P-F19b}
\]
with the Cartan multiplication. Every homogeneous system of parameters
of $T^{(n)}$ contains a parameter of stretch degree at least
$2^{n-1}$. In particular its canonical modular index satisfies
\[
 \chi(T^{(n)})\ge2^{n-1}.
 \tag{P-F19c}
\]
The parameter of GCT VI, Hypothesis 1.6.7, is
\[
 h=m^2+\operatorname{ht}(\lambda)+\operatorname{ht}(\pi)
   =\binom{n+2}{3}^{\!2}+n+1=O(n^6).
 \tag{P-F19d}
\]
Consequently the canonical modular-index clause of that hypothesis
is false already for ordinary plethysm with the fixed inner shape
$\mu=(3)$. This statement concerns the h.s.o.p. of the actual canonical algebra.
The parameter argument below alone does not assert that every positive
rational expression for its Hilbert series has a large denominator
weight. P-F31 later proves that stronger conclusion using the full
support theorem and an actual Hilbert-series pole. The superpolynomial positivity and saturation indices
require the paired-wave argument and quantitative estimates below.
\end{proposition}

\begin{proof}
Write $V=\mathbb C^n$ and identify
\[
 R=\mathbb C[W^*]^{SL_n}=(\operatorname{Sym}W)^{SL_n}.
 \tag{P-F19f}
\]
Thus points of $W^*=\operatorname{Sym}^3V^*$ are cubic forms in
the linear coordinate functions $x_1,\ldots,x_n$ on $V$.
This choice avoids identifying a representation with its dual.
The usual action on such a form is $(a\cdot F)(x)=F(a^{-1}x)$.
The ring $R$ has its ordinary coefficient grading, denoted $R_D$.

First we check the canonical grading and multiplication in (P-F19b).
If $R_D\ne0$, a scalar $\zeta I_n\in SL_n$ with $\zeta^n=1$
acts on $\operatorname{Sym}^D W$ by $\zeta^{3D}$, whence
$n\mid3D$. This is equivalent to $D=d_0s$ for some integer
$s\ge0$. On $R_{d_0s}$ the entire $GL_n$ acts by the character
$(\det)^{3s/g}$. Indeed, write $a=zh$ with $h\in SL_n$ and
$z^n=\det a$. Since $h$ fixes this subspace and $zI_n$ acts by
$z^{3d_0s}$, its action is $(\det a)^{3s/g}$, independently of
the choice of the $n$th root $z$.

Choose a nonzero volume vector $\varepsilon\in\bigwedge^nV$.
Evaluation of an intertwiner at $\varepsilon^{\otimes3s/g}$
identifies the degree-$s$ space in (P-F19b) with $R_{d_0s}$.
The Cartan product of these one-dimensional determinant modules
sends the compatible volume vectors to their tensor product.
Evaluation consequently respects multiplication in
$\operatorname{Sym}W$. We obtain an isomorphism of graded
algebras
\[
 T^{(n)}\simeq\bigoplus_{s\ge0}R_{d_0s},
 \qquad \deg_{T^{(n)}} R_{d_0s}=s.
 \tag{P-F19g}
\]
All nonzero ordinary degrees of $R$ occur in this display, so this
is the whole invariant ring with its grading divided by $d_0$,
not a proper Veronese subring. The canonical ring and its
multiplication are those in GCT VI, Section 4.1.2, equation (4.5),
printed p.~69, and equations (4.2)--(4.3), printed p.~67.
For the canonical modular index we use the detailed $T$-ring
definition of Section 4.1.1, printed p.~68.

Consider the explicit point
\[
 F_n=x_1^2x_2+x_2^2x_3+\cdots+x_{n-1}^2x_n+x_n^3
       \in W^*.
 \tag{P-F19h}
\]
Its projective hypersurface is smooth. To see this directly, its
partial derivatives are
\[
 \partial_{x_1}F_n=2x_1x_2,\qquad
 \partial_{x_i}F_n=x_{i-1}^2+2x_ix_{i+1}\ (1<i<n),\qquad
 \partial_{x_n}F_n=x_{n-1}^2+3x_n^2.
\]
Suppose they all vanish at a nonzero vector, and let $j$ be the
smallest index with $x_j\ne0$. If $j<n$, the $j$th derivative
gives $x_{j+1}=0$, and the $(j+1)$st derivative is then $x_j^2$,
which is nonzero, a contradiction. If $j=n$, the last derivative
is $3x_n^2\ne0$, also a contradiction. Thus the partials have
no common nonzero zero, proving smoothness over $\mathbb C$.

Let $\Delta$ be the discriminant of cubic forms. It is an
$SL_n$-invariant homogeneous polynomial of coefficient degree
$n2^{n-1}$, and it is nonzero precisely on smooth projective
hypersurfaces. An exact primary locator for these properties is
McCrory--Shifrin--Varley, \emph{Siegel modular forms generated by
invariants of cubic hypersurfaces}, arXiv:alg-geom/9403005v3,
Section 3.1.2(i), printed p.~15: the discriminant is the resultant
of the partial derivatives, with degree $n(k-1)^{n-1}$ for forms
of degree $k$.
For a direct resultant treatment see also Bus\'e--Jouanolou,
\emph{A computational approach to the discriminant of homogeneous
polynomials}, arXiv:1210.4697v1, Definition 4.6 and
Proposition 4.7(i), printed p.~40, and Proposition 4.13,
printed pp.~42--43, for the coordinate-change character.
Hence
\[
 \Delta(F_n)\ne0,\qquad
 \Delta\in R_{n2^{n-1}}
      =T^{(n)}_{g2^{n-1}}.
 \tag{P-F19i}
\]
In particular $F_n$ is outside the null cone. No claim about
stability beyond this nonvanishing is needed.

We next prove the decisive divisibility statement:
\[
 P\in T^{(n)}_s,\quad P(F_n)\ne0
       \quad\Longrightarrow\quad 2^{n-1}\mid s.
 \tag{P-F19j}
\]
We give two independent calculations.

\emph{Coefficient-weight calculation.}
Write $D=d_0s$ for the ordinary coefficient degree. A polynomial
invariant is a sum of coefficient monomials of torus weight zero:
the diagonal torus acts diagonally on the monomial basis, so a
nonzero weight cannot occur in an invariant polynomial.
If $P(F_n)\ne0$, at least one monomial occurring in $P$ uses
only the $n$ coefficients displayed in (P-F19h), because every
other coefficient is zero at $F_n$. Let its nonnegative exponent
tuple be $(a_1,\ldots,a_n)$; then $\sum_i a_i=D$.
Triviality of its character on the determinant-one diagonal torus
means its exponent vector in $x_1,\ldots,x_n$ is constant.
Thus, for one integer $c$,
\[
 2a_1=c,\qquad
 a_{i-1}+2a_i=c\quad(2\le i<n),\qquad
 a_{n-1}+3a_n=c.
 \tag{P-F19k}
\]
Adding these equations gives $nc=3D$, so $c=3s/g$.
The first $n-1$ equations solve to
\[
 a_i=\frac{c\bigl(2^i-(-1)^i\bigr)}{3\,2^i}
       \qquad(1\le i<n),
 \tag{P-F19l}
\]
as is checked by substitution and the first equation.
The numerator factor $2^i-(-1)^i$ is odd. Since $a_{n-1}$
is an integer, (P-F19l) forces $2^{n-1}\mid c$.
But $c=(3/g)s$ and $3/g$ is odd, so $2^{n-1}\mid s$.
This proves (P-F19j). Possible cancellation among the surviving
monomials causes no problem: nonzero evaluation guarantees that
at least one such monomial occurs, and every one obeys the same
divisibility condition.

\emph{Diagonal-stabilizer calculation.}
Let $\zeta$ be a primitive root of unity of order $N=3\,2^{n-1}$
and define
\[
 A=\operatorname{diag}
  (\zeta,\zeta^{-2},\zeta^{(-2)^2},\ldots,
                         \zeta^{(-2)^{n-1}})\in GL_n.
\]
Each chain monomial in (P-F19h) is fixed by $A$, since
$2(-2)^{i-1}+(-2)^i=0$; the last monomial is fixed since
$3(-2)^{n-1}$ is a multiple of $N$. The same is true with
$A^{-1}$, so $A\cdot F_n=F_n$ for the specified action on forms.
Set
\[
 S_n=1-2+(-2)^2+\cdots+(-2)^{n-1}
     =\frac{1-(-2)^n}{3}.
\]
The integer $S_n$ is odd and satisfies $S_n\equiv n\pmod3$.
Consequently
\[
 \operatorname{ord}(\det A)
 =\frac{N}{\gcd(N,S_n)}
 =\frac{3\,2^{n-1}}{g}.
 \tag{P-F19m}
\]
An $SL_n$-invariant polynomial $P$ of coefficient degree
$D=d_0s$ satisfies, as a function on $W^*$,
\[
 P(a\cdot F)=(\det a)^{-3s/g}P(F).
 \tag{P-F19n}
\]
This follows by writing $a=zh$ as above: $h$ does not change
$P$, while a scalar sends $F$ to $z^{-3}F$.
If $P(F_n)\ne0$, evaluation at its stabilizer $A$ therefore
requires $(\det A)^{3s/g}=1$. By (P-F19m),
$(3\,2^{n-1}/g)\mid(3s/g)$, which is again equivalent to
$2^{n-1}\mid s$. This verifies (P-F19j) independently of
the coefficient recurrence.

Finally, let $q_1,\ldots,q_r$ be any homogeneous system of
parameters of $T^{(n)}$. The algebra is finite as a module over
$\mathbb C[q_1,\ldots,q_r]$, so
$T^{(n)}/(q_1,\ldots,q_r)$ is finite-dimensional and graded.
Every positive-degree homogeneous element in this quotient is
nilpotent: its powers eventually have degree beyond the largest
nonzero graded piece. In particular
$\Delta^M\in(q_1,\ldots,q_r)$ for some integer $M\ge1$.
If every $q_i(F_n)$ were zero, evaluation of this ideal membership
would give $\Delta(F_n)^M=0$, contradicting (P-F19i).
Thus some $q_i$ is nonzero at $F_n$. Its positive stretch degree
is divisible by $2^{n-1}$ by (P-F19j), and is at least that
large. This applies to every h.s.o.p., including one with a
lexicographically minimal degree sequence, proving (P-F19c).

The inner module is the single irreducible $S^{(3)}V$, and
the partitions in (P-F19a) have heights one and $n$.
Thus (P-F19d) follows directly from the source's stated
parameter. Exponential growth of $2^{n-1}$ exceeds every
polynomial in $h=O(n^6)$, proving the asserted failure of the
canonical modular-index clause.
\end{proof}

\paragraph{Exact scope of the refutation.}
GCT VI, Hypothesis 1.6.7 first asks for some positive rational
form of polynomial modular index, and then explicitly asks for
this bound on the minimal form arising from a lexicographically
minimal h.s.o.p. of its canonical ring. P-F19 refutes the latter,
stronger clause. Its proof concerns values of actual invariants
at a particular semistable point. A positive rational expression
for the Hilbert series need not arise from parameters of that
ring, so the first existence clause is not refuted by this
argument alone. P-F30--P-F31 subsequently close that gate by proving
an actual primitive-$2^{n-1}$ pole of the ordinary Hilbert series,
and thereby refute the first clause as well. The paired-wave argument below supplies the additional
estimates needed to prove superpolynomial positivity and saturation indices.

This is an irreducible-inner ordinary plethysm family, unlike
the reducible torus embeddings in P-F14--P-F15. Determinant
twisting is absent: the inner shape is always $(3)$, and only
the number of variables grows. No historical
priority is asserted for the chain forms, their stabilizers,
or the resulting invariant-degree obstruction. The same chain
forms, after reversing the variables, already appear in
Bus\'e--Jouanolou, Example 4.12, printed p.~42, with attribution
to Demazure; the statement here gives their explicit consequence
for the specified canonical plethysm parameter bound.



% Source: branches/plethysm/sections/positivity_noncancellation.tex
% Reproduced inherited proof; finding P-F23.
\subsection{A geometric condition that prevents a cyclotomic cancellation}

\begin{proposition}[P-F23: a unique smooth generic fixed component gives an actual pole]
\label{prop:fixedcomponentpole}
Let $R=\bigoplus_{d\ge0}R_d$ be a finitely generated positively
graded complex domain with $R_0=\mathbb C$. Fix $M\ge2$ and a
primitive $M$th root of unity $\zeta$, and let $g$ act on $R_d$
by multiplication by $\zeta^d$. Let $X=\operatorname{Spec}R$ and
let $X^g$ be the scheme-theoretic fixed locus. Suppose that
$X^g$ has a unique irreducible component $Y$ of maximal dimension
$\delta\ge1$, and that $X$ is smooth at the generic point of $Y$.
Then $\operatorname{Hilb}_R(t)$ has a pole of order exactly
$\delta$ at $t=\zeta$.

More precisely, give $Y$ its reduced induced structure, write
$\mathbb C[Y]=\bigoplus_{k\ge0}\mathbb C[Y]_{Mk}$, and set
\[
 h_Y(z)=\sum_{k\ge0}\dim\mathbb C[Y]_{Mk}\,z^k,
 \qquad
 e_Y=\lim_{z\to1}(1-z)^\delta h_Y(z)>0.
\]
At the generic point of $Y$, the conormal space of $Y$ in $X$
has no $g$-eigenvalue equal to one. If its eigenvalues, with
multiplicity, are $\nu_1,\ldots,\nu_c$, where
$c=\dim X-\delta$, then
\[
 \lim_{t\to\zeta}(1-t^M)^\delta\operatorname{Hilb}_R(t)
       =\frac{e_Y}{\prod_{j=1}^c(1-\nu_j)}\ne0.
 \tag{P-F23a}
\]
In particular every positive rational expression
$H(t)/\prod_i(1-t^{a_i})$ for this Hilbert series has some
denominator weight divisible by $M$.
\end{proposition}

\begin{proof}
Choose positive-degree homogeneous algebra generators of $R$.
List those whose degrees are not divisible by $M$ as
$f_1,\ldots,f_l$, with $\deg f_i=w_i$. Their ideal $I$
is exactly the fixed-locus ideal. Indeed $f_i$ itself is a
nonzero scalar multiple of $g(f_i)-f_i$, and after killing
these generators every remaining algebra generator is fixed
by $g$. Thus
\[
 X^g=\operatorname{Spec}(R/I),
 \qquad (R/I)_d=0\quad\hbox{if }M\nmid d.
 \tag{P-F23b}
\]
The ideal is homogeneous, so every irreducible component of
the fixed locus is graded; equivalently the connected grading
torus cannot permute a finite set of components nontrivially.
In particular $Y$ is graded and its coordinate ring is
concentrated in degrees divisible by $M$.

Let $K_\bullet=K_R(f_1,\ldots,f_l)$ be the graded Koszul
complex, with its $i$th exterior generator given internal degree
$w_i$ and $g$-character $\zeta^{w_i}$. Denote its homology
modules by $N_j$. Multiplication by each $f_i$ on the Koszul
complex is homotopic to zero, so every $N_j$ is annihilated
by $I$ and is a finite graded $R/I$-module. The degreewise
Euler characteristic identity is
\[
 \operatorname{Hilb}_R(t)\prod_{i=1}^l(1-t^{w_i})
       =\sum_{j=0}^l(-1)^j\operatorname{Hilb}_{N_j}(t).
 \tag{P-F23c}
\]
All factors on the left except the Hilbert series are nonzero
at $\zeta$.

We describe carefully the leading contribution of the modules
$N_j$. Every finite graded module over $R/I$ splits into its
internal-degree classes modulo $M$. Each class, shifted by
its residue, is a graded module for the regrading in which
degree $M$ becomes degree one. Its support has dimension at
most $\delta$, so its Hilbert series has pole order at most
$\delta$ at $t=\zeta$. A component of dimension less than
$\delta$ has smaller pole order and contributes zero after
multiplication by $(1-t^M)^\delta$.

For a module $N$ whose localization at the generic point
$\eta$ of $Y$ is a vector space over $\mathbb C(Y)$, the
leading term is
\[
 \lim_{t\to\zeta}(1-t^M)^\delta\operatorname{Hilb}_N(t)
       = e_Y\operatorname{Tr}(g\mid N_\eta).
 \tag{P-F23d}
\]
Here $g$ acts trivially on the field $\mathbb C(Y)$, since it
fixes $Y$ pointwise. One way to verify (P-F23d), including its
grading shifts, is to take a graded prime filtration separately
in each residue class of $N$. All factors of dimension less
than $\delta$ disappear in the displayed limit. Since $Y$ is
the unique maximal-dimensional component, the remaining factors
are shifted copies of $\mathbb C[Y]$. A factor with generator
of internal degree $a$ contributes
$t^a h_Y(t^M)$, so its leading contribution is
$\zeta^a e_Y$. Upon localization, these factors form a
filtration of $N_\eta$ by one-dimensional character spaces of
character $\zeta^a$, giving exactly the trace. This also shows
why lower-dimensional components and torsion cannot cancel
that leading coefficient. The existence and positivity of
$e_Y$ is the ordinary positive leading Hilbert coefficient
of the positively graded domain obtained by regrading
$\mathbb C[Y]$; equivalently take any common Veronese and
use its Hilbert polynomial.

We now compute the generic Koszul homology. Smoothness of
$X$ along a dense open subset of $Y$ implies that the fixed
scheme is smooth there. For clarity, this familiar fact can
be seen directly at a smooth closed fixed point: diagonalize
the finite-order action on the cotangent space and lift each
eigenvector to the maximal ideal by the corresponding averaging
projector. These lifted eigenvectors are regular formal
parameters, and the completed local ring is the power series
ring in them, with the indicated diagonal action. Its fixed
scheme is cut out by exactly the parameters whose eigenvalues
are not one. Thus the fixed scheme is reduced and regularly
embedded at such points, and its conormal representation has
no trivial character. This applies at the generic point of
$Y$ as well.

Put $K=\mathbb C(Y)$. In the regular local ring $R_\eta$,
the localized ideal $I_\eta$ is the maximal ideal, generated
by a regular sequence of length $c=\dim X-\delta$. There
is a natural exact sequence of $K$-representations
\[
 0\longrightarrow L\longrightarrow
 E:=\bigoplus_{i=1}^l K e_i
 \longrightarrow V:=I_\eta/I_\eta^2\longrightarrow0,
 \qquad e_i\longmapsto f_i\bmod I_\eta^2.
 \tag{P-F23e}
\]
The characters on $E$ are $\zeta^{w_i}\ne1$.
The conormal space $V$ has the characters
$\nu_1,\ldots,\nu_c\ne1$, so the kernel $L$ also has
only nontrivial characters.

The localized Koszul homology is naturally
\[
 (N_j)_\eta\cong\bigwedge^j_K L.
 \tag{P-F23f}
\]
For an explicit verification, choose a minimal subcollection of
the $f_i$ that generates $I_\eta$. It is a regular sequence;
subtract its linear combinations from the other generators.
This invertible change of the Koszul presentation gives the
Koszul complex of that regular sequence tensored with the
Koszul complex on $l-c$ zeros. Its homology is an exterior
algebra on $l-c$ copies of $K$. The natural map in homological
degree one sends a Koszul cycle to its coefficient vector modulo
$I_\eta$ and identifies this space with the kernel in
(P-F23e). The exterior-algebra identification is natural and
therefore $g$-equivariant, even though no equivariant choice
of the intermediate change of generators is needed.
In particular the modules satisfy the vector-space hypothesis
used in (P-F23d).

Taking alternating traces gives
\[
 \sum_j(-1)^j\operatorname{Tr}(g\mid (N_j)_\eta)
       =\det(1-g\mid L).
\]
Multiplying (P-F23c) by $(1-t^M)^\delta$, applying
(P-F23d), and using (P-F23e), we obtain
\begin{align*}
 \lim_{t\to\zeta}(1-t^M)^\delta\operatorname{Hilb}_R(t)
 &=e_Y\frac{\det(1-g\mid L)}{
                         \prod_i(1-\zeta^{w_i})}\\
 &=\frac{e_Y}{\det(1-g\mid V)},
\end{align*}
which is (P-F23a). None of its factors is zero. The module
argument already bounded the pole order by $\delta$, and
the nonzero limit makes it exactly $\delta$.

If every $a_i$ in a rational expression were not divisible
by $M$, then every denominator factor $1-t^{a_i}$ would
be nonzero at a primitive $M$th root. Its polynomial numerator
could not create a pole. This proves the final assertion;
numerator positivity is not needed for that step.
\end{proof}

\paragraph{Standard algebra used in the proof.}
The Koszul change of generators, null homotopies, and tensor-product
properties used above are stated in the Stacks Project,
Section~15.29, Lemmas~15.29.3--15.29.6 and 15.29.12
(\url{https://stacks.math.columbia.edu/tag/0621}). The vanishing
of positive Koszul homology for a regular sequence is
Lemma~15.31.2, Tag~062F
(\url{https://stacks.math.columbia.edu/tag/062F}). The
equivariant excess calculation and the leading-character formula
needed here are proved explicitly above, rather than inferred
from an unqualified fixed-point or purity slogan.

\paragraph{Several components and the cubic application.}
The same proof, when every maximal-dimensional fixed component
is generically in the smooth locus of $X$, gives a sum
\[
 \sum_{Y:\,\dim Y=\delta}
       \frac{e_Y}{\det(1-g\mid N^*_{Y/X,\eta_Y})}
 \tag{P-F23g}
\]
for the limit in (P-F23a). Its terms can have different complex
phases. The uniqueness hypothesis avoids that additional
cancellation problem. Generic smoothness is also substantial:
P-F21 and P-F22 show that a universal factor by itself can
disappear completely from the ordinary Hilbert series even
for normal Gorenstein rational domains.

For the chain cubic family, P-F24 below proves generic smoothness
along the one-dimensional chain component, including the full
stabilizer calculation. P-F26 identifies the entire smooth-hypersurface
fixed locus with its orbit and supplies finite two-power lifts.
P-F27 excludes graph-support boundary competitors, even when the
lift order was not assumed to be $2^{n-1}$. The later full support
theorem P-F30 also excludes the remaining singular competitors,
including disconnected graph parts with triple monomials. P-F31
then proves that the fixed scheme has the unique chain curve as
its support and applies the present proposition to obtain the
actual ordinary-series pole. Thus this proposition supplies the
noncancellation step of that application, while its abstract
multiple-component caution remains necessary.



% Source: branches/plethysm/sections/positivity_chain_smoothness.tex
% Reproduced inherited proof; finding P-F24.
\subsection{The chain contributes a smooth fixed component}

\begin{proposition}[P-F24: full chain stabilizer and local fixed dimension]
\label{prop:chainsmoothcomponent}
Let $n\ge4$, $M=2^{n-1}$, $g=\gcd(n,3)$, and let $F_n$ be the
chain cubic of P-F19. Its full projective linear automorphism group
is the diagonal cyclic group of order $M$. Its stabilizer as an
actual cubic form under $SL_n$ is exactly the scalar subgroup
$\boldsymbol\mu_g$, which acts trivially on all cubic forms.

Let $X=\operatorname{Spec}T^{(n)}$, with the stretch grading of
P-F19, and let $p$ be the image of $F_n$ in this quotient. Then
$X$ is smooth at $p$. The fixed scheme for the grading action of
$\boldsymbol\mu_M$ is smooth of dimension one at $p$.
The closure of the grading orbit of $p$ is consequently an
irreducible one-dimensional component of this fixed scheme,
and its nonzero grading orbit consists of smooth points of $X$.
If $3\nmid n$, this component has coordinate ring
$\mathbb C[z]$ with $\deg z=M$; its regraded leading Hilbert
coefficient $e_Y$ in P-F23 is therefore one.
This does not classify the other components of the fixed scheme.
\end{proposition}

\begin{proof}
We first determine the full projective automorphism group; computing
only its diagonal subgroup would not suffice. Half the Hessian matrix
of $F_n$ is the symmetric tridiagonal matrix with diagonal
\[
 (x_2,x_3,\ldots,x_n,3x_n)
\]
and successive off-diagonal entries $x_1,\ldots,x_{n-1}$.
We claim that its projective rank-at-most-two locus consists exactly of
\[
 [e_1],\qquad
 [a e_{n-1}+b e_n]\quad
 \text{with }a=0\text{ or }a^2=3b^2.
 \tag{P-F24a}
\]
The latter are three distinct points on a line not containing $[e_1]$.

Here is a full rank check. A tridiagonal block of length $l$ with
all its $l-1$ off-diagonal entries nonzero has kernel dimension at
most one: its first equation determines the second coordinate from
the first, and the subsequent equations determine the rest in order.
The zero off-diagonal entries divide the whole matrix into such
blocks. Its rank is therefore at least the number of nonzero
entries among $x_1,\ldots,x_{n-1}$. At rank at most two, at most
two of these coordinates are nonzero.

If precisely two are nonzero, say $x_i,x_j$ with $i<j$, first
suppose $j\ge i+2$. The block on indices $i,i+1$ has determinant
$-x_i^2$: its first diagonal entry $x_{i+1}$ is zero. It has
rank two. The separate block containing the other nonzero
off-diagonal has rank at least one, a contradiction. If
$j=i+1$ and $i\ge2$, the singleton on index $i-1$ has nonzero
diagonal $x_i$ and is separate from the length-three block on
$i,i+1,i+2$. Their ranks are at least one and two. If $i=1$,
then for $n\ge5$ the block on $1,2,3$ has determinant
$-x_2^3\ne0$. For $n=4$ this remains true when $x_4=0$;
when $x_4\ne0$, the singleton on index four has rank one and
the length-three block has rank at least two. This excludes
every case with two nonzero off-diagonal entries.

If the only nonzero one is $x_j$ with $2\le j\le n-2$,
the block on $j,j+1$ has determinant $-x_j^2$, and the
singleton on $j-1$ has diagonal $x_j$, again giving rank
at least three. If $j=1$, its length-two block already has
rank two; $x_n\ne0$ would contribute nonzero diagonal entries
on the disjoint last two indices. Thus only $[e_1]$ remains.
If $j=n-1$, put $a=x_{n-1}$ and $b=x_n$. The only nonzero
part is the singleton $a$ on index $n-2$ and the block
\[
 \begin{pmatrix}b&a\\a&3b\end{pmatrix}.
\]
Its total rank is two exactly when $a^2=3b^2$.
Finally, if there are no nonzero off-diagonal entries, the
projective point is $[e_n]$ and the rank is two. This proves
(P-F24a).

Under a linear change of coordinates the Hessian transforms by
congruence, so any projective linear automorphism of $F_n$
preserves (P-F24a). The three collinear points form its unique
collinear triple. Hence its remaining point $[e_1]$ is fixed.
Write a representing linear transformation in coordinate forms
$\ell_1,\ldots,\ell_n$, with
\[
 F_n(\ell_1,\ldots,\ell_n)=cF_n(x_1,\ldots,x_n),\qquad c\ne0.
\]
Fixing $[e_1]$ means $\ell_1=ax_1+L(x_2,\ldots,x_n)$,
$a\ne0$, while the other $\ell_i$ do not involve $x_1$.
Comparing the coefficients of $x_1^2$ gives
$a^2\ell_2=cx_2$, so $\ell_2=bx_2$ with $b\ne0$.
The coefficient of $x_1$ is $2aL\ell_2$ and must vanish.
Thus $L=0$. Restriction to $x_1=0$ leaves an automorphism
of the chain cubic in $n-1$ variables, whose first coordinate
form is a scalar multiple of its first variable.

This proves by induction that all the coordinate forms are
diagonal once the base case $n=4$ is checked. For that case
write the remaining ternary cubic as
\[
 G(y,z,w)=y^2z+z^2w+w^3.
\]
Its transformation has coordinate forms $by$, $uy+A(z,w)$,
and $vy+B(z,w)$, where $A,B$ are independent linear forms.
The coefficient of $y$ in its transformed cubic is
\[
 vA^2+2uAB+3vB^2.
\]
It is zero identically, so independence of $A,B$ implies
$u=v=0$. Comparison of the coefficient of $y^2$ now gives
$b^2A=cz$, so $A$ is a nonzero multiple of $z$.
Write $B=\alpha z+\beta w$, with $\beta\ne0$. The
coefficient of $zw^2$ is $3\alpha\beta^2$; hence
$\alpha=0$. This is the required diagonal form and completes
the induction.

For diagonal projective transformations choose a scalar
representative with last entry one. The last monomial then
forces the common multiplier to be one, and the other
monomials give $a_{i+1}=a_i^{-2}$. The last-entry condition is
$a_1^{(-2)^{n-1}}=1$. Thus the full projective group is
cyclic of order $M$. For linear transformations preserving
$F_n$ exactly, rather than its projective class, the same
recurrence instead ends with $a_n^3=1$. They form the
cyclic group of order $3M$, with entries
\[
 a_i=a_1^{(-2)^{i-1}},\qquad a_1^{3M}=1.
\]
Their determinant is $a_1^S$, where
\[
 S=1-2+4-\cdots+(-2)^{n-1}
   =\frac{1-(-2)^n}{3}.
\]
The integer $S$ is odd and satisfies $S\equiv n\pmod3$.
Consequently $\gcd(S,3M)=g$. The determinant-one kernel
has order $g$; when $g=3$, its entries are all the same
cube root of unity, because $(-2)^i\equiv1\pmod3$.
This proves the asserted $SL_n$ stabilizer.

By P-F19 the chain hypersurface is smooth. The standard
stability theorem for smooth hypersurfaces of degree at
least three implies that its nonzero affine $SL_n$ orbit
is closed. An exact source is Huybrechts, \emph{The geometry
of cubic hypersurfaces}, Chapter~3, Corollary~1.6, printed
pp.~149--150, together with Exercise~1.4 on p.~149 for the
affine formulation. The cited degree and dimension conditions
hold here, since $n\ge4$. Passing to $SL_n/\boldsymbol\mu_g$
does not change the quotient and makes the stabilizer trivial.
Luna's \emph{Slices \'etales}, Section~III.1, the unnumbered
slice theorem on p.~97 and Remark~1$^\circ$ on pp.~97--98, identify an
\'etale neighborhood of the quotient point with a neighborhood
in the quotient of a smooth normal slice by its stabilizer.
The stabilizer here is trivial, so the slice and the quotient
are smooth. Thus $X$ is smooth at $p$.

It remains to determine the local fixed dimension; smoothness
alone does not give it. Let $\zeta$ be primitive of order $M$,
and put
\[
 A=\operatorname{diag}
       (\zeta,\zeta^{-2},\ldots,\zeta^{(-2)^{n-1}}).
\]
This transformation fixes $F_n$. Its determinant is
$\zeta^S$ with $S$ odd. On degree-$s$ elements of the
canonical algebra $T^{(n)}$, it acts by the character
$\det(A)^{-3s/g}$, up to replacing $A$ by its inverse
according to the coordinate-action convention. Since
$3/g$ and $S$ are odd, this character generates exactly
the grading $\boldsymbol\mu_M$ action.

The only degree-three monomials fixed by $A$ are the $n$
monomials of $F_n$. To see this, the nonzero residue weights
$(-2)^i\pmod M$, $0\le i\le n-2$, have distinct
2-adic valuations, and the remaining weight is zero.
In a sum of three such weights that is zero modulo $M$,
the least valuation cannot occur only once. Its index
therefore occurs twice. If it is $i<n-1$, the third
weight must be $-2(-2)^i=(-2)^{i+1}$, giving exactly
a chain monomial. At the last nonzero weight, twice that
weight is zero and the third is the zero weight. Three
zero weights give $x_n^3$; three equal nonzero weights
never sum to zero because three is odd. This exhausts
the possibilities. Thus the fixed subspace of the cubic
representation has dimension $n$.

All eigenvalues of $A$ on the underlying $n$-space are
distinct, so its fixed subspace on $\mathfrak{sl}_n$
under conjugation consists of the traceless diagonal
matrices and has dimension $n-1$. The map
$\mathfrak{sl}_n\to T_{F_n}(SL_nF_n)$ is an isomorphism,
because the stabilizer is finite. Luna's smooth quotient
description, or the corresponding differential of the
local principal bundle, gives the exact tangent sequence
\[
 0\longrightarrow\mathfrak{sl}_n F_n
 \longrightarrow W^*
 \longrightarrow T_pX\longrightarrow0.
\]
It is $A$-equivariant. Taking fixed vectors is exact for
this finite-order action in characteristic zero. Hence
\[
 \dim(T_pX)^{\boldsymbol\mu_M}=n-(n-1)=1.
 \tag{P-F24b}
\]
The fixed scheme of a finite-order automorphism of a
smooth complex variety is smooth near a fixed point:
averaging local parameters into eigenspaces and passing
to the completed local ring identifies its fixed ideal
with the parameters of nontrivial character. This is
also the local argument used in P-F23. It follows from
(P-F24b) that the fixed scheme is smooth of dimension one
at $p$.

The positive-degree discriminant is nonzero at $p$, so
the grading orbit of $p$ has dimension one. All of it
is fixed by $\boldsymbol\mu_M$, since the grading group
is abelian. Its irreducible closure is contained in a
component of the fixed scheme through $p$. That component
has dimension one locally at its smooth point $p$, and
therefore has dimension one globally. The two irreducible
closed sets agree. Grading translations preserve smoothness
of $X$, so the entire nonzero orbit consists of smooth
points.

For the last assertion assume $3\nmid n$, so $g=1$.
Restriction to the grading orbit is the homomorphism
\[
 T^{(n)}\longrightarrow\mathbb C[u,u^{-1}],\qquad
 f\in T_s\longmapsto f(p)u^s.
\]
All exponents are nonnegative, and P-F19 shows that a
nonzero image has $M\mid s$. The discriminant has stretch
degree exactly $M$ and is nonzero at $p$. The image therefore
contains a nonzero multiple of $u^M$ and is exactly
$\mathbb C[u^M]$. This image is the coordinate ring of
the orbit closure, proving the claim. Its Hilbert series is
$1/(1-t^M)$, so $e_Y=1$. This completes the proof.
\end{proof}

\paragraph{The geometric gate isolated at this step.}
The diagonal subgroup itself is a known computation:
Zheng, \emph{On abelian automorphism groups of hypersurfaces},
arXiv:\texttt{2004.09008v3}, Lemma~3.2, printed p.~5,
computes the type-$T_n$ diagonal group. The Hessian argument
above verifies the full-group assertion needed here.
No historical priority is claimed. The local result establishes
the generic-smoothness hypothesis of P-F23 for the chain component.
What remains is to control all the other components of
$X^{\boldsymbol\mu_M}$, including those represented by
singular polystable cubics, sufficiently to prevent their
leading contributions from cancelling the chain contribution.
Uniqueness of the maximal-dimensional component would suffice.
This local proposition alone does not assert it; P-F30--P-F31
below establish it by a complete support classification and
thereby close the global noncancellation gate.

Primary sources for the cited geometric theorems are
\url{https://www.math.uni-bonn.de/people/huybrech/Notes.pdf}
and \url{https://www.numdam.org/item/10.24033/msmf.110.pdf}.



% Source: branches/plethysm/sections/positivity_finite_lift.tex
% Reproduced inherited proof; finding P-F26.
\subsection{Finite lifts and the entire smooth fixed locus}

\begin{proposition}[P-F26: finite two-power lifts and the smooth chain locus]
\label{prop:finitecubiclift}
Let $n\ge4$, $M=2^{n-1}$, $g=\gcd(n,3)$, and
$X=\operatorname{Spec}T^{(n)}$ with the canonical stretch grading
of P-F19. Write $\Delta$ for the discriminant, viewed as a
homogeneous function on $X$, and let $p_n$ be the quotient point
of the chain cubic $F_n$. Then the following statements hold.

\begin{enumerate}
\item Every nonzero closed point $p\in X^{\boldsymbol\mu_M}$
has a representative cubic $F$ with closed $SL_n$ orbit and an
exact-form stabilizer element $A\in GL_n$ of finite two-power
order such that $\det A$ has order exactly $M$.
\item If $\Delta(p)\ne0$, the lift $A$ can be chosen of order
exactly $M$, and $F$ is $GL_n$-equivalent to $F_n$.
\item Scheme-theoretically, the part of the fixed scheme in
the smooth-hypersurface locus is the single grading orbit:
\[
 X^{\boldsymbol\mu_M}\cap D(\Delta)
      =\mathbb G_m\cdot p_n.
 \tag{P-F26a}
\]
This is a smooth reduced curve. Every irreducible component of
$X^{\boldsymbol\mu_M}$ other than the chain component of P-F24
is contained in $V(\Delta)$.
\end{enumerate}
The first statement does not assert that the lift has order $M$
when $F$ is singular.
\end{proposition}

\begin{proof}
Put $\gamma=3/g$, an odd integer. Recall the actual ring
identification from P-F19:
\[
 T^{(n)}_s=\bigl(\operatorname{Sym}^{(n/g)s}
                         (\operatorname{Sym}^3\mathbb C^n)\bigr)^{SL_n}.
\]
The $GL_n$-character on this subspace is $(\det)^{\gamma s}$.
On points, with the convention $(B\cdot F)(x)=F(B^{-1}x)$,
one consequently has
\[
 P(B\cdot F)=(\det B)^{-\gamma s}P(F)
       \quad(P\in T^{(n)}_s).
 \tag{P-F26b}
\]
Fix a primitive $M$th root $\zeta$. Since $\gamma$ is odd,
there is a unique $\eta\in\boldsymbol\mu_M$ with
$\eta^{-\gamma}=\zeta$, and $\eta$ is primitive. Choose
$B\in GL_n$ with determinant $\eta$. Thus the action of $B$
on the quotient is the grading action with parameter $\zeta$.
This choice obtains a determinant of order $M$ directly; a
scalar-only lift could instead introduce an unnecessary odd
factor into its determinant order.

Every fiber of the affine reductive quotient contains exactly
one closed orbit. A precise primary source for this fact is
Luna, \emph{Slices \'etales}, Section~I.1, printed p.~84.
Choose $F$ in that closed orbit over $p$. The form $B\cdot F$
has a closed $SL_n$ orbit, since $B$ normalizes $SL_n$.
It has the same quotient point by (P-F26b) and the fixed-point
hypothesis. Hence $B\cdot F=S\cdot F$ for some $S\in SL_n$.
Replacing $B$ by $S^{-1}B$ produces a matrix fixing $F$
exactly and still of determinant $\eta$. Denote this matrix
temporarily by $C$.

We next replace $C$ by a finite two-power element without
changing its determinant. Let $C=C_sC_u$ be its commuting
semisimple and unipotent parts. Their actions on cubic forms
are again semisimple and unipotent. Decompose $F$ into the
eigenspaces of $C_s$ in that representation. On an eigenspace
of eigenvalue $a\ne1$, the operator $aC_u-1$ is invertible,
because every eigenvalue of $C_u$ is one. Since $C F=F$,
the components in those eigenspaces are zero. It follows that
$C_sF=F$; also $\det C_s=\det C=\eta$.

Diagonalize $C_s$, with eigenvalues $\alpha_1,\ldots,\alpha_n$
on the underlying vector space. Let
$\Gamma=\langle\alpha_1,\ldots,\alpha_n\rangle\subset\mathbb C^*$.
This is a finitely generated abelian group. Its torsion
subgroup is finite cyclic, and a choice of a free complement
gives a group homomorphism
\[
 \rho:\Gamma\longrightarrow\Gamma_{\mathrm{tors},2}
\]
that is the identity on its two-primary torsion subgroup,
and kills its odd torsion and the chosen free complement.
Here $\eta=\prod_i\alpha_i$ belongs to that two-primary
torsion subgroup, so $\rho(\eta)=\eta$.

Every monomial $x_1^{e_1}\cdots x_n^{e_n}$ with nonzero
coefficient in $F$ satisfies $\prod_i\alpha_i^{e_i}=1$;
inverting all the eigenvalues, as in the action on forms,
gives the same equality. Applying $\rho$ preserves this
relation. Therefore
\[
 A=\operatorname{diag}(\rho(\alpha_1),\ldots,
                       \rho(\alpha_n))
 \tag{P-F26c}
\]
fixes $F$ exactly, has finite two-power order, and has
determinant $\eta$. This proves the first statement.

Now assume $\Delta(p)\ne0$. The representative $F$ defines
a smooth hypersurface. Write the order of $A$ as $2^\ell$.
Because its determinant has order $M$, one has $\ell\ge n-1$.
At least one eigenvalue of $A$ has order $2^\ell$.
We give a direct smoothness argument for the reverse bound
and its equality case.

Suppose the eigenvalue on variable $x_i$ has nontrivial
two-power order. The coefficient of $x_i^3$ in an invariant
cubic must be zero. Thus the coordinate point $[e_i]$ lies
on the hypersurface. Smoothness there implies the presence
of a monomial $x_i^2x_j$ with nonzero coefficient for some
$j\ne i$. Invariance forces
\[
 \alpha_j=\alpha_i^{-2},
 \tag{P-F26d}
\]
where in this paragraph $\alpha_i$ denotes the eigenvalue
of the finite-order matrix $A$ itself. Its order is half
the order of $\alpha_i$. Starting with an eigenvalue of
order $2^\ell$ and iterating gives eigenvalues of orders
\[
 2^\ell,2^{\ell-1},\ldots,2,1.
\]
They occupy $\ell+1$ distinct variables. Hence $\ell+1\le n$.
Consequently $\ell=n-1$, and the chosen lift has order
exactly $M$.

In the equality case these are all the variables. After
permutation their eigenvalues are
\[
 \xi,\xi^{-2},\ldots,\xi^{(-2)^{n-1}}=1,
 \qquad \operatorname{ord}(\xi)=M.
 \tag{P-F26e}
\]
For completeness, the invariant degree-three monomials for
this spectrum are exactly those of the chain. The nonzero
residue weights $1,-2,4,\ldots,(-2)^{n-2}$ have distinct
two-adic valuations. In a vanishing sum of three weights,
a least nonzero valuation cannot occur only once. If it
occurs three times, the sum is not zero, since three is odd.
Thus it occurs twice, and the third weight must be minus
twice that weight, which is precisely the next chain weight.
At the last nonzero weight the next is zero. The all-zero
case gives the final cube. Therefore
\[
 F=c_1x_1^2x_2+\cdots+c_{n-1}x_{n-1}^2x_n+c_nx_n^3.
 \tag{P-F26f}
\]
Every $c_i$ is nonzero. For $i<n$ this follows again from
smoothness at $[e_i]$; if $c_n=0$, the final coordinate
point is singular as well. Starting with $b_n^3c_n=1$
and successively solving $b_i^2b_{i+1}c_i=1$ gives nonzero
complex numbers $b_i$. The diagonal substitution
$x_i\mapsto b_ix_i$ turns (P-F26f) into $F_n$. This proves
the second statement, including the equality classification.

The $GL_n$-action on the quotient is precisely the grading
action by (P-F26b): the map $\mathbb C^*\to\mathbb C^*$,
$z\mapsto z^{-\gamma}$, is surjective. Thus every fixed
closed point in $D(\Delta)$ belongs to the grading orbit
of $p_n$. Conversely that entire orbit is fixed by
$\boldsymbol\mu_M$ and lies in $D(\Delta)$, by P-F24 and
discriminant nonvanishing.

P-F24 also shows that the fixed scheme is smooth of dimension
one at every point of this orbit. Its intersection with
$D(\Delta)$ consequently is reduced and has exactly the
closed points of that orbit. Both carry their reduced
locally closed scheme structures, so they agree, proving
(P-F26a). More explicitly, any nonreduced or additional
finite-type locus would have a closed complex point,
contradicting the preceding local smoothness and point
classification. If a different irreducible fixed component
met $D(\Delta)$, it would pass through a smooth point of
the fixed scheme on this orbit. That point has a unique
component, namely the chain component of P-F24. This is a
contradiction. Thus all other components lie in $V(\Delta)$.
\end{proof}

\paragraph{Known order bound, and the precise remaining lift issue.}
For smooth hypersurfaces the two-power order bound also follows
from Zheng, \emph{On abelian automorphism groups of hypersurfaces},
arXiv:\texttt{2004.09008v3}, Theorem~4.8, printed pp.~11--12.
An exact-form stabilizer of two-power order has the same
projective order, since its scalar kernel is contained in
$\boldsymbol\mu_3$.
At degree three its factors $|1-(-2)^a|$ are odd, and the only
two-power factors in the theorem are $2^{b-1}$ with $b\le n$.
The elementary coordinate-point argument above proves the needed
bound and makes its equality case explicit. Theorem~4.3 alone
states an abstract subgroup embedding into the diagonal group
of a simple hypersurface; it should not be cited by itself as
conjugacy of the original representation. No historical priority
is claimed for the order bound or its chain realization.

For a singular closed-orbit representative $F$, put
$H=\operatorname{Stab}_{GL_n}(F)$. Once the determinant image
contains $\boldsymbol\mu_M$, an exact-order-$M$ lift exists
if and only if the sequence of complex groups
\[
 1\longrightarrow H\cap SL_n\longrightarrow
   \det^{-1}(\boldsymbol\mu_M)\cap H
   \xrightarrow{\det}\boldsymbol\mu_M\longrightarrow1
 \tag{P-F26g}
\]
has a group-theoretic section. Indeed an order-$M$ element
mapping to a generator defines a section, and the image
of a generator under a section has order exactly $M$.
Such splitting is not automatic for abstract group extensions:
the squaring map from a cyclic group of order $2M$ onto a
cyclic group of order $M$ has no such section when $M$ is a
power of two. This abstract example is not asserted to occur
as the stabilizer extension of a cubic. The proof above
establishes splitting in the smooth case. At this step it
confines all possible competing fixed components to the
discriminant boundary without assuming a splitting there.
P-F30 subsequently excludes every nonzero singular competitor
at this extremal determinant order, even for larger-order lifts;
P-F31 then proves uniqueness and noncancellation. The abstract
extension warning explains why that additional work is needed.

The primary texts are
\url{https://www.numdam.org/item/10.24033/msmf.110.pdf}
and \url{https://arxiv.org/pdf/2004.09008}.



% Source: branches/plethysm/sections/positivity_graph_support.tex
% Reproduced inherited proof; finding P-F27.
\subsection{The singular support gate requires three distinct variables}

\begin{proposition}[P-F27: extremal determinant order for graph supports]
\label{prop:cubicgraphsupport}
Let $n\ge2$, and let $F\in\mathbb C[x_1,\ldots,x_n]$ be a
cubic whose nonzero monomials all have the forms $x_i^3$ or
$x_i^2x_j$ with $i\ne j$. Let $B$ be the integer matrix whose
columns are their exponent vectors. Assume that the vector
$\mathbf1=(1,\ldots,1)^t$ belongs to the rational column span
of $B$. This hypothesis holds, in particular, if $F$ is
$SL_n$-semistable.

Suppose a diagonal matrix $A$ of finite two-power order fixes
$F$ exactly. Then the order of $\det A$ is at most $2^{n-1}$.
If that order is $2^{n-1}$, the support is, after a permutation
of variables, exactly
\[
 x_1^2x_2,\ x_2^2x_3,\ldots,\ x_{n-1}^2x_n,\ x_n^3.
 \tag{P-F27a}
\]
All their coefficients are nonzero, the matrix $A$ itself has
order $2^{n-1}$, and a diagonal change of variables turns $F$
into the smooth chain cubic $F_n$.

Consequently, for $n\ge4$, every nonzero fixed quotient point
in the discriminant boundary of P-F26 has the following
property: in any eigenbasis for any finite two-power
exact-form lift with determinant of order $2^{n-1}$, its
closed-orbit representative contains a nonzero monomial
$x_ix_jx_k$ with $i,j,k$ pairwise distinct.
\end{proposition}

\begin{proof}
We first explain the rational-span assumption. If a
positive-degree $SL_n$-invariant polynomial in the
coefficients of $F$ does not vanish at $F$, at least one of
its coefficient monomials has nonzero value. Invariance
under the diagonal determinant-one torus means that the
sum of the cubic exponent vectors, counted with that
coefficient monomial's nonnegative integer multiplicities,
has all coordinates equal. The common coordinate is
positive, because the coefficient degree is positive.
Dividing by it expresses $\mathbf1$ as a nonnegative
rational combination of columns of $B$. Only membership
in their rational span will be used below.

Write the eigenvalues of $A$ as $\alpha_1,\ldots,\alpha_n$.
They lie in one finite cyclic two-group. The map
\[
 \psi:\mathbb Z^n\longrightarrow\mathbb C^*,\qquad
 (u_i)\longmapsto\prod_i\alpha_i^{u_i}
 \tag{P-F27b}
\]
annihilates the lattice $L$ generated by the columns of $B$,
since every supported monomial is fixed. Thus it factors
through $\mathbb Z^n/L$, and
$\psi(\mathbf1)=\det A$. This is the precise lattice
interpretation of the determinant character; the order
of the determinant need not equal the order of $A$.

Make a graph on the variables. A supported monomial
$x_i^2x_j$ gives an edge with coefficient two at vertex
$i$ and coefficient one at vertex $j$; a cube gives a loop
with coefficient three. Distinct oriented monomials may
give two edges on the same unordered pair. Decompose the
graph into connected components, ignoring loops for
connectivity. The exponent lattice and rational span
split over these components. In a component with $v$
vertices, therefore, its all-ones vector $\mathbf1_v$
belongs to the rational span of its support columns.
The determinant is the product of the images of these
component vectors under $\psi$.

We prove that each component contributes a determinant
character of order at most $2^{v-1}$, with a stronger
bound $2^{v-2}$ in the rank-deficient case when $v\ge2$.
An isolated vertex must have a cube, since otherwise
$\mathbf1_v$ would not belong to the column span. Its
relation is $3e_i=0$, and every character into a two-group
is trivial on $e_i$. This handles $v=1$.

For $v\ge2$, choose a spanning tree of the underlying
connected graph, and let $T$ be the $v\times(v-1)$ matrix
of its support columns. Its columns are independent:
removing a leaf removes a row with exactly one nonzero
entry, equal to one or two, and induction proves rank
$v-1$. Hence the full component matrix has rank either
$v-1$ or $v$.

\emph{The rank-deficient case.}
Suppose its rank is $v-1$. Then its rational column span
is that of $T$, so $Tc=\mathbf1_v$ has a unique rational
solution $c$. For every vertex $r$, let $T^{(r)}$ be the
square matrix obtained by deleting row $r$. Repeatedly
remove leaves other than $r$. This proves
\[
 D_r:=\det T^{(r)}=\pm2^{e_r},\qquad 0\le e_r\le v-1.
 \tag{P-F27c}
\]
Cramer's rule applied after deleting row $r$ shows that
the least positive integer $q$ with $qc$ integral divides
$|D_r|$. It therefore divides the greatest common divisor
of all these maximal minors. Also
$q\mathbf1_v=T(qc)$ lies in the tree lattice. Equivalently,
the torsion class of $\mathbf1_v$ in the tree cokernel has
order $q$, and its image in the full support cokernel has
order dividing $q$. This proves the needed maximal-minor
assertion directly, without assuming full rank.

The signed cofactor vector of $T$ is a nonzero vector
in the kernel of $T^t$. Along an edge with coefficients
two and one, its two endpoint coordinates consequently
satisfy $2u_i+u_j=0$. Their absolute values differ by a
factor of two. Since their absolute values are the
$|D_r|$, two adjacent exponents in (P-F27c) differ by one.
As all these exponents are at most $v-1$, their minimum
is at most $v-2$. Thus
\[
 q\mid 2^{\min_r e_r},\qquad \min_r e_r\le v-2.
 \tag{P-F27d}
\]
The component determinant character has order at most
$2^{v-2}$, as asserted. Notice that free directions in
the cokernel may allow the matrix eigenvalues themselves
to have much larger order; the span hypothesis makes
$\mathbf1_v$ a torsion class and is what controls its
character.

\emph{The full-rank case.}
Suppose instead that the component matrix has rank $v$.
One support column outside the span of $T$ extends its
columns to a square nonsingular matrix $C$. Its support
graph is connected with $v$ edges and $v$ vertices, and
has exactly one cycle, possibly a loop or a double edge.
The finite cokernel of $C$ has order $|\det C|$; one may
also use the adjugate identity to see directly that
$\det C$ kills every cokernel class. Since its image under
$\psi$ is a two-group, the component determinant order
is at most $2^{v_2(\det C)}$.

Let the unique cycle have length $l$. Expanding the
matrix successively along leaf rows contributes one
factor of either one or two for each of its $v-l$
tree edges. If $l=1$, the remaining loop contributes
three, and hence
\[
 v_2(\det C)\le v-1.
 \tag{P-F27e}
\]
If $l\ge2$, the cycle determinant has exactly two
permutation terms. For some $0\le a\le l$ they give,
up to an overall sign,
\[
 2^a+(-1)^{l-1}2^{l-a}.
 \tag{P-F27f}
\]
If $a=l-a$, then $l$ is even and the determinant is zero,
which is excluded. Otherwise its two-adic valuation is
$\min(a,l-a)\le\lfloor(l-1)/2\rfloor$. Therefore
\[
 v_2(\det C)
 \le v-l+\left\lfloor\frac{l-1}{2}\right\rfloor
 \le v-2\qquad(l\ge2).
 \tag{P-F27g}
\]
This proves the component bound in every case, including
multiple edges and rank deficiency.

Because all component characters have two-power order,
their product has order at most their largest order.
Every component has at most $n$ vertices. This proves
$\operatorname{ord}(\det A)\le2^{n-1}$.

Now suppose equality holds. At least one component must
have $n$ vertices, since every smaller component has
bound at most $2^{n-2}$. Thus the support graph is
connected. It has full rank by (P-F27d), and the chosen
unicyclic basis must have a loop by (P-F27g). Moreover,
equality in (P-F27e) requires every leaf expansion factor
to be two. Root its remaining tree at the loop vertex.
Each basis edge therefore has coefficient two at its
child and one at its parent. The eigenvalue relations are
\[
 \alpha_{\mathrm{root}}^3=1,\qquad
 \alpha_{\mathrm{parent}}=\alpha_{\mathrm{child}}^{-2}.
 \tag{P-F27h}
\]
Since all eigenvalues have two-power order, the root
eigenvalue is one. A vertex at depth $d$ then has
eigenvalue order at most $2^d$. If the maximum tree
depth is $h$, the order of every eigenvalue, and therefore
of their product, is at most $2^h$. Equality for the
determinant forces $h\ge n-1$. A rooted tree with $n$
vertices has $h\le n-1$, with equality only for a path
rooted at an endpoint. Thus this basis is the chain
support in (P-F27a).

The deepest eigenvalue must have order $2^{n-1}$: all
other vertices have smaller depth, and a product of
elements of smaller two-power order could not have the
required determinant order. Write this eigenvalue as
$\xi$, and order the path from its deepest vertex to the
root. Then the entire eigenvalue list is
\[
 \xi,\ \xi^{-2},\ \xi^4,\ldots,
 \xi^{(-2)^{n-1}}=1,
 \qquad \operatorname{ord}(\xi)=2^{n-1}.
 \tag{P-F27i}
\]
In particular, $A$ has this same order, although this
was not assumed.

These eigenvalues permit no additional cubic support.
Indeed their nonzero residue weights modulo $2^{n-1}$
have distinct two-adic valuations. In a zero sum of
three weights, the least nonzero valuation cannot occur
once, and cannot occur three times because three is
odd. It occurs twice at the same vertex; the third
weight must be its negative double, the next chain
weight. The all-zero case gives the root cube. Thus the
fixed cubic monomials are exactly (P-F27a). They are all
present, since they were the columns of the chosen basis.
Their nonzero coefficients can be made one by first
rescaling the root variable using a cube root and then
successively rescaling its predecessors using square
roots. Hence $F$ is the smooth chain cubic up to a
diagonal substitution, by the smoothness proof in P-F19.

Finally, P-F26 supplies a finite two-power lift with
primitive determinant for every nonzero fixed quotient
point. Its closed-orbit representative is semistable,
so it satisfies the span hypothesis in every coordinate
basis. If that representative is singular and any such
lift eigenbasis had only graph-support monomials, the
result just proved would identify it with a smooth
chain. This contradiction proves the last assertion.
\end{proof}

\paragraph{Connected graph part, even with additional triple monomials.}
The same equality classification holds for an arbitrary cubic satisfying
the rational-span hypothesis if its graph made from its square-linear
monomials is connected; cubes are retained as loops. Thus the presence
of a three-distinct-variable monomial by itself is not enough to produce
a competitor. In the discriminant boundary of P-F26, this graph must be
disconnected in every finite-lift eigenbasis.

Here is the additional proof. Let $G$ be the matrix consisting only
of the cube and square-linear support columns. Since its graph is
connected, its rank is $n$ or $n-1$. If it has rank $n$, its graph
subform satisfies the span hypothesis by itself. The proposition
therefore gives the chain eigenvalue list (P-F27i), which permits no
additional triple monomials. The full form is the chain.

Suppose instead that $G$ has rank $n-1$, and take the spanning-tree
matrix $T$ used above. Its signed cofactor vector has the form
\[
 u_i=\varepsilon 2^e(-2)^{h_i},\qquad
 \varepsilon\in\{1,-1\},\quad e\ge0,\quad
 h_i\ge0,\quad \min_i h_i=0.
 \tag{P-F27j}
\]
Indeed its coordinates are signed powers of two and the relation
$2u_i+u_j=0$ along each tree edge fixes both the change in exponent
and the change in sign. Connectivity makes the sign formula in
(P-F27j) consistent. Put $H=\max_i h_i$. The tree-minor calculation
gives $e+H\le n-1$.

If every triple support column lies in the rational span of $T$,
then so does $\mathbf1$, and the rank-deficient argument already
bounds the determinant order by $2^{n-2}$. Otherwise choose a triple
column $\beta=e_i+e_j+e_k$ outside that span, with $i,j,k$ distinct.
Cofactor expansion gives a nonsingular square matrix with
\[
 \det(T,\beta)=\pm 2^e
       \bigl((-2)^{h_i}+(-2)^{h_j}+(-2)^{h_k}\bigr).
 \tag{P-F27k}
\]
A nonzero sum of three powers of $-2$ has two-adic valuation strictly
less than the largest exponent occurring, unless all three exponents
are equal. To verify this, a unique smallest exponent gives exactly
that valuation. If two smallest exponents are $r$ and the third is
$t>r$, the sum is $(-2)^r(2+(-2)^{t-r})$: it vanishes when $t=r+1$,
and has valuation $r+1<t$ when $t\ge r+2$. Three equal exponents
$r$ give valuation $r$. Consequently the valuation in (P-F27k) is
at most $e+H$, with equality possible only if all three chosen
vertices have height $H$.

If $e+H\le n-2$, this already contradicts determinant order
$2^{n-1}$, since the determinant of the square support submatrix
annihilates the determinant character. In the only remaining case,
$e+H=n-1$, choose a vertex $r$ of height $H$. Its deleted-row tree
minor has absolute value $2^{n-1}$. Root the tree at $r$ and perform
the leaf expansion for this minor. Every one of its $n-1$ factors
must be two. Thus each tree edge has coefficient two at its child
and one at its parent, and the cofactor relation says
$h_{\rm parent}=h_{\rm child}+1$. The root is the unique vertex of
height $H$. Three distinct vertices cannot all have that height.
Therefore the valuation in (P-F27k) is again at most $n-2$, the
same contradiction. This completes the connected-graph extension.

\paragraph{Scope of the reduction.}
The proof treats graph supports and, by the extension just proved,
arbitrary supports with connected graph part. It does not classify
all singular semistable cubics. It does not infer
that the determinant order bounds an arbitrary finite lift
before imposing the extremal graph-support case. Three
distinct-variable monomials are exactly the remaining
cubic support type, and the unresolved case has a disconnected
graph part. Thus P-F27 removes the rank-deficiency and
larger-order-lift obstructions for the stated classes, while
leaving the disconnected hypergraph support problem unresolved by
this proposition alone. P-F30 subsequently resolves that remaining
case by component contraction and a determinant bound, and P-F31
uses it to control the global fixed support and ordinary pole.
No historical priority is claimed for this elementary
lattice and graph reduction.



% Source: branches/plethysm/sections/positivity_chain_residue.tex
% Reproduced inherited proof; finding P-F28.
\subsection{The exact local coefficient of the chain component}

\begin{proposition}[P-F28: chain multiplicity and cyclotomic valuation]
\label{prop:chainlocalcoefficient}
Let $n\ge4$, $M=2^{n-1}$, $g=\gcd(n,3)$, and
$m=\binom{n+2}{3}$. Let $Y\subset X=\operatorname{Spec}T^{(n)}$
be the reduced chain orbit closure of P-F24. Its grading
stabilizer at every nonzero point is exactly
$\boldsymbol\mu_M$. In the notation of P-F23 its leading
graded multiplicity is
\[
 e_Y=\lim_{z\to1}(1-z)
           \sum_{s\ge0}\dim\mathbb C[Y]_{Ms}z^s=1.
 \tag{P-F28a}
\]

Fix a primitive $M$th root $\zeta$, and let $g_\zeta$
denote its grading action. Normalize the cyclotomic
valuation above $2$ by $v(1-\zeta)=1$. The generic
conormal representation of $Y$ in $X$ has dimension
$m-n^2$, and
\[
 v\!\left(\det(1-g_\zeta\mid N^*_{Y/X,\eta_Y})\right)
       =3M-\binom{n+2}{2}.
 \tag{P-F28b}
\]
Consequently the local coefficient attached to this
component by the calculation of P-F23,
\[
 \mathcal C_n(\zeta):=
   \frac{e_Y}{\det(1-g_\zeta\mid N^*_{Y/X,\eta_Y})},
\]
is nonzero and satisfies
\[
 v(\mathcal C_n(\zeta))=\binom{n+2}{2}-3M.
 \tag{P-F28c}
\]
This is a local coefficient. No claim is made that it
equals a Laurent coefficient of the entire Hilbert series,
or dominates the contributions of other possible fixed components
by this proposition alone. P-F30--P-F31 below prove that it is
the actual normalized leading coefficient.
\end{proposition}

\begin{proof}
Put $\gamma=3/g$. By P-F24, the full group of linear
transformations fixing the chain form exactly is cyclic
of order $3M$, and its determinant image has order
$3M/g$. On the quotient, a matrix $B\in GL_n$ acts by
the grading parameter $(\det B)^{-\gamma}$, with the
point-action convention of (P-F26b). Thus the image
of this exact-form stabilizer in the grading torus has
order
\[
 \frac{3M/g}{\gcd(3M/g,\,3/g)}=M.
 \tag{P-F28d}
\]
This image is the entire grading stabilizer of the
quotient point. To see the converse inclusion, let a
grading parameter $u$ fix that point and choose a
matrix $B$ such that $(\det B)^{-\gamma}=u$.
The chain form and $B$ applied to it have closed
$SL_n$ orbits and the same quotient point. The
unique-closed-orbit property of the affine reductive
quotient therefore gives $S\in SL_n$ such that
$B\cdot F_n=S\cdot F_n$. The matrix $S^{-1}B$
fixes the form exactly and induces the same parameter
$u$. This proves the converse and hence the exact
stabilizer assertion. Grading translations give it
at every point of the same nonzero orbit.

Restrict homogeneous functions to the map
$\mathbb G_m\to Y$, $u\mapsto u\cdot p_n$.
A function of degree $d$ restricts to $P(p_n)u^d$.
The image of the coordinate ring is consequently
\[
 \mathbb C[Y]=\mathbb C[u^d:d\in S],
 \tag{P-F28e}
\]
where $S\subset\mathbb Z_{\ge0}$ is a finitely generated
additive semigroup containing a positive integer.
Finite generation follows by restricting a finite
set of homogeneous algebra generators; positivity
follows, for example, from the nonzero discriminant.
The common stabilizer of these monomials is
$\boldsymbol\mu_{\gcd S}$, so (P-F28d) gives
$\gcd S=M$.
This description also shows that there are no other
nonzero points of $Y$: the map $\mathbb A^1\to Y$
from (P-F28e) is finite and surjective, since $u$
is integral over the ring containing some $u^a$,
and its points with $u\ne0$ are exactly the grading
orbit.

After dividing its degrees by $M$, this semigroup
has greatest common divisor one and contains every
sufficiently large integer. Here is the elementary
argument for the latter assertion. Choose a positive
element $a$ of the semigroup. The residues of the
semigroup elements generate all of the finite group
$\mathbb Z/a\mathbb Z$, because their greatest common
divisor with $a$ is one. A subsemigroup of a finite
group is a subgroup: the inverse of an element of
order $b$ is its $(b-1)$-fold sum. Hence there is a
semigroup element in every residue class modulo $a$.
Adding nonnegative multiples of $a$ to these finitely
many representatives proves the assertion. Each
nonzero graded piece in (P-F28e) is one-dimensional.
The regraded Hilbert series therefore differs from
$1/(1-z)$ by a polynomial. This proves $e_Y=1$.

We next compute the conormal valuation. Write
$r=n-1$, so $M=2^r$. For any nonzero residue
$b\bmod M$, one has
\[
 v(1-\zeta^b)=2^{v_2(b)},
 \qquad 0\le v_2(b)\le r-1.
 \tag{P-F28f}
\]
For completeness, the polynomial
$\Phi_{2^r}(1+Z)=(1+Z)^{2^{r-1}}+1$
is Eisenstein at $2$: its constant term is $2$,
and reduction modulo $2$ leaves only the leading
term. Thus $1-\zeta$ is a uniformizer of the
cyclotomic extension at $2$.
If $b=2^j c$ with $c$ odd, the quotient
$(1-\zeta^{2^j c})/(1-\zeta^{2^j})$ is a unit,
since its geometric sum is congruent to the odd
integer $c$ modulo the uniformizer; one may choose
$c$ positive in its residue class. Moreover
\[
 1-\zeta^{2^j}
   =\prod_{k=0}^{2^j-1}
       \left(1-\zeta^{\,1+(M/2^j)k}\right).
 \tag{P-F28g}
\]
Every exponent on the right is odd, and every factor
has valuation one. This gives (P-F28f).
In particular this valuation is unchanged on
replacing $\zeta$ by any primitive $M$th root,
or on replacing a character by its inverse.

Choose a primitive $M$th root $\xi$ so that the
diagonal chain automorphism
\[
 A=\operatorname{diag}
    (\xi,\xi^{-2},\xi^4,\ldots,\xi^{(-2)^{n-1}})
\]
induces $g_\zeta$ on the quotient. This is possible:
its determinant exponent
$S_0=(1-(-2)^n)/3$ and the integer $\gamma$ are
both odd, so the quotient parameter
$\xi^{-\gamma S_0}$ runs through all primitive
$M$th roots. Put $w_i=(-2)^i\bmod M$ for
$0\le i\le n-1$. The weights $w_0,\ldots,w_{n-2}$
are nonzero and have respective two-adic valuations
$0,\ldots,n-2$, while $w_{n-1}=0$.

By P-F24, the tangent space of $X$ at the chain point
is the quotient of the cubic representation by the
tangent space of the $SL_n$ orbit. Its fixed part is
one-dimensional, namely the tangent direction of
the grading orbit. Taking moving parts and then
duals therefore gives the conormal representation
of $Y$. The cubic space has dimension $m$, the
$SL_n$ orbit has dimension $n^2-1$, and the fixed
quotient direction has dimension one. The conormal
dimension is consequently $m-n^2$.
The same character data holds at the generic point
of $Y$, since its nonzero grading orbit is dense
and the grading torus commutes with $g_\zeta$.

It suffices to subtract the determinant valuation
of the moving adjoint representation from that
of the moving cubic representation. Duals do not
change these valuations by (P-F28f).
The off-diagonal adjoint weights are $w_i-w_j$
for $i\ne j$. If $i<j$, their two-adic valuation
is $i$, including the case $j=n-1$ of the zero
weight. The adjoint determinant valuation is thus
\[
 V_{\mathrm{ad}}
   =2\sum_{i=0}^{n-2}(n-1-i)2^i
   =4M-2n-2.
 \tag{P-F28h}
\]

For cubic monomials, fix the smallest index $i$
of a nonzero weight that occurs. There are three
cases according to its multiplicity. If it occurs
once, the other two indices are greater than $i$,
and there are $\binom{n-i}{2}$ unordered choices
with repetition. The weight sum has two-adic
valuation $i$. If it occurs three times, there is
one monomial, also with valuation $i$.
If it occurs twice, the remaining index is some
$j>i$. The choice $j=i+1$ is a fixed chain
monomial and supplies no moving factor. Each
$j\ge i+2$ instead gives weight $2w_i+w_j$ of
valuation $i+1$; there are $n-i-2$ such choices
for $0\le i\le n-3$. The final all-zero-weight
cube is fixed. These cases exhaust the cubic
monomials, so
\begin{align*}
 V_{\mathrm{cub}}
 &=\sum_{i=0}^{n-2}
        \left[\binom{n-i}{2}+1\right]2^i
      +\sum_{i=0}^{n-3}(n-i-2)2^{i+1}\\
 &=7M-\frac{n^2+7n+6}{2}.
 \tag{P-F28i}
\end{align*}
For the finite sums used here, the identities are
\[
 \sum_{i=0}^{n-2}(n-1-i)2^i=2^n-n-1,\qquad
 \sum_{i=0}^{n-2}\binom{n-i}{2}2^i
    =2^{n+1}-\frac{n^2+3n+4}{2}.
 \tag{P-F28j}
\]
They follow by induction on $n$: the left sides
are twice their values at $n-1$, plus respectively
$n-1$ and $\binom n2$, and both identities begin
at $n=2$. The second sum in (P-F28i) is twice the
first sum of (P-F28j) with $n$ replaced by $n-1$.

Subtracting (P-F28h) from (P-F28i) yields
\[
 V_{\mathrm{cub}}-V_{\mathrm{ad}}
     =3M-\frac{n^2+3n+2}{2}
     =3M-\binom{n+2}{2}.
 \tag{P-F28k}
\]
This is (P-F28b). Each conormal eigenvalue is
nontrivial by P-F24, so its determinant factor
is nonzero. Combining it with $e_Y=1$ proves
(P-F28c).
\end{proof}

\paragraph{Role in the subsequent global argument.}
The local calculation does not compare possible boundary
components or establish their absence. P-F30 supplies
that separate support classification, and P-F31 then
verifies the global hypotheses of P-F23. Only at that
stage does the local coefficient computed here become
the actual normalized leading coefficient of the whole
Hilbert series. Thus the proof of the global pole uses
classification of the fixed locus, without an assumed
inequality between cyclotomic valuations of different
components.



% Source: branches/plethysm/sections/positivity_full_cubic_support.tex
% Reproduced inherited proof; finding P-F30.
\subsection{The full cubic support theorem}

\begin{theorem}[P-F30: extremal two-power determinant character forces the chain]
\label{thm:fullcubicsupport}
Let $n\ge2$, let $F\in\mathbb C[x_1,\ldots,x_n]$ be
a cubic, and let $B$ be the matrix of exponent vectors
of its nonzero monomials. Suppose
$\mathbf1=(1,\ldots,1)^t\in\operatorname{im}_{\mathbb Q}B$.
If a diagonal matrix $A$ of finite two-power order
fixes $F$ exactly, then
\[
 \operatorname{ord}(\det A)\le2^{n-1}.
 \tag{P-F30a}
\]
If equality holds, then, after a permutation of variables,
the support of $F$ is exactly
\[
 x_1^2x_2,\ x_2^2x_3,\ldots,\ x_{n-1}^2x_n,\ x_n^3.
 \tag{P-F30b}
\]
Its coefficients can be made one by a diagonal substitution,
and $A$ itself has order $2^{n-1}$. In particular, the
form is smooth. The bound is attained by the chain.

The rational-span hypothesis holds for every
$SL_n$-semistable cubic. Thus this theorem applies
to every eigenbasis of the finite lifts in P-F26;
it imposes no prior bound on the order of the lift.
\end{theorem}

\begin{proof}
The exponent character
\[
 \psi:\mathbb Z^n\longrightarrow\mathbb C^*,
 \qquad (a_i)\longmapsto\prod_i\alpha_i^{a_i},
\]
where $\alpha_i$ are the eigenvalues of $A$, annihilates
the support lattice and has two-power image. Its value
on $\mathbf1$ is $\det A$.
We will use the following elementary denominator fact
throughout. If $K$ is a selection of support columns
forming a basis of the rational support span, then
$Kz=\mathbf1$ has a unique rational solution $z$.
For any nonzero maximal square minor $K_I$, Cramer's
rule shows that the least positive integer $d$ for
which $dz$ is integral divides $|\det K_I|$.
Since $d\mathbf1=K(dz)$ is in the support lattice,
the order of $\det A$ divides $d$. Consequently
\[
 \operatorname{ord}(\det A)\le2^{v_2(\det K_I)}.
 \tag{P-F30c}
\]
This argument requires neither full rank of $B$
nor a bound on the orders of the individual $\alpha_i$.

Form the graph whose edges are the supported monomials
$x_i^2x_j$, with coefficient two at $i$ and one at $j$,
and whose loops are the supported cubes. Include all
$n$ vertices, even those isolated in this graph.
Let $G$ be the matrix of these graph-support columns.
Every other support column has the form
$e_i+e_j+e_k$ for three distinct vertices.

In a connected graph component on $v$ vertices,
a spanning-tree matrix has rank $v-1$. Thus the
graph-component matrix has rank $v$ or $v-1$.
This includes an isolated vertex: it has rank one
if its cube occurs, and rank zero otherwise.
Write $f$ for the number of full-rank components,
$c$ for the number of rank-deficient components,
and $v_i$ for their vertex counts. Then
\[
 \operatorname{rank}G=n-c.
 \tag{P-F30d}
\]

Choose a square support basis $C_i$ in each full-rank
component by adjoining one independent graph column
to a spanning tree. The unicyclic determinant
calculation in P-F27 gives
\[
 v_2(\det C_i)\le v_i-1.
 \tag{P-F30e}
\]
For precision, leaf expansion contributes a factor
one or two for each edge off the unique cycle.
A loop contributes the odd factor three.
A cycle of length $l\ge2$ contributes, up to sign,
$2^a+(-1)^{l-1}2^{l-a}$; if nonzero, its valuation
is at most $\lfloor(l-1)/2\rfloor$. These alternatives
give (P-F30e), including loops and double edges.

Choose a spanning-tree matrix $T_i$ in each deficient
component. Its signed cofactor vector has the form
\[
 2^{e_i}u_i,\qquad
 (u_i)_v=\varepsilon_i(-2)^{h_{iv}},
 \quad \varepsilon_i\in\{1,-1\},\quad
 \min_v h_{iv}=0.
 \tag{P-F30f}
\]
All exponents here are nonnegative integers.
Indeed every maximal tree minor is, up to sign,
a power of two, by successive leaf expansion.
The cofactor relation along an edge with
coefficients two and one is $2u_v+u_w=0$.
Thus adjacent heights differ by one and their
signs are opposite; connectivity gives the
single sign convention in (P-F30f).
For a single isolated deficient vertex, the
empty minor is one and we use $e_i=h_{iv}=0$.

Put
\[
 b_i=v_i-1-e_i,\qquad H_i=\max_v h_{iv}.
\]
Since a tree minor is the product of $v_i-1$
factors, each one or two, one has
\[
 0\le H_i\le b_i.
 \tag{P-F30g}
\]
We need one additional equality property: if
$H_i=b_i$, the vertex of height $H_i$ is unique.
To prove it, choose such a vertex as root.
Its deleted-row minor has absolute value
$2^{e_i+H_i}=2^{v_i-1}$. Expand this minor by
removing leaves toward that root. Every factor
must be two. Thus each edge has coefficient
two at its child and one at its parent.
The cofactor relation gives
$h_{\rm parent}=h_{\rm child}+1$.
Heights strictly decrease away from the root,
so no other vertex has its height.
The assertion is also immediate for an isolated vertex.

Normalize the entries in each of these cofactor
vectors by $2^{b_i}$. For sums over distinct
vertices of one deficient component, the following
bounds hold:
\[
 \left|\sum_{v\in E}\frac{(u_i)_v}{2^{b_i}}\right|
 \le
 \begin{cases}
  1,&|E|=1,\\
  5/4,&|E|=2,\\
  3/2,&|E|=3.
 \end{cases}
 \tag{P-F30h}
\]
Here a case is used only when the component has
enough distinct vertices. If $H_i<b_i$, every
entry has absolute value at most $1/2$, which
proves all three bounds. If $H_i=b_i$, there is
only one entry of absolute value one. Change
the overall sign to make this entry positive.
Every entry of absolute value $1/2$ is negative,
because its height is $H_i-1$; every remaining
entry has absolute value at most $1/4$.
A pair avoiding the top entry has absolute sum
at most one; a pair containing it has absolute
sum at most $5/4$. A triple avoiding it has
absolute sum at most $3/2$. For a triple containing
it, either the other two entries are both $-1/2$,
giving zero, exactly one is $-1/2$, giving absolute
sum at most $3/4$, or neither is, giving at most
$1+1/4+1/4=3/2$. This proves (P-F30h).

The chosen full-component bases and deficient-component
trees together form a rational basis for the columns
of $G$. Extend it to a basis for all support columns
by adjoining $q$ independent triple columns
$\beta_1,\ldots,\beta_q$. Then
\[
 0\le q\le c,\qquad
 r:=\operatorname{rank}B=n-c+q.
 \tag{P-F30i}
\]
If $q=0$, the graph subform itself satisfies the
rational-span hypothesis. P-F27 applies directly:
it gives (P-F30a), and equality gives the chain
eigenvalue list and the chain graph support.
That eigenvalue list admits no other cubic
monomials, so there can be no additional triple
monomials in $F$. This proves all the required
claims in the case $q=0$.
We therefore assume $q>0$ for the rest of the argument.

Define the $c\times q$ compressed matrix
\[
 D_{ij}=u_i^t(\beta_j|_i),
 \tag{P-F30j}
\]
where $\beta_j|_i$ is the restriction to the vertices
of deficient component $i$. The map on the entire
vertex space that sends a vector to these $c$
cofactor pairings is surjective, and its kernel
is precisely the rational graph-column span:
on a full component the graph columns span all
coordinates, while on a deficient component the
kernel of $u_i^t$ is the tree-column span.
Consequently $D$ has rank $q$. Choose $q$
deficient-component indices $J$ for which the
square matrix $D_J$ has nonzero determinant.

Let $K$ be the $n\times r$ original support-basis
matrix just selected. We now specify an actual
nonzero $r\times r$ minor of it. Keep every row
in each full component and in each deficient
component indexed by $J$. In every other deficient
component, delete one row at a vertex of height zero,
and keep the remaining rows. The resulting row
set $I$ has size $n-(c-q)=r$. Its determinant is
\[
 \det K_I
   =\pm\left(\prod_{\text{full }i}\det C_i\right)
          2^{\sum_{\text{deficient }i}e_i}
          \det D_J.
 \tag{P-F30k}
\]
We justify this factorization explicitly.
In a full component the graph basis is square;
eliminating the triple-column entries there
using its graph columns and expanding factors
out $\det C_i$. In a deficient component not
in $J$, the retained tree block is square.
Its deleted row had height zero, so its determinant
is $\pm2^{e_i}$; eliminate the triple entries
there and factor out that determinant as well.
These column eliminations do not alter other
components, since every graph column is supported
in its own component.
The remaining matrix consists of the tall
tree blocks for $i\in J$ and $q$ triple columns.
Expansion by the maximal minors of each tree
block leaves one row pairing for each such
component. Its signed cofactor vector is
$2^{e_i}u_i$. The remaining $q\times q$
determinant is exactly the pairings (P-F30j),
giving the remaining factors
$2^{\sum_{i\in J}e_i}\det D_J$.
This proves (P-F30k), and proves that the
specified original minor is nonzero.

Divide row $i$ of $D_J$ by $2^{b_i}$.
Each resulting column is formed from at most
three distinct vertices of the original graph.
If they belong to three selected components,
its squared Euclidean norm is at most three.
If two lie in one selected component and one
in another, (P-F30h) gives squared norm at most
\[
 (5/4)^2+1=41/16<3.
\]
If all three lie in one component, the squared
norm is at most $(3/2)^2<3$. Ignoring vertices
in unselected or full components only leaves
cases with fewer than three contributing vertices,
which satisfy the same bound. Every normalized
column therefore has Euclidean norm at most
$\sqrt3$. Hadamard's determinant inequality yields
\[
 0<|\det D_J|
       \le 2^{\sum_{i\in J}b_i}3^{q/2}.
 \tag{P-F30l}
\]

Suppose first that $q<c$, or that $q=c$ and $f>0$.
Using (P-F30e), (P-F30k), and (P-F30l) gives
\begin{align*}
 v_2(\det K_I)
 &\le \sum_{\text{full }i}(v_i-1)
      +\sum_{\text{deficient }i}e_i
      +\sum_{i\in J}b_i+\frac q2\log_2 3\\
 &=n-f-c-\sum_{\substack{\text{deficient }i\\i\notin J}}b_i
              +\frac q2\log_2 3\\
 &<n-f-c+q.
 \tag{P-F30m}
\end{align*}
The last inequality is strict because $q>0$
and $\sqrt3<2$; the omitted $b_i$ are nonnegative.
If $q<c$, its right side is at most $n-1$.
If $q=c$ and $f>0$, it is likewise at most $n-1$.
The valuation is an integer, so in both cases
\[
 v_2(\det K_I)\le n-2.
 \tag{P-F30n}
\]

The only case remaining is $q=c$ and $f=0$.
Here $K$ is a square $n\times n$ support matrix,
and $I$ is its full row set. Every original
support column has sum three. Hence
$\mathbf1^tK=3\mathbf1^t$, and reduction modulo
three shows that $3$ divides its nonzero determinant.
By (P-F30k), which now has no full-component
factors, $3$ also divides $\det D$.
The two-part of this determinant is consequently
at most $|\det D|/3$. Since
$\sum_i(e_i+b_i)=\sum_i(v_i-1)=n-c$,
we obtain
\begin{align*}
 2^{v_2(\det K)}
 &=2^{\sum_i e_i}\,2^{v_2(\det D)}\\
 &\le 2^{\sum_i e_i}\frac{|\det D|}{3}\\
 &\le 2^{n-c}3^{c/2-1}
  <\frac{2^n}{3}
  <2^{n-1}.
 \tag{P-F30o}
\end{align*}
Here $c=q>0$ and $3^{c/2}<2^c$ justify the
first strict inequality. Thus (P-F30n) holds
in this case as well.

Combining (P-F30c) and (P-F30n) proves the
stronger bound
$\operatorname{ord}(\det A)\le2^{n-2}$
whenever $q>0$. All extremal cases therefore
occur in the already treated case $q=0$,
where P-F27 proves precisely (P-F30b), the
coefficient normalization, and the exact order
of $A$. Conversely the chain is fixed by
the diagonal matrix with eigenvalues
$\xi,\xi^{-2},\ldots,\xi^{(-2)^{n-1}}$
for a primitive $2^{n-1}$th root $\xi$.
Its determinant exponent
$1-2+4-\cdots+(-2)^{n-1}$ is odd, so its
determinant has order $2^{n-1}$. The bound
is therefore attained.

Finally, P-F27's coefficient-monomial argument
shows that any positive-degree $SL_n$ invariant
nonzero at $F$ expresses $\mathbf1$ as a
nonnegative rational combination of the
support exponents. This supplies the
rational-span hypothesis for every semistable
cubic, in every coordinate basis, and completes
the proof.
\end{proof}

\paragraph{The structural conclusion.}
The theorem closes the mixed-support case left by
P-F27 and P-F29. It uses the actual determinant
character, rather than identifying that character's
order with an arbitrary lift order. Its proof
allows rank-deficient exponent matrices and uses
an explicitly constructed maximal minor in those
cases. Combining it with P-F26 rules out
nonzero singular representatives in the
$\boldsymbol\mu_{2^{n-1}}$ fixed quotient locus.
The ensuing global fixed-locus and Hilbert-pole
conclusions are stated separately.



% Source: branches/plethysm/sections/positivity_actual_cubic_pole.tex
% Reproduced inherited proof; finding P-F31.
\subsection{Actual poles at roots of exponential order in ordinary plethysm}

\begin{theorem}[P-F31: the chain pole refutes the positive-form clause of PH3]
\label{thm:actualcubicpole}
Let $n\ge4$, put
\[
 g=\gcd(n,3),\qquad d_0=n/g,\qquad
 M=2^{n-1},\qquad m=\binom{n+2}{3},
\]
and use the ordinary plethysm ray of P-F19,
\[
 \mu=(3),\qquad \lambda=(d_0),\qquad
 \pi=((3/g)^n).
\]
Write
\[
 f_n(s)=\dim\operatorname{Hom}_{GL_n}
 \left((\det)^{3s/g},\operatorname{Sym}^{d_0s}
                    (\operatorname{Sym}^3\mathbb C^n)\right),
 \qquad
 H_n(t)=\sum_{s\ge0}f_n(s)t^s.
 \tag{P-F31a}
\]
Let $T^{(n)}$ be its canonical Cartan multiplicity algebra,
$X=\operatorname{Spec}T^{(n)}$, and let $Y\subset X$ be the
reduced closure of the chain grading orbit. Then:

\begin{enumerate}
\item The support of the fixed scheme $X^{\boldsymbol\mu_M}$
is exactly $Y$. Away from the graded origin this fixed scheme
is the smooth chain orbit, and $X$ is smooth there. Possible
nilpotent or embedded structure at the origin is not excluded.
\item The ordinary Hilbert series $H_n(t)$ has a simple pole
at every primitive $M$th root of unity $\zeta$. More precisely,
with $g_\zeta$ the grading action,
\[
 \lim_{t\to\zeta}(1-t^M)H_n(t)
   =\frac{1}{\det(1-g_\zeta\mid N^*_{Y/X,\eta_Y})}
   \ne0.
 \tag{P-F31b}
\]
\item In every expression
\[
 H_n(t)=\frac{Q(t)}{\prod_{j=1}^{r}(1-t^{a_j})^{b_j}},
 \qquad Q(t)\in\mathbb C[t],\quad a_j,b_j\in\mathbb Z_{>0},
 \tag{P-F31c}
\]
some denominator weight $a_j$ is divisible by $M$. Thus every
such expression, in particular every positive rational form,
has maximum denominator weight at least $2^{n-1}$.
\end{enumerate}
Since the source parameter is
\[
 h=\dim GL_m+\operatorname{ht}(\lambda)+\operatorname{ht}(\pi)
   =\binom{n+2}{3}^{\!2}+1+n=O(n^6),
 \tag{P-F31d}
\]
the first, positive-form existence clause of GCT VI,
Hypothesis~1.6.7 is false for this ordinary irreducible-inner
plethysm family. The stronger canonical parameter clause was
already refuted by P-F19. No superpolynomial positivity or
saturation index is asserted here.
\end{theorem}

\begin{proof}
By P-F19 the actual canonical algebra, including multiplication,
is the entire cubic invariant ring
\[
 T^{(n)}\simeq
   \bigoplus_{s\ge0}
      \mathbb C[(\operatorname{Sym}^3\mathbb C^n)^*]^{SL_n}_{d_0s},
 \tag{P-F31e}
\]
with the coefficient grading divided by $d_0$. Scalar central
characters force every nonzero invariant degree to be a multiple
of $d_0$, so this is a regrading of the whole ring, rather than
a smaller subring. In particular $T^{(n)}$ is a finitely generated
positively graded complex domain with degree-zero part $\mathbb C$.
Its Hilbert series is exactly (P-F31a).

Consider a nonzero closed point $p$ of the fixed scheme.
P-F26 gives a representative cubic $F$ with closed $SL_n$ orbit
and an exact-form stabilizer element $A$ of finite two-power order
whose determinant has order $M$. Since $p$ is not the graded
origin, some positive-degree invariant is nonzero at $F$.
The supported cubic exponent vectors therefore span
$\mathbf1$ rationally, by the coefficient-monomial argument
in P-F27. The complete support theorem P-F30 applies to this
$F$ and $A$, without assuming that the lift has order $M$
or that the support is a graph. Its equality case identifies
$F$, after a permutation and diagonal rescaling in an eigenbasis,
with the smooth chain cubic $F_n$.

Consequently $p$ lies on the $GL_n$ orbit of the chain quotient
point. On the quotient this is exactly the grading orbit, by
(P-F26b) and the surjectivity of the character
$z\mapsto z^{-3/g}$ of $\mathbb G_m$. Conversely P-F24 shows
that the entire nonzero chain grading orbit is fixed by
$\boldsymbol\mu_M$.

The graded origin is also fixed and is in the closure of this
orbit. Indeed the map $u\mapsto u\cdot p_n$ extends from
$\mathbb G_m$ to $\mathbb A^1$: each positive-degree homogeneous
function restricts to a constant times a positive power of $u$,
and thus takes value zero at $u=0$. It follows that the closed
points of the fixed scheme are exactly those of $Y$. A
finite-type closed subset over $\mathbb C$ is determined by
its closed points, so the support of $X^{\boldsymbol\mu_M}$
is $Y$. This is an irreducible curve.

By P-F24 the ambient quotient $X$ and the fixed scheme are
smooth at every point of the nonzero chain orbit, and the fixed
scheme there has dimension one. Since these are all its
nonzero closed points, any nonreduced structure or embedded
associated point is confined to the origin. It is unnecessary
to assert reducedness at that origin: the fixed scheme has a
unique maximal-dimensional component $Y$, and $X$ is smooth
at its generic point. These are precisely the hypotheses
needed in P-F23.

Apply P-F23 with dimension $\delta=1$. It gives
\[
 \lim_{t\to\zeta}(1-t^M)H_n(t)
   =\frac{e_Y}{\det(1-g_\zeta\mid N^*_{Y/X,\eta_Y})}.
 \tag{P-F31f}
\]
All conormal eigenvalues are nontrivial. The exact grading
stabilizer and semigroup calculation in P-F28 give $e_Y=1$.
Thus the right hand side is nonzero, proving (P-F31b).
The factor $1-t^M$ has a simple zero at $\zeta$, so this is
an actual pole of order one of $H_n$.

This application includes the possible structure at the
origin. In the Koszul proof of P-F23, terms supported at
that zero-dimensional point have finite-dimensional graded
homology and hence polynomial Hilbert series. They vanish
after multiplication by $1-t^M$ in the limit. They cannot
cancel the one-dimensional generic contribution. No global
smoothness or reducedness of the fixed scheme was used.

For additional precision, P-F28 now computes the actual
normalized leading coefficient in (P-F31b). If the cyclotomic
valuation above two is normalized by $v(1-\zeta)=1$, it has
valuation
\[
 \binom{n+2}{2}-3\,2^{n-1}.
 \tag{P-F31g}
\]
The ordinary residue of $H_n(t)\,dt$ at $t=\zeta$ is
$-\zeta/M$ times the right hand side of (P-F31b), because
$(1-t^M)'|_{t=\zeta}=-M\zeta^{-1}$. The nonvanishing
in (P-F31b), rather than any comparison of phases, is
what is needed below.

Finally suppose (P-F31c) is any rational expression of
the indicated form. Its numerator is regular at $\zeta$.
If none of its denominator factors vanished there, the
expression would also be regular there, contradicting
the proved pole. Hence $\zeta^{a_j}=1$ for some $j$.
Primitivity of $\zeta$ is exactly the assertion that
$M\mid a_j$. This argument does not require a nonnegative
numerator, a prescribed number of denominator factors,
or a presentation arising from an h.s.o.p.

The exact source definition of a positive form is
GCT VI, Definition~1.2.7, printed p.~15, equation~(1.4);
its modular index is the maximum of the weights $a_j$.
Hypothesis~1.6.7, printed p.~25, first asks for the
existence of such a form with modular index polynomial
in $\dim G+\operatorname{ht}(\lambda)+\operatorname{ht}(\pi)$,
and then makes the stronger assertion for its canonical
minimal parameter form. The parameter in this family is
(P-F31d), whereas every allowed form has a weight at
least $2^{n-1}$. For any fixed polynomial in $h$, its
value at (P-F31d) is bounded by a fixed polynomial in
$n$, which is eventually smaller than $2^{n-1}$.
Therefore even the first existence clause fails.
\end{proof}

\paragraph{Consequences and remaining scope.}
The result closes the ordinary-pole gate that P-F19 alone
left open. It uses the full cubic support theorem and the
unique generically smooth fixed component; it does not
infer an ordinary pole merely from an orbit degree period,
a large parameter degree, a universal denominator, or
cohomological purity. The cancellation examples P-F21--P-F22
remain valid and explain why those shorter arguments
would not have sufficed.

Every quasipolynomial period of $f_n$ is divisible by $M$.
Indeed a quasipolynomial of period $L$ has a generating
function with denominator a power of $1-t^L$: sum its
constituent polynomials over each arithmetic progression,
using the elementary generating functions for powers.
The pole in (P-F31b) forces $M\mid L$. This does not say
that $M$ is the entire minimal period or classify its
other prime-power factors.

The implication from bounded positive-form weights to
PH2 in P-F12 remains correct, but failure of its premise
does not refute PH2 or SH. Nor do large denominator
weights alone establish a counting-complexity lower bound,
an obstruction search algorithm, or a failure of every
possible revised geometric formulation. Here the inner
shape is the single irreducible cubic shape $(3)$, and
only the number of variables grows; reducible torus
embeddings and determinant twists are not responsible
for the obstruction. No historical priority claim is
made for the support classification or the resulting
Hilbert-pole consequence.

The exact primary source for the conjecture and its
positive-form convention is
\url{https://arxiv.org/pdf/0704.0229},
Definition~1.2.7, p.~15, and Hypothesis~1.6.7, p.~25.


\section{The maximal dyadic wave and its recursion}

% Source: branches/plethysm/sections/paired_wave_chain.tex
\subsection{An isolated dyadic wave and an actual-coefficient index test}

The simple poles of P-F31 permit a strict-index test using a specified
cyclotomic trace and an elementary bound on an actual coefficient.
The reduction is reproduced with its complete proof.  The conditions
explicitly stated as implications here are discharged by PNEW-005 below.

\begin{proposition}[P-F56: maximal two-part poles and paired constituents]
\label{prop:cubicpairedwave}
Use the ordinary cubic multiplicity algebra $T^{(n)}$ and its
Hilbert function $f_n$ from P-F31, with $n\ge4$, and put
\[
 M=2^{n-1},\qquad g=\gcd(n,3),\qquad d_0=n/g,
 \qquad m=\binom{n+2}{3}.
\]
Every pole of $H_n(t)$ is a root of unity whose order has
two-part at most $M$. Among roots whose order has two-part exactly
$M$, the only poles are the primitive $M$th roots, and each of
those poles is simple. More precisely, if $\eta$ is a root of
unity of order $q$ with $v_2(q)\ge n-1$, then
\[
 \eta\text{ is a pole of }H_n
       \quad\Longleftrightarrow\quad q=M;
 \qquad q=M\Longrightarrow\text{the pole is simple}.
 \tag{P-F56a}
\]

Let $L$ be the least quasipolynomial period of $f_n$, and write
$f_n(s)=q_r(s)$ for $s\equiv r\pmod L$, with
$0\le r<L$. Then
\[
                     v_2(L)=n-1.                            \tag{P-F56b}
\]
Fix $\zeta=e^{2\pi i/M}$ and use the following actual-residue normalization:
\[
 \mathcal C_n(\zeta)=\lim_{t\to\zeta}(1-t^M)H_n(t),
 \qquad A_n(\zeta)=\mathcal C_n(\zeta)/M.
\]
For each integer $r$ define the rational, $M$-periodic wave
\[
 \begin{split}
 W_n(r)
   &=\sum_{\substack{1\le a<M\\a\text{ odd}}}
                      A_n(\zeta^a)\zeta^{-ar}\\
   &=\frac1M\operatorname{Tr}_{\mathbb Q(\zeta)/\mathbb Q}
                  \bigl(\mathcal C_n(\zeta)\zeta^{-r}\bigr).
 \end{split}                                                 \tag{P-F56c}
\]
For every residue $r$ and every complex argument $X$ one has
the exact polynomial identity
\[
               q_r(X)-q_{r+L/2}(X)=2W_n(r),                 \tag{P-F56d}
\]
where subscripts are read modulo $L$. In particular,
\[
 \begin{split}
 f_n\bigl(p(f_n)\bigr)&\ge2W_n\bigl(p(f_n)\bigr),\\
 f_n\bigl(s(f_n)+1\bigr)&\ge2W_n\bigl(s(f_n)+1\bigr).
 \end{split}                                                 \tag{P-F56e}
\]
These are necessary conditions for the literal indices of
GCT VI, Definitions~1.2.2 and 1.2.4, printed pp.~13--14.
They require no bound on the average polynomial $B_n$ or on
Fourier terms of smaller two-part.
More generally, the inequality $f_n(x)\ge2W_n(x)$ holds
at every integer $x\ge p(f_n)$ and at every integer
$x\ge s(f_n)+1$. Thus a single integer $x\ge0$ satisfying
\[
          2W_n(x)>f_n(x)
       \quad\Longrightarrow\quad p(f_n)>x,
                                  \quad s(f_n)\ge x.
 \tag{P-F56e'}
\]

For every fixed real $K\ge0$ and every integer $0\le x\le n^K$,
\[
 f_n(x)\le\binom{m+d_0x-1}{m-1}
     \le\bigl(2n^{\max\{3,K+1\}}\bigr)^{n^3}.
 \tag{P-F56f}
\]
Consequently the following would suffice to refute both narrow
polynomial-index hypotheses on this cubic family: there exist
$\epsilon>0$ and an unbounded set $\mathcal N$ of dimensions
such that, for every fixed $K\ge0$ and every sufficiently large
$n\in\mathcal N$, there is an integer $x$ satisfying
\[
 \lceil n^K\rceil\le x\le\lceil n^K\rceil+2,
                         \qquad W_n(x)\ge2^{\epsilon n^4}.
 \tag{P-F56g}
\]
Under this additional condition, both $p(f_n)$ and $s(f_n)$
exceed every fixed polynomial in $n$ along $\mathcal N$.
PNEW-005 below supplies a witness in a window of width at most nine,
which is sufficient for the same conclusion. In the present conditional criterion the wave can have
either sign a priori, and the absolute value of one Fourier
coefficient alone does not establish (P-F56g).
The window of three consecutive integers is a sufficient
target, not a proved oscillation pattern; any bounded-width
window, or suitable polynomially bounded witnesses tending
beyond every fixed polynomial threshold, would work as well.
\end{proposition}

\begin{proof}
We first prove a fixed-point criterion for absence of a Hilbert
pole, including possible nonreduced structure at the origin.
Let $R$ be any finitely generated positively graded complex
algebra with $R_0=\mathbb C$, and let $\eta$ have order $q$.
The ideal defining its grading fixed scheme is
\[
                    I_q=(R_d:q\nmid d).                     \tag{P-F56h}
\]
Indeed a homogeneous element of degree $d$ changes by
the scalar $\eta^d$, so its difference from its transform
is zero exactly when $q\mid d$.
Choose homogeneous ideal generators $a_1,\ldots,a_l$ of
$I_q$, each of degree $e_i$ not divisible by $q$.
Such a finite choice is possible by Noetherianity, starting
with the indicated homogeneous generating set.
If the fixed scheme is supported only at the graded origin,
then $\sqrt{I_q}=R_{>0}$ and $R/I_q$ is a finite-dimensional
complex algebra. To see this directly, a power of the finitely
generated ideal $R_{>0}$ lies in $I_q$, and only finitely many
monomials of bounded length in homogeneous algebra generators
survive.

The homology modules $N_j$ of the graded Koszul complex
$K_R(a_1,\ldots,a_l)$ are finite $R$-modules annihilated by
$I_q$: multiplication by each $a_i$ is null-homotopic on
the Koszul complex. Thus all $N_j$ are finite-dimensional.
The degreewise Euler identity gives
\[
 H_R(t)\prod_{i=1}^l(1-t^{e_i})
                    =\sum_j(-1)^j\operatorname{Hilb}_{N_j}(t).
 \tag{P-F56i}
\]
The right side is a polynomial and every factor on the left
is nonzero at $t=\eta$. Hence $H_R(t)$ is regular at $\eta$.
This is the zero-dimensional case of the Koszul calculation
in P-F23; reducedness of the fixed scheme is unnecessary.

Apply this criterion to $R=T^{(n)}$ and $X=\operatorname{Spec}R$.
Let $q=2^ab$ with $b$ odd and $a\ge n-1$.
Suppose a nonzero closed point $p\in X$ is fixed by a primitive
$q$th root $\eta$. It is then fixed by
$\eta_2=\eta^b$, which has order $2^a$.
We need the closed-orbit lift for this pure two-primary
subgroup, with no assumption about the odd part of $q$.
Here are the details extending the first part of P-F26.

Put $\gamma=3/g$, an odd integer. The canonical grading
character on quotient points is
$(\det B)^{-\gamma}$ for $B\in GL_n$, as in (P-F26b).
There is a primitive $2^a$th root $\theta$ such that
$\theta^{-\gamma}=\eta_2$.
Choose $B\in GL_n$ with $\det B=\theta$, and choose
a cubic $F$ in the unique closed $SL_n$ orbit over $p$.
Both $F$ and $B\cdot F$ have closed orbits and the same
quotient point. Therefore $B\cdot F=S\cdot F$ for some
$S\in SL_n$, and $C=S^{-1}B$ fixes $F$ exactly with
$\det C=\theta$.
The unique closed orbit in each affine reductive quotient
fiber is the standard input already used in P-F26; an exact
primary locator is Luna, \emph{Slices \'etales},
Section~I.1, printed p.~84:
\url{https://www.numdam.org/item/10.24033/msmf.110.pdf}.

The semisimple part $C_s$ also fixes $F$ and has determinant
$\theta$. Indeed in the commuting Jordan decomposition,
an eigenspace of $C_s$ with eigenvalue different from one
cannot contain a fixed vector of $C_sC_u$, because the
unipotent part $C_u$ has all eigenvalues one.
Diagonalize $C_s$ and let $\Gamma\subset\mathbb C^*$ be
the finitely generated group generated by its eigenvalues.
Project $\Gamma$ to its finite two-primary torsion subgroup,
using a decomposition into its finite cyclic torsion subgroup
and a free complement. This projection fixes $\theta$.
Applying it to every eigenvalue preserves every monomial
relation expressing $C_sF=F$.
The resulting diagonal matrix $A$ fixes $F$, has finite
two-power order, and has determinant of order exactly $2^a$.

Since $p$ is nonzero, a positive-degree invariant is nonzero
at $F$. In any eigenbasis, one of its nonzero coefficient
monomials uses only monomials supported in $F$.
Invariance under the diagonal $SL_n$ torus forces the
corresponding nonzero nonnegative integral combination of supported
cubic exponent vectors to be a positive rational multiple
of $\mathbf1$. Thus $\mathbf1$ belongs to the rational
span of the support, precisely the hypothesis of P-F30.
Theorem P-F30 gives
\[
                       2^a=\operatorname{ord}(\det A)\le M.
 \tag{P-F56j}
\]
If $a>n-1$ this is a contradiction, so no nonzero fixed
point exists. If $a=n-1$, equality in P-F30 makes $F$
equivalent to the chain cubic, hence $p$ lies on its
nonzero grading orbit. The exact stabilizer of every
point of that orbit is $\mu_M$ by P-F28, including the
converse lift argument in (P-F28d).
The original root $\eta$ can fix $p$ only if $\eta\in\mu_M$.
Since its order is $q=Mb$, this requires $b=1$.

It follows that for every $q$ in (P-F56a) except $q=M$
the fixed scheme has no nonzero closed point. Its support
is consequently just the graded origin, because a closed
subset of an affine scheme of finite type over $\mathbb C$
is determined by its closed points. Equations
(P-F56h)--(P-F56i) show that there is no Hilbert pole there.
For $q=M$, P-F31 proves an actual simple pole at every
primitive root. This proves all of (P-F56a).

For completeness, the least quasipolynomial period is the
least common multiple of the orders of the Hilbert poles.
Indeed the finite Fourier transforms are
\[
 P_\rho(X)=\frac1L\sum_{r=0}^{L-1}\rho^r q_r(X),\qquad
 q_r(X)=\sum_{\rho^L=1}\rho^{-r}P_\rho(X).                 \tag{P-F56k}
\]
The generating function of a nonzero polynomial times
$\rho^{-s}$ has its only pole at $t=\rho$, with order
one more than its polynomial degree. Thus $P_\rho\ne0$
if and only if $\rho$ is a pole of $H_n$.
Let $L_0$ be the least common multiple of those root orders.
The second expression in (P-F56k) is periodic in $r$
with period $L_0$, so it exhibits $L_0$ as a quasipolynomial
period. Conversely $L_0\mid L$. Minimality gives $L=L_0$.
Now (P-F56a) and the primitive $M$th poles give (P-F56b).

In particular $L/M$ is odd. If $\rho$ is any Hilbert pole
other than a primitive $M$th root, its order has two-part
at most $M/2$ and divides $L/2$. Its contribution in the
difference of (P-F56k) at $r$ and $r+L/2$ therefore
vanishes. For $\rho=\zeta^a$ with $a$ odd,
$\rho^{-L/2}=-1$.
At a simple pole $\rho$, the term with its principal part is
$A_n(\rho)/(1-t/\rho)$: indeed
$\lim_{t\to\rho}(1-t^M)/(1-t/\rho)=M$.
Its coefficient at degree $s$ is $A_n(\rho)\rho^{-s}$,
so $P_{\zeta^a}(X)=A_n(\zeta^a)$, independent of $X$.
Subtracting the two Fourier formulas proves (P-F56d).
It also shows that every nonconstant coefficient of the
constituent polynomials has period dividing $L/2$.
The rational Hilbert series is defined over $\mathbb Q$,
so Galois conjugation sends $\mathcal C_n(\zeta)$ to
$\mathcal C_n(\zeta^a)$. This proves the trace expression
and rationality in (P-F56c). Its periodicity is immediate;
also $W_n(r+M/2)=-W_n(r)$.

After the coefficient-positivity shift $p=p(f_n)$, all
constituents evaluated at $p$ are nonnegative, since these
values are their shifted constant coefficients.
Take the residue $r\equiv p\pmod L$ in (P-F56d).
Its first term at $X=p$ is the actual value $f_n(p)$,
and its second term is nonnegative.
After a strict-saturation shift $s=s(f_n)$, every nonzero
constituent at $s+1$ is positive, and every identically
zero constituent has value zero. Taking
$r\equiv s+1\pmod L$ proves the second inequality of
(P-F56e). This uses exactly the source conventions, not
positivity only on each constituent's own progression.
These tests persist at every later integer argument.
For positivity, each polynomial $q_r(p+X)$ has nonnegative
coefficients and hence is nonnegative for all real $X\ge0$.
For saturation, the definition directly gives
$q_r(s+j)\ge0$ for every integer $j\ge1$.
Applying (P-F56d) at the residue of any such later argument
proves (P-F56e'). No derivative positivity is being inferred
from strict saturation.

Finally the multiplicity $f_n(x)$ is bounded by the
dimension of the ambient symmetric power of an
$m$-dimensional space. Since $m\le n^3$ and
$d_0x\le n^{K+1}$,
\[
 \binom{m+d_0x-1}{m-1}
      \le(m+d_0x)^{m-1}
      \le\bigl(2n^{\max\{3,K+1\}}\bigr)^{n^3},
\]
including $x=0$. This proves (P-F56f).
Its base-two logarithm is $O_K(n^3\log n)$.
The witness in (P-F56g) is at most $n^{K+1}$ for all
sufficiently large $n$. Apply (P-F56f) with $K+1$.
Its actual coefficient is smaller than $2W_n(x)$ for
sufficiently large $n\in\mathcal N$, so (P-F56e') gives
$p(f_n)>x\ge n^K$ and $s(f_n)\ge x\ge n^K$.
Applying this with a larger fixed exponent proves that
neither index is bounded by any polynomial in $n$.
The GCT parameter is $h=m^2+1+n=\Theta(n^6)$ by P-F31,
so a polynomial bound in that parameter would also fail.
This establishes the conditional criterion. The analytic estimates
and bounded-window witness are proved in PNEW-005 below.\hfill$\square$
\end{proof}


\paragraph{An exact positive multiset for the chain normal representation.}
There is also a numerator-free product specifying every residue
in the remaining estimate. Regard
$w_0,\ldots,w_{n-1}$ as residues modulo $M$, with the
$w_i=(-2)^i$ for $i<n-1$ and $w_{n-1}=0$, and form the multiset
\[
 \begin{gathered}
 \mathcal E_n=
 \{\{0,1,k\}:2\le k<n\}\cup\{\{0,2,3\}\},\\
 \mathcal S_n=\{w_i:0\le i\le n-2\}\ \sqcup\
 \{w_i+w_j+w_k:0\le i<j<k<n,\
       \{i,j,k\}\notin\mathcal E_n\}.
 \end{gathered}
 \tag{P-F56u}
\]
Repeated weights are retained. Every weight is nonzero modulo
$M$, and $|\mathcal S_n|=D=\binom n3$.
With the actual Hilbert-residue convention, the exact identity is
\[
 \mathcal C_n(\zeta)=
                 \prod_{b\in\mathcal S_n}(1-\zeta^{gb})^{-1}.
 \tag{P-F56v}
\]
This is an equality of complex numbers, including their phases,
not only an absolute-value identity.

To prove it, put $\eta=\zeta^g$ and let
$\alpha_i=\eta^{w_i}$ be the natural eigenvalues of the diagonal
chain automorphism. Write $V$ for their representation.
Its action on polynomial coordinate functions is the usual
$GL_n$ action on $\operatorname{Sym}^3 V$.
On an invariant of coefficient degree $d_0s$ the character is
$(\det)^{\gamma s}$, where $\gamma=3/g$.
Since $3\sum_i w_i\equiv1\pmod M$, one has
$(\prod_i\alpha_i)^\gamma=\zeta$.
Thus its action on the canonical multiplicity algebra is precisely
the ring automorphism multiplying degree $s$ by $\zeta^s$.
This choice fixes the convention required in (P-F56v).

The smooth quotient tangent sequence in P-F24, dualized to
coordinate functions, shows that the moving conormal character is
\[
 \operatorname{Sym}^3 V-V\otimes V^*.
 \tag{P-F56w}
\]
Indeed the removed orbit cotangent character is the adjoint
character $V\otimes V^*-\mathbf1$, and the remaining fixed
grading direction contributes the other $\mathbf1$.
This is an identity in the representation ring of the finite
cyclic group, with its left side realized by the actual
conormal representation.

Let $h_j$ and $e_j$ be the complete and elementary symmetric
characters in the $\alpha_i$.
The identity
$\prod_i(1+t\alpha_i)=
 \prod_i(1-t^2\alpha_i^2)/\prod_i(1-t\alpha_i)$
gives $h_3=e_3+h_1\sum_i\alpha_i^2$.
The multisets $\{\alpha_i^2\}$ and $\{\alpha_i^{-1}\}$
differ by replacing $\alpha_0^{-1}$ with one: this follows
from $\alpha_i^2=\alpha_{i+1}^{-1}$ for $i<n-1$
and $\alpha_{n-1}=1$.
Consequently (P-F56w) is
\[
                       e_3+h_1-\alpha_0^{-1}h_1.
 \tag{P-F56x}
\]
The trivial weight in $h_1$ cancels the term with $i=0$
in $\alpha_0^{-1}h_1$.
For every $2\le k<n$, the remaining subtracted weight
$w_k-w_0=w_k-1$ equals the weight of the triple
$\{0,1,k\}$, because $w_0+w_1=-1$.
The last subtracted weight is $w_1-w_0=-3$.
It equals the weight of $\{0,2,3\}$:
for $n\ge5$ the integer sum is $1+4-8=-3$;
for $n=4$ it is $1+4+0=5\equiv-3\pmod8$.
These removed triples are distinct.
The remaining character is exactly the positive multiset
(P-F56u). The count is
$(n-1)+\binom n3-(n-2)-1=D$.
Nonvanishing of its singleton weights is immediate, and
a distinct triple sum is nonzero by its unique smallest
two-adic valuation among the nonzero summands.
The residue formula (P-F31b) now gives (P-F56v).

For reference, put
$a_0=g(n+1)(n+2)/6\in\mathbb Z$.
The determinant of $\operatorname{Sym}^3 V$ is
$(\det V)^{\binom{n+2}2}$, whereas
$\det(V\otimes V^*)=1$.
It follows from (P-F56w) and $3\sum_i w_i\equiv1\pmod M$
that the sum of the exponents $gb$ in (P-F56v) is congruent
to $a_0$ modulo $M$.
Inverting each factor therefore gives the exact relations
\[
 \mathcal C_n(\zeta^{-1})
       =(-1)^D\zeta^{a_0}\mathcal C_n(\zeta),
 \qquad
 W_n(s)=(-1)^D W_n(-s-a_0).
 \tag{P-F56y}
\]
They fix a reciprocity symmetry, but do not prescribe the sign
of the wave at a specified nonnegative integer.

Finally, if $h=M/2$ and
$\mathcal C_n(\zeta)=\sum_{j=0}^{h-1}c_{n,j}\zeta^j$
is its unique representative in
$\mathbb Q[t]/(t^h+1)$, then for every integer $s$,
\[
             2W_n(s)=
       (-1)^{\lfloor s/h\rfloor}c_{n,\,s\bmod h}.
 \tag{P-F56z}
\]
Indeed the primitive dyadic trace of $\zeta^b$ is zero unless
$h\mid b$, and is $h(-1)^{b/h}$ otherwise. This follows
directly by summing the geometric series over the odd
exponents modulo $M$.
Only $j=s\bmod h$ survives in (P-F56c), giving (P-F56z).
In particular, at $0\le x<h$ the exact one-witness
obstruction is $f_n(x)<c_{n,x}$.
This is a coefficient reformulation of the unresolved signed
wave question, with no positivity of those coefficients assumed.

\paragraph{A contraction gate involving only quadratic-many factors.}
The unresolved mass comparison can be reduced further to a
dimension-increment inequality. This reduction proves the error
from the two selected sibling branches; it does not assume or
prove the remaining uniform contraction inequality.

For $n\ge4$ write $M_n=2^{n-1}$, let
$b_n=3^{-1}\pmod {M_n}$, and use the positive multiset
$\mathcal S_n$ of (P-F56u). For every odd residue $a\bmod M_n$
define
\[
 P_n(a)=
 \prod_{b\in\mathcal S_n}
       |1-\exp(2\pi iab/M_n)|^{-1},\qquad
 Q_n=\sum_{\substack{a\bmod M_n,\ a\ {\rm odd}\\
                               a\ne b_n,-b_n}}
                              \frac{P_n(a)}{P_n(b_n)}.
 \tag{P-F56aa}
\]
Every factor is nonzero. Multiplication of the primitive-root
index by the odd number $g=\gcd(n,3)$ permutes the odd residues.
Thus (P-F56v) implies exactly $Q_n=E_n/|Z_n|$; the common
normalization $1/M_n$ cancels.

For $n\ge5$, form the multiset of residues modulo $M_n$
\[
 \begin{split}
 \mathcal T_n={}&\{1\}\\
 &\sqcup\{1+w_j+w_k:2\le j<k<n,\ (j,k)\ne(2,3)\}\\
 &\sqcup\{2+w_k:3\le k<n\}\sqcup\{6\},
 \qquad |\mathcal T_n|=\binom{n-1}2,
 \end{split}
 \tag{P-F56ab}
\]
where $w_i=(-2)^i$ for $i<n-1$ and $w_{n-1}=0$,
as before. Define the positive increment product
$J_n(a)=\prod_{b\in\mathcal T_n}
 |1-\exp(2\pi iab/M_n)|^{-1}$.
The exact recursion is
\[
             P_n(a)=J_n(a)P_{n-1}(-a\bmod M_{n-1}).
 \tag{P-F56ac}
\]

Here is a sufficient uniform contraction hypothesis, involving
only the $\binom{n-1}2=O(n^2)$ factors of each increment:
there exist a fixed $0<\rho<1$ and an integer $n_0\ge5$ such that,
for every $n\ge n_0$ and every odd $a\bmod M_n$ with
$a\bmod M_{n-1}\notin\{b_{n-1},-b_{n-1}\}$,
\[
 J_n(a)+J_n(a+M_{n-1})\le \rho J_n(b_n).
 \tag{P-F56ad}
\]
PNEW-002--005 below prove this inequality with explicit constants.
We first show that it implies $Q_n\longrightarrow0$. The phase-window
argument at the end of this section then produces the positive
wave values used in PNEW-005.
In this implication there is no unproved finite base
certificate: any finite starting value of $Q_n$ suffices.

The selected sibling error is bounded unconditionally:
\[
 \delta_n:=
 \frac{P_n(b_n+M_{n-1})}{P_n(b_n)}
 \le
 2^{-n^3/6+4n^2-107n/6+19},
 \qquad
 \delta_n\le2^{-n^3/12}\quad(n\ge48).
 \tag{P-F56ae}
\]
The exact implication of (P-F56ad) is
\[
                         Q_n\le\rho Q_{n-1}+2\delta_n.
 \tag{P-F56af}
\]

\emph{Proof.}
Remove the first variable from the chain weights.
The remaining weights $w_1,\ldots,w_{n-1}$ are $-2$
times the weights of the $(n-1)$-variable chain, in the
same order. In the multiset (P-F56u), all triples with
smallest index at least one occur. The shifted smaller
multiset removes from these exactly the triples
$\{1,2,k\}$ for $3\le k<n$ and $\{1,3,4\}$.
The additional triples with smallest index zero are
$\{0,j,k\}$ for $2\le j<k<n$, except $\{0,2,3\}$.
The additional singleton is $w_0=1$.
These extra weights are (P-F56ab): the first removed
family has weight $w_1+w_2+w_k=2+w_k$, and the last
triple has weight six modulo $M_n$.
For $n\ge6$ its integer weight is $-2-8+16=6$;
for $n=5$ it is $-2-8+0=-10\equiv6\pmod {16}$.
Thus the multiset decomposition is exact, including
multiplicities and the smallest allowed dimension.
Its cardinality is
$1+(\binom{n-2}2-1)+(n-3)+1=\binom{n-1}2$.
All its factors are nonzero because it is a submultiset
of $\mathcal S_n$. Substitution of the root
$\exp(-4\pi ia/M_n)$ in the smaller product proves
(P-F56ac).

Put $M=M_n$ and $\eta=\exp(2\pi ib_n/M)$.
The two roots indexed by $b_n$ and $b_n+M/2$
have the same smaller-dimensional factor in (P-F56ac).
In the increment, every even weight gives the same
factor at $\eta$ and $-\eta$.
The odd weights of $\mathcal T_n$ are exactly
\[
 \mathcal O_n=
 \{1\}\sqcup
 \{1+w_j+w_k:2\le j<k<n,\ (j,k)\ne(2,3)\}.
 \]
It follows that
\[
                 \delta_n
       =\prod_{b\in\mathcal O_n}
                     \left|\tan\frac{\pi b_n b}{M}\right|.
 \tag{P-F56ag}
\]
No tangent is singular, since $b_n$ and every $b$ here
are odd and $M\ge16$.

Among these odd weights, the bulk pairs
$2\le j<k\le n-2$, excluding $(2,3)$,
give
$B=\binom{n-3}2-1$ weights divisible by three.
For such a weight $b=1+w_j+w_k$, one has
$|b|<2^{k+1}$ and $|b|\le M$.
Since $3b_n\equiv1\pmod M$,
$\tan(\pi b_n b/M)=\tan(\pi b/(3M))$.
On $[-\pi/3,\pi/3]$, the bounds
$|\sin u|\le|u|$ and $\cos u\ge1/2$ give
\[
 \left|\tan\frac{\pi b_n b}{M}\right|
       \le\frac{2\pi |b|}{3M}
       \le\frac{8\,2^k}{M}.
 \tag{P-F56ah}
\]
The sum of $n-1-k$ over these bulk pairs is
$\binom{n-2}3-(n-4)$: before the one excluded pair,
it counts triples $2\le j<k<\ell\le n-1$.
There are $n-2$ exceptional odd factors, namely the
singleton one and the pairs using the zero-weight
variable. Each is at most $M/2$.
Indeed an odd residue modulo $M$ is at distance at least
one from $M/2$, so its cosine in (P-F56ag) has absolute
value at least $\sin(\pi/M)\ge2/M$, while its sine has
absolute value at most one.
Combining these estimates gives
\[
 \begin{split}
 \log_2\delta_n
 &\le 3\left(\binom{n-3}2-1\right)
        -\binom{n-2}3+(n-4)+(n-2)^2\\
 &=-n^3/6+4n^2-107n/6+19 .
 \end{split}
 \tag{P-F56ai}
\]
For $n\ge48$, $4n^2\le n^3/12$ and
$-107n/6+19<0$, proving the last assertion of
(P-F56ae).

Finally group the odd residues modulo $M_n$ into their
two lifts above an odd residue modulo $M_{n-1}$.
Reduction of $b_n$ is $b_{n-1}$.
The two selected parent classes therefore have the two
selected children $\pm b_n$ and the two unselected
children $\pm b_n+M_{n-1}$.
Complex conjugation makes the latter two contributions
equal, with total normalized mass $2\delta_n$.
For every other parent, (P-F56ac) and (P-F56ad)
bound its two children by $\rho$ times its normalized
mass in $Q_{n-1}$.
Here
$P_n(b_n)=J_n(b_n)P_{n-1}(b_{n-1})$,
using conjugation for the minus sign in (P-F56ac).
This proves (P-F56af), with no factor of two lost
in the nonselected branches.

Since $\delta_n\to0$ and $\rho<1$, the recursion forces
$Q_n\to0$. Explicitly, for any $\epsilon>0$, choose
$N$ so that $2\delta_n\le\epsilon(1-\rho)/2$
for all $n>N$. Iterating gives
\[
 Q_n\le \rho^{n-N}Q_N+\epsilon/2
                  \quad(n\ge N),
 \tag{P-F56aj}
\]
and the first term tends to zero.
The starting value $Q_N$ is finite by its definition.
This completes the conditional contraction reduction
and the unconditional selected-sibling bound.



\paragraph{The bounded phase window, with its normalization.}
Put $r=3g\in\{3,9\}$, choose $a_*r\equiv1\pmod M$, and set
$Z_n=A_n(\zeta^{a_*})$ and
$E_n=\sum_{a\ {
m odd},\,a\ne\pm a_*}|A_n(\zeta^a)|$.
Then $Q_n=E_n/|Z_n|$ by (P-F56v) and (P-F56aa).
Write $ra_*=1+kM$.  Here $1\le k<r$ and $\gcd(k,r)=1$.
For any integer $X$ the $r$ numbers
$Z_n\zeta^{-a_*X}\exp(-2\pi i k j/r)$, $0\le j<r$,
are equally spaced on the circle of radius $|Z_n|$.
One has real part at least $\cos(\pi/r)|Z_n|\ge |Z_n|/2$.
Replacing that phase by $\zeta^{-a_*j}$ changes its real part
by at most $2\pi|Z_n|/M\le |Z_n|/4$ when $n\ge6$.
The selected conjugate pair therefore contributes at least
$|Z_n|/2$ at some $x\in[X,X+r-1]$.
If $Q_n\le1/4$, the remaining terms have absolute sum at most
$|Z_n|/4$, and hence $W_n(x)\ge |Z_n|/4$.
PNEW-004 below gives the required selected-amplitude lower bound
directly from the exact positive residue product.



% Source: branches/plethysm/quantitative_refinement.tex
\section{An explicit cubic witness and a quantitative refinement}
\label{sec:plethysm-quantitative}

\paragraph{Status and dependence.}
The cubic fixed-locus and residue statements below are inherited from
ALPOGE\_GCT, P-F19, P-F23--P-F31, and P-F56.  The finite four-variable
calibration is ALPOGE\_GCT, P-F35.  We do not claim those results as new.
The revised dyadic estimates in PNEW-002--PNEW-005 are new quantitative
consequences of the same residue product; their proofs are given here.
All indices use the constituentwise conventions of GCT VI,
Definitions 1.2.2 and 1.2.4, printed pp.\ 13--14.

For $n\ge4$, put $g=\gcd(n,3)$, $d_0=n/g$, $M_n=2^{n-1}$,
$m_n=\binom{n+2}{3}$, and
\[
 f_n(s)=\dim\operatorname{Hom}_{GL_n}
  ((\det)^{3s/g},\operatorname{Sym}^{d_0s}(\operatorname{Sym}^3\mathbb C^n)).
\]
Set $w_i=(-2)^i$ for $0\le i<n-1$, and $w_{n-1}=0$.
Let $\mathcal S_n$ be the multiset consisting of the $n-1$ weights
$w_0,\ldots,w_{n-2}$ and all distinct-index triple sums
$w_i+w_j+w_k$, except the triples $\{0,1,k\}$ ($2\le k<n$)
and $\{0,2,3\}$.  Thus $|\mathcal S_n|=\binom n3$.
For a primitive $M_n$th root $\zeta$, the inherited exact residue is
\[
 C_n(\zeta)=\lim_{t\to\zeta}(1-t^{M_n})H_n(t)
       =\prod_{b\in\mathcal S_n}(1-\zeta^{gb})^{-1}.
 \tag{PNEW-R}
\]
Every factor is nonzero.  Define
\[
 P_n(a)=\prod_{b\in\mathcal S_n}|1-e^{2\pi iab/M_n}|^{-1},
 \qquad b_n=3^{-1}\pmod {M_n},
 \qquad
 Q_n=\sum_{\substack{a\bmod M_n\text{ odd}\\a\ne\pm b_n}}
            \frac{P_n(a)}{P_n(b_n)}.
\]
This is exactly $E_n/|Z_n|$ in P-F56: multiplication by the odd
number $g$ permutes primitive-root indices.  For $n\ge5$, the
increment $I_n(a)$ is the reciprocal product on
\[
 \mathcal T_n=\{1\}\sqcup
 \{1+w_j+w_k:2\le j<k<n,\ (j,k)\ne(2,3)\}
 \sqcup\{2+w_k:3\le k<n\}\sqcup\{6\}.
\]
Its cardinality is $\binom{n-1}{2}$, and cancellation of the shifted
smaller multiset gives
\[
 P_n(a)=I_n(a)P_{n-1}(-a\bmod M_{n-1}).
 \tag{PNEW-rec}
\]

\begin{proposition}[PNEW-001: a smallest-rank explicit witness]
For the cubic family one has
\[
                  2W_4(19)=\frac18>0=f_4(19).
\]
Consequently $p(f_4)>19$ and $s(f_4)\ge19$.  In fact the inherited
exact calibration gives $p(f_4)=s(f_4)=23$.  Rank four is the least
rank at which the primitive-$2^{n-1}$ wave gives an obstruction
$2W_n(x)>f_n(x)$ with $x\ge0$.
\end{proposition}
\begin{proof}
Salmon's cubic-surface invariant-ring theorem gives
\[
 H_4(t)=\frac{1+t^{25}}{(1-t^2)(1-t^4)(1-t^6)(1-t^8)(1-t^{10})}.
\]
See Huybrechts, \emph{The geometry of cubic hypersurfaces}, Chapter 4,
Theorem 4.1, p.\ 196, and the proof of Corollary 4.2; the complete
constituent calculation is ALPOGE\_GCT P-F35.  In
$\mathbb Q[z]/(z^4+1)$ the residue at a primitive eighth root is
\[
 C_4(z)=\frac{1+z^{25}}{(1-z^2)(1-z^4)(1-z^6)(1-z^{10})}
       =\frac{1+z+z^2+z^3}{8}.
\]
The primitive-dyadic trace formula, P-F56z, says that
$2W_4(x)=1/8$ for $x\equiv0,1,2,3\pmod8$, and $-1/8$ otherwise.
The numerator of $H_4$ has no odd term below degree 25, so
$f_4(19)=0$.  The strict-index conclusion is P-F56e'.

For $n=3$, the ternary-cubic invariant ring has generators of
ordinary degrees 4 and 6, so $H_3(t)=1/((1-t^4)(1-t^6))$.
An exact primary locator is McCrory--Shifrin--Varley (1994),
Section 3.2, printed p.~16.
At the primitive fourth roots $C_3=1/2$, hence
$2W_3(x)=\tfrac12\cos(\pi x/2)$.  Its positive values occur only
when $4\mid x$, where $f_3(x)\ge1$.  At $n=2$ one has
$H_2(t)=1/(1-t^2)$.  An elementary character check is as follows.
Put $d(j)=\#\{(a,b)\in\mathbb Z_{\ge0}^2:2a+3b=j\}$, with
$d(j)=0$ for $j<0$.  The difference of the weight-zero and
weight-two multiplicities in $\operatorname{Sym}^{2s}
(\operatorname{Sym}^3\mathbb C^2)$ is, from its Gaussian
polynomial character,
\[
 f_2(s)=d(3s)-d(s-1)-d(s-2)-d(s-3).
\]
The first term is $\lfloor s/2\rfloor+1$.  The sum of the last
three is $\lceil s/2\rceil$, since its generating function is
$(t+t^2+t^3)/((1-t^2)(1-t^3))=t/((1-t)(1-t^2))$.
Thus $f_2(s)$ equals one for even $s$ and zero for odd $s$, and
$2W_2(x)=(-1)^x\le f_2(x)$.  Rank one has no nontrivial primitive
dyadic subgroup.  These classical smaller-rank descriptions show
minimality; no asymptotic claim is inferred from the finite witness.
\end{proof}

\begin{lemma}[PNEW-002: a sharper near-increment estimate]
Let $n\ge5$, let $3\le j<M_n$ be odd, put
$a_j\equiv j3^{-1}\pmod{M_n}$, $t=\log_2j$, and
$R=\lceil\log_2(16j)\rceil$.  With
$B=(n-1)(n-4)/2$ one has
\[
 \left|\log_2\frac{I_n(a_j)}{I_n(b_n)}
        +Bt-\mathbf1_{3\mid j}\binom n2\right|
              \le \frac{16}{3}nR^2.
 \tag{PNEW-near}
\]
\end{lemma}
\begin{proof}
Write $l=n-1$.  The $B$ weights divisible by three are the active
weights $1+w_u+w_v$, $2\le u<v<l$, excluding $(2,3)$; the weights
$2+w_k$ for $3\le k<l$; and 6.  The remaining $l$ weights are
$1,2$, and $1+w_k$ for $2\le k<l$.  Give a weight its largest
active dyadic exponent $k$ and set $q=l-k$.  The exceptional
weights have one occurrence of each $q=1,\ldots,l$.
At each $q$ there are at most $n$ divisible weights: for a largest
index $k$ there are at most $k-2$ bulk pairs, one additional weight
$2+w_k$, and the separately assigned weight 6 at $k=2$.

Here are the elementary bounds needed for the comparison.  A sum
of at most $r\le3$ distinct signed dyadic terms has phase
$\kappa/3+(j/3)\sum_{i=1}^r\epsilon_i2^{-q_i}$, with
$q_1<\cdots<q_r$ and $\epsilon_i\in\{\pm1\}$.  Stop the initial
block before the first gap exceeding $R$.  Its numerator over
$3\cdot2^{q_u}$ is odd, so its distance from an integer is at least
$1/(3\cdot2^{q_u})$.  The discarded tail has magnitude at most
$1/(48\cdot2^{q_u})$, because $2^R\ge16j$.  Therefore the full
distance is at least $1/(6\cdot2^{q_1+(r-1)R})$.
Using $|1-e^{2\pi iy}|\ge4\|y\|$, a factor with $q_1\le R$ is
at least $2^{-3R-1}$; for two terms it is at least $2^{-2R-1}$.

Call $q>R$ low.  On a low divisible weight the cube-root shift
vanishes, and the target phase is $j$ times the reference phase.
Its magnitude is at most $(2j/3)2^{-q}<1/24$.
The inequalities $1-y^2/6\le\sin(y)/y\le1$ give error at most
$16j^2 4^{-q}$ from the main logarithm $t$.  Summing over at most
$n$ weights at each $q$ gives total error at most $n/48$.
For a high divisible weight the reference factor is at least
$2^{-q}/3$.  The preceding target bound and the upper bound 2
show that the absolute logarithm of the factor ratio is at most
$3R+2$.  Restoring its main term costs at most $t\le R$.
There are at most $nR$ high factors.  Thus the divisible-weight
error is bounded by
\[
                        4nR^2+2nR+n/48.
\]

Every exceptional reference factor lies in $[1/3,2]$.  If
$3\nmid j$, a low target factor has a nontrivial cube-root shift
perturbed by less than $1/24$, so its logarithmic ratio has
absolute value at most 3.  There are at most
$H=\min(R,l)$ high exceptional factors, each with logarithmic
ratio of absolute value at most $3R$.  Their total error is at
most $3n+3nR$.

If $3\mid j$, a low exceptional factor has logarithmic
reciprocal ratio $q-t$ with error at most 4.  Restoring the
missing high $q$ terms costs at most $H(H+1)/2\le nR$;
the high logarithmic ratios cost at most $3nR$; discarding
the low $-t$ terms costs at most $nR$.  Hence the error from
the exceptional main term $\sum_{q=1}^lq=\binom n2$ is at
most $5nR+4n$.

Since $j\ge3$, $R\ge6$.  In the second, larger case the total
error divided by $nR^2$ is at most
\[
                4+\frac7R+\frac{193}{48R^2}
       \le4+\frac76+\frac{193}{1728}<\frac{16}{3}.
\]
The first case is smaller.  This proves the lemma with all
constants uniform in $n$ and $j$.
\end{proof}

\begin{lemma}[PNEW-003: a sharper far-increment estimate]
For odd $a\bmod M_n$, put $j\equiv3a\pmod{M_n}$,
$1\le j<M_n$, and $J=\min(j,M_n-j)$.  Then
\[
 \log_2\frac{I_n(a)}{I_n(b_n)}
       \le \frac{25}{24}n^2+3n-\frac{(\log_2J-8)_+^3}{24}.
 \tag{PNEW-far}
\]
\end{lemma}
\begin{proof}
Set $l=n-1$, $\theta=a/2^l$, and
$y_{uv}=\theta(1+w_u+w_v)$ for $2\le u<v<l$.
Distinct subsets of the active dyadic weights have distinct sums
modulo $2^l$: their difference has a unique smallest dyadic
valuation.  This remains true after multiplication by odd $a$.
For $0\le h<l$, put $r=l-h$ and count the pairs with
$\|y_{uv}\|<2^{-h}$.  There are $e=(r-2)_+$ early indices.
Early--early pairs number at most $\binom e2$.  For each early
index, good late partners lie in one residue class modulo $2^r$;
the open interval $(-2^r,2^r)$ contains at most two representatives
of that class.  Good late--late pairs similarly number at most
two.  Thus the layer count is at most
$\binom e2+2e+2$, which is $\binom r2+1$ for $r\ge2$ and 2
for $r=1$.  Its sum over all layers is $\binom n3+n$.

If $J>3\cdot2^r$, pair the early indices consecutively into
$\lfloor e/2\rfloor$ disjoint pairs.  For two such pairs,
\[
 y_{uv}+y_{u+1,v}+y_{u,v+1}=3\theta\pmod1.
\]
At least one of the four associated cross factors is outside the
current layer.  Distinct blocks use disjoint factors, and the
omitted factor $(2,3)$ is internal to an index-pair.  The layer
count therefore loses at least $\binom{\lfloor e/2\rfloor}{2}$.
Summing for $r\le(\lfloor\log_2J\rfloor-2)_+$ gives loss at
least $(\log_2J-8)_+^3/24$, by summing
$\binom{\lfloor(r-2)/2\rfloor}{2}\ge(r-5)^2/8$ for $r\ge6$.

The remaining increment factors are exactly the two families
\[
       \{c\}\sqcup\{c+w_k:2\le k<l\},\qquad c=1,2.
\]
In each layer there are at most $(r-2)_+$ early terms and at
most two late terms, including the singleton $c$ in the late
class.  Each family's total layer count is at most
$\binom n2+1$.

For every nonzero dyadic phase $y$,
\[
 \log_2|1-e^{2\pi iy}|^{-1}
   \le -\log_2\|y\|-2
   \le \sum_{h=0}^{l-1}\mathbf1_{\{\|y\|<2^{-h}\}}-2.
\]
Add the omitted bulk factor only to the nonnegative layer counts,
and subtract 2 for each of the actual $\binom{n-1}{2}$ increment
factors.  The resulting upper bound is
\[
 \log_2 I_n(a)\le\binom n3+3n
                    -\frac{(\log_2J-8)_+^3}{24}.
\]
For the selected increment, all $B=(n-1)(n-4)/2$ weights
divisible by three can be estimated at a small phase.  Write
$\alpha=13/12$.  The classical inequality $\pi<22/7$ and
the exact integer inequality $(44/21)^{12}<2^{13}$ give
$2\pi/3<2^\alpha$.  For an active bulk weight with largest
index $k$ and other index $j$, its selected factor is at most
\[
  2^\alpha\frac{2^k}{M_n}
                       (1+2^{j-k}+2^{-k}).
\]
For $2+w_k$, $3\le k<n-1$, replace the parenthesis by
$1+2^{1-k}$; for the remaining weight 6, assign $k=2$ and
use the parenthesis $3/2$.  The inequality
$\log_2(1+u)\le2u$ for $u\ge0$ follows from
$\ln(1+u)\le u$ and $\ln2>1/2$.
For the bulk corrections,
$\sum_{j=2}^{k-1}2^{j-k}<1$ and
$\sum_{k=3}^\infty(k-2)2^{-k}=1/2$ give a total at most
$2(n-5)+1$; omitting the pair $(2,3)$ only improves it.
The $2+w_k$ corrections sum to at most 1, and the correction
for 6 is less than 1.  Their combined sum is therefore at
most $2n$ for every $n\ge5$.

The sum of $n-1-k$ over all divisible weights is exactly
\[
 \binom{n-2}{3}-(n-4)+\binom{n-3}{2}+(n-3)
                  =\binom{n-1}{3}-n+4.
\]
Each of the $n-1$ exceptional factors is at most 2.  Hence
\[
 \log_2 I_n(b_n)\ge
       \binom{n-1}{3}-\frac{13}{12}B-4n+5.
\]
Subtracting this from the preceding layer bound leaves
$25n^2/24+67n/24-11/6$, which is at most
$25n^2/24+3n$.  This proves the lemma.
\end{proof}

\begin{lemma}[PNEW-004: a cubic bound for the starting mass]
For every $n\ge5$,
\[
                  Q_n\le2^{\,2n^3/3+10n/3-4}.
 \tag{PNEW-mass}
\]
\end{lemma}
\begin{proof}
Use layers as in the preceding proof, now with active indices
$0,\ldots,n-2$.  For distinct active triple sums, the layer
count is at most $\binom r3+2\binom r2+2r+2$, by fixing the
early subset and using at most two late representatives in each
residue class.  For distinct active pair sums it is at most
$\binom r2+2r+2$.  For singleton weights it is at most $r$:
a nonzero late multiple of $2^r$ cannot lie in the open interval
$(-2^r,2^r)$.  Add omitted weights only to the nonnegative layer
counts.  Subtracting 2 for the actual $\binom n3$ factors gives
\[
 \log_2P_n(a)\le\binom n4+\binom n3+5\binom n2+4(n-1).
\]
Here is a direct selected-amplitude estimate from the exact product,
avoiding any further determinant formula.  Among all active distinct
triples $0\le i<j<k\le n-2$, each weight is divisible by three and
has absolute value less than $2^{k+1}$.  Its selected factor is at
most $8\cdot2^k/M_n$, using $\pi<4$.  Summing $n-1-k$
over all such triples gives $\binom n4$.  To obtain the actual
bulk in $\mathcal S_n$, delete $\{0,1,k\}$ for $2\le k\le n-2$
and $\{0,2,3\}$.  The sum of their contributions $n-1-k-3$
is $(n-3)(n-8)/2+(n-7)$.  Bound each remaining nonbulk factor
by 2.  There are $\binom n3$ actual factors altogether.  Thus
\[
 \begin{split}
 \log_2P_n(b_n)&\ge
 \binom n4-2\binom{n-1}{3}-\binom n3
                                  -\frac{n^2}{2}+\frac{7n}{2}-3\\
 &=\frac{n^4}{24}-\frac{3n^3}{4}
                        +\frac{59n^2}{24}-\frac{3n}{4}-1.
 \end{split}
 \tag{PNEW-amplitude}
\]
There are fewer than $2^{n-1}$ summands in $Q_n$.  Combining
these bounds yields
\[
                  \log_2Q_n\le\frac23n^3+\frac{10}{3}n-4.
\]
If $Q_n=0$, the asserted inequality is immediate.
\end{proof}

\begin{theorem}[PNEW-005: an explicit uniform threshold]
For every integer $n\ge79000$,
\[
                  p(f_n)>2^{n/8},\qquad s(f_n)\ge2^{n/8}.
\]
Thus the inherited plethysm refutations of PH2 and SH(a) have an
explicit all-rank threshold of $79000$ by elementary estimates.
This is a sufficient threshold, not a claim of optimality.
\end{theorem}
\begin{proof}
Put $N=39500$ and $c=1168/12500$.  We first show that for every $n\ge N$ and
every nonselected odd $a\bmod M_n$,
\[
                   I_n(a)/I_n(b_n)\le2^{-2n^2/7}.
\]
By conjugation use the least positive absolute residue
$J$ of $3a$ and set $t=\log_2J$.  If $J\ne1$, then
$t>19/12$, since $3^{12}>2^{19}$.  For $t\le cn$, PNEW-002 and $R\le t+5$ give
the upper bound
\[
 F_n(t)=-\frac{(t-1)n^2}{2}
          +\frac52nt+\frac{16}{3}n(t+5)^2.
\]
This convex quadratic takes its maximum on
$[19/12,cn]$ at an endpoint.  The first endpoint gives
\[
 F_n(19/12)=-\frac{7n^2}{24}+\frac{50783}{216}n.
\]
After division by $n^2$, both endpoint upper bounds decrease
with $n$.  At $N$ their margins below $-2/7$ are exactly
\[
 -\frac27-\frac{F_N(19/12)}{N^2}
       =\frac{19}{59724000}>0,\qquad
 -\frac27-\frac{F_N(cN)}{N^2}
       =\frac{8303507}{86406250}>0.
\]
For $t\ge cn$, PNEW-003 and the exact rational inequality
\[
 \frac{(cN-8)^3}{24N^2}-\frac{25}{24}-\frac3N-\frac27
       =\frac{1676026865671}{255978515625000}>0
\]
give the same bound.  The left side is increasing with $N$
replaced by any real $n\ge N$: the derivative of
$(cn-8)^3/n^2$ is $(cn-8)^2(cn+16)/n^3>0$,
and $-3/n$ is increasing.

Grouping the two lifts of every nonselected parent in
PNEW-rec gives the exact recurrence
\[
 Q_n\le2^{1-2n^2/7}Q_{n-1}+2\delta_n,\qquad
 \delta_n\le2^{-n^3/12}\quad(n\ge48),
\]
where the last bound is the selected-sibling product P-F56ae.
In particular every contraction factor is at most $1/2$ for
$n\ge N$.  PNEW-004 and iteration therefore give
\[
 Q_n\le
 2^{(2/3)N^3+(10/3)N+(n-N+1)-2(n^3-(N-1)^3)/21}
                  +4\cdot2^{-N^3/12}.
\]
We used
$\sum_{k=N}^n k^2\ge\int_{N-1}^{n}x^2\,dx$.
At $n=2N$ the first exponent is exactly
$-445603261$, hence less than $-3$; it decreases for
all larger $n$.  The second summand is less than $1/8$.
Thus $Q_n<1/4$ for every $n\ge2N=79000$.

For completeness, the conversion to an index bound uses the
following amplitude and phase facts.  If $r=3g\in\{3,9\}$ and
$a_*r\equiv1\pmod{M_n}$, then
\[
 |Z_n|\ge2^{n^4/24-2n^3}.
\]
Indeed $|Z_n|=P_n(b_n)/M_n$, and PNEW-amplitude is stronger.
Among any $r$ consecutive integer arguments, the selected
conjugate pair has a value at least $|Z_n|/2$: its phases differ
from $r$ equally spaced phases by at most $2\pi/M_n$, and
$\cos(\pi/r)\ge1/2$.  For $n\ge6$ the error is at most
$1/4$.  Subtracting the remaining mass $Q_n|Z_n|\le|Z_n|/4$
gives an $x$ in
$[\lceil2^{n/8}\rceil,\lceil2^{n/8}\rceil+8]$ with
\[
                    2W_n(x)\ge2^{n^4/24-2n^3-1}.
\]
On the other hand
$f_n(x)\le\binom{m_n+d_0x-1}{m_n-1}$.
For $n\ge1024$, use $m_n\le n^3/5$,
$\log_2n\le n/16$, $2^{n/8}\ge n^2$, and
$x\le2^{n/8+1}$ to obtain
\[
 \log_2\max(1,f_n(x))
 \le\frac{n^3}{5}(n/8+\log_2n+2)
 \le\frac{3n^4}{80}+\frac{2n^3}{5}
                  <\frac{n^4}{24}-2n^3-1.
\]
The last difference is $n^4/240-12n^3/5-1>0$.
Hence $2W_n(x)>f_n(x)$.  The exact paired-constituent identity
$q_a(X)-q_{a+L/2}(X)=2W_n(a)$ implies that any positivity shift
at most $x$, or saturation shift below $x$, is impossible.
This proves the theorem for the displayed sufficient threshold.
\end{proof}


% Source: branches/plethysm/exterior_witnesses.tex
\section{Exterior-power bounds and a positive initial exception}
\label{sec:exterior-witnesses}

\begin{proposition}[PNEW-006: an exterior ambient bound]
Put $g=\gcd(n,3)$, $k=3s/g$, and $d=ns/g$, for $s\ge1$.
Then
\[
 f_n(s)=\dim\operatorname{Hom}_{GL_k}
       ((\det)^n,\bigwedge^d(\bigwedge^3\mathbb C^k))
       \le \binom{\binom{k}{3}}{d}.
 \tag{PNEW-ext}
\]
The binomial coefficient is zero when its lower argument exceeds
its upper argument.  In particular the exact dyadic residue gives
\[
 2W_9(7)=\frac{782264335828}{64}
       >2319959400=\binom{35}{21}\ge f_9(7),
 \tag{PNEW-nine}
\]
so $p(f_9)>7$ and $s(f_9)\ge7$.
There is also a witness with a proved nonzero actual coefficient:
\[
 2W_9(6)=\frac{1819142906755}{64}>1=f_9(6).
 \tag{PNEW-nine-positive}
\]
At rank 15 the certificate gives the larger finite witness
\[
 2W_{15}(60)>2^{881}>2^{866}>\binom{979}{300}\ge f_{15}(60),
 \tag{PNEW-fifteen-sixty}
\]
and hence $p(f_{15})>60$, $s(f_{15})\ge60$.
These finite consequences combine the inherited exact residue with
the exterior identity already proved in ALPOGE\_GCT P-F32.
The contribution here is the explicit comparison and its certificate;
no novelty is claimed for the exterior identity or its top-power case.
\end{proposition}
\begin{proof}
Work first in the ring of symmetric functions over $\mathbb Q$.
Let $\omega$ be the involution $p_r\mapsto(-1)^{r-1}p_r$,
so $\omega(s_\lambda)=s_{\lambda'}$.  The power-sum formulas
\[
 h_3=\frac{p_1^3+3p_1p_2+2p_3}{6},\qquad
 e_3=\frac{p_1^3-3p_1p_2+2p_3}{6}
\]
show directly that
$\omega(p_r[h_3])=(-1)^{r-1}p_r[e_3]$.
Expanding $h_d=\sum_{\alpha\vdash d}p_\alpha/z_\alpha$
and $e_d=\sum_{\alpha\vdash d}(-1)^{d-\ell(\alpha)}
p_\alpha/z_\alpha$ therefore gives
\[
                         \omega(h_d[h_3])=e_d[e_3].
\]
The target partition $(k^n)$ transposes to $(n^k)$.
Specialization to $k$ variables preserves the coefficient of that
Schur function: terms with more than $k$ rows vanish, and the
remaining Schur polynomials are linearly independent.  The last
partition represents $\det^n$ of $GL_k$.  The right hand side is
the character of the exterior power in (PNEW-ext).  Complete
reducibility and its ambient dimension give the inequality.

For $(n,s)=(9,7)$, one has $(k,d)=(7,21)$ and
$\dim\bigwedge^3\mathbb C^7=35$.  The exact residue computation
described below produces the two displayed rational values.  Their
comparison with the integer $\binom{35}{21}=2319959400$ is exact.
For $(n,s)=(9,6)$, put $U=\bigwedge^3\mathbb C^6$.
Then $\dim U=20$ and $\det U=(\det\mathbb C^6)^{10}$,
because each variable occurs in $\binom52=10$ basis triples.
Exterior duality gives
\[
 \bigwedge^{18}U\simeq(\bigwedge^2 U)^*\otimes\det U.
\]
The alternating pairing $U\otimes U\to\bigwedge^6\mathbb C^6$
given by exterior multiplication is nondegenerate.  It is alternating
because three is odd.  The irreducibility of $\bigwedge^3\mathbb C^6$
and Schur's lemma show that this equivariant pairing is unique up
to a scalar.  Hence the multiplicity of $\det$ in $\bigwedge^2U$
is one, and the multiplicity of $\det^9$ in $\bigwedge^{18}U$
is one.  This proves $f_9(6)=1$.
For the rank-15 comparison at $s=60$, use the original symmetric
ambient space: $\dim\operatorname{Sym}^3\mathbb C^{15}=680$ and
$d=15s/3=300$, so complete reducibility gives
$f_{15}(60)\le\dim\operatorname{Sym}^{300}(\mathbb C^{680})
=\binom{979}{300}$.  The two strict power-of-two comparisons in
(PNEW-fifteen-sixty) are verified as integer inequalities in the
certificate below.
\end{proof}

\paragraph{The exact finite computation.}
The script \texttt{scripts/certify\_plethysm.py} uses the positive
multiset $\mathcal S_n$ and performs arithmetic in
$\mathbb Q[z]/(z^{M/2}+1)$ only.  Multiplication by $z^a$ is a
signed permutation of the monomial basis.  Each cycle of that
permutation has total sign $-1$, because $M$ is a power of two
and $a\not\equiv0\pmod M$.  Thus $1-z^a$ is invertible, and the
linear recurrence around each cycle solves
$(1-z^a)v_{\rm new}=2v_{\rm old}$ over the integers.
The script checks that identity after every factor, before
cancelling powers of two.  It independently compares the positive
multiset against the character
$\operatorname{Sym}^3V-V\otimes V^*$ with all multiplicities.
Consequently its computed coefficient is exactly the residue
coefficient in (P-F56v), including the phase.  Formula (P-F56z)
then gives $2W_n(x)$ with no numerical approximation.
The complete rational results and integer inequality margins are
in \texttt{certificates/plethysm\_exact.json}.  Eight witnesses and
all 663 individual inverse-factor identities pass; that count
refers to the four dimensions $4,9,10,15$, not to an extrapolation.

\begin{proposition}[PNEW-007: a nonzero input with strict saturation index at least seven]
Let
\[
 \mu=(3),\qquad\lambda=(35),\qquad\pi=(7^{15}),
 \qquad h(s)=a_{s\lambda,\mu}^{s\pi}=f_{15}(7s).
\]
Then
\[
                      h(1)=1,\qquad p(h)>7,\qquad s(h)\ge7.
 \tag{PNEW-positive-start}
\]
Thus this is an exception inside the nonzero-coefficient population
used in narrow SH(b), rather than merely an input whose initial
coefficient vanishes.
\end{proposition}
\begin{proof}
At $n=15,s=7$, (PNEW-ext) identifies the coefficient with the
multiplicity of $\det^{15}$ in
$\bigwedge^{35}(\bigwedge^3\mathbb C^7)$.  The inner space has
dimension 35, and its determinant is $\det^{\binom62}=\det^{15}$.
The top exterior power is exactly this one-dimensional representation,
so $h(1)=1$.

For $s=14$, the same bound gives
\[
             f_{15}(14)\le\binom{364}{70}<2^{253}.
\]
The exact residue certificate gives
$2W_{15}(14)=A/64$, where $2^{888}<A<2^{889}$.
For completeness, $A$ is the integer obtained by concatenating
the following four digit strings, without spaces or line breaks:
\begin{center}\footnotesize\ttfamily
3333954305645752920064485456028591317103939322607541526148896903832899860\\
4542313254224303373186420275533728219695881859566923177720354053196664566\\
8443320869445013925724828682556583506896582799729547390014645903631244200\\
6703694838860313433573900791230758229801554002547
\end{center}
The script verifies the full cyclotomic recurrence and the strict
integer inequalities, not logarithmic floating-point surrogates.
It follows that $2W_{15}(14)>f_{15}(14)$.

We must check the effect of the odd dilation on constituents, since
an even dilation could discard the primitive dyadic wave.  Let $L$
be a period of $f_{15}$, write $D=\gcd(L,7)$ and $L'=L/D$.
Then $L'$ is a period of $h$, and $v_2(L')=14$.
On residue $a$ modulo $L'$ its constituent is
$q_{7a\bmod L}(7X)$.  Increasing $a$ by $L'/2$ increases
the old residue by $(7/D)L/2$, an odd multiple of $L/2$.
Thus (P-F56d) gives exactly
\[
 q_{7a}(7X)-q_{7(a+L'/2)}(7X)=2W_{15}(7a).
\]
At $a=X=2$ the first value is $h(2)=f_{15}(14)$ and
the second is strictly negative.  No positivity shift at most 2
can make every constituent coefficient nonnegative, and no
saturation shift below 2 can make every nonzero constituent
positive at all positive integers.  This already proves
$p(h)>2$ and $s(h)\ge2$.
For the stronger bounds stated above, the same exact certificate gives
\[
       2W_{15}(49)>2^{875}>2^{766}>
               \binom{924}{245}\ge f_{15}(49).
\]
Here the last bound is the symmetric ambient dimension
$\dim\operatorname{Sym}^{245}(\mathbb C^{680})$.
Apply the identical paired-constituent argument at $a=X=7$.
All displayed power-of-two comparisons are exact integer tests
in the accompanying script.
\end{proof}


\section{Complete four-variable calibration}

% Source: branches/plethysm/sections/positivity_quaternary_indices.tex
% Reproduced inherited proof; finding P-F35.
\subsection{Exact indices for the four-variable cubic calibration}

\begin{proposition}[P-F35: both quaternary-cubic indices are twenty-three]
\label{prop:quaternaryindices}
For the ordinary cubic-invariant stretching function of P-F31 with
$n=4$,
\[
 f_4(s)=\dim\left(\operatorname{Sym}^{4s}
                  (\operatorname{Sym}^3\mathbb C^4)\right)^{SL_4},
\]
the strict positivity and saturation indices of GCT VI are exactly
\[
                          p(f_4)=s(f_4)=23.
 \tag{P-F35a}
\]
Its minimal quasipolynomial period is $120$. In particular the
primitive eighth-root poles of P-F31 occur in the classical
explicit series and coexist with these finite indices. This
single dimension gives no asymptotic index bound.
\end{proposition}

\begin{proof}
Salmon's invariant-ring theorem gives algebraically independent
invariants of ordinary degrees $8,16,24,32,40$, together with
one invariant of degree $100$ whose square lies in the polynomial
ring they generate. An exact inspected locator is Huybrechts,
\emph{The geometry of cubic hypersurfaces}, Chapter~4,
Theorem~4.1, printed p.~196; the following proof of Corollary~4.2
also states the resulting graded decomposition. The ordinary
degrees in the polynomial subring are divisible by eight, whereas
those in its degree-$100$ translate are congruent to four modulo
eight. Thus the sum of these two modules is direct; it spans
by the monic quadratic relation. Dividing ordinary degrees by
four gives
\[
 H_4(t)=\sum_{s\ge0}f_4(s)t^s
       =\frac{1+t^{25}}{(1-t^2)(1-t^4)(1-t^6)(1-t^8)(1-t^{10})}.
 \tag{P-F35b}
\]
This also directly checks the nonzero primitive eighth-root
pole: the factor $1-t^8$ vanishes simply there, the other four
factors do not vanish, and $1+\zeta^{25}=1+\zeta\ne0$.

Let $P(y)$ be the global quasipolynomial with generating series
\[
                  \sum_{j\ge0}P(j)z^j
                       =\frac1{\prod_{a=1}^5(1-z^a)}.
\]
We will use its constituent polynomials explicitly. Put
\[
 Q(y)=\frac{y^4}{2880}+\frac{y^3}{96}
            +\frac{31y^2}{288}+\frac{85y}{192}
            +\frac{50651}{86400}.
\]
For an integer residue $r$, define the periodic numbers
\[
 a_r=\begin{cases}1&r\equiv0,1\pmod4,\\-1&r\equiv2,3\pmod4,\end{cases}
 \quad
 b_r=\begin{cases}2&3\mid r,\\-1&3\nmid r,\end{cases}
 \quad
 c_r=\begin{cases}4&5\mid r,\\-1&5\nmid r.\end{cases}
\]
Then the period-$60$ constituent on $r$ is
\[
 P_r(y)=Q(y)+\frac{(-1)^r(2y+15)}{128}
                   +\frac{a_r}{16}+\frac{b_r}{27}+\frac{c_r}{25}.
 \tag{P-F35c}
\]
For completeness, this formula follows by coefficient extraction
from the following exact partial-fraction identity, which can also
be checked by multiplying through by $\prod_{a=1}^5(1-z^a)$:
\begin{align*}
 \frac1{\prod_{a=1}^5(1-z^a)}
 &=\frac1{120(1-z)^5}+\frac1{24(1-z)^4}
   +\frac{31}{288(1-z)^3}+\frac3{16(1-z)^2}
   +\frac{20831}{86400(1-z)}\\
 &\quad+\frac1{64(1+z)^2}+\frac{13}{128(1+z)}
   +\frac{z+1}{16(z^2+1)}
   +\frac{z+2}{27(z^2+z+1)}\\
 &\quad+\frac{z^3+2z^2+3z+4}
                  {25(z^4+z^3+z^2+z+1)}.
 \tag{P-F35d}
\end{align*}
The first line gives $Q$ after using the binomial-series formula.
The two $1+z$ terms give $(-1)^r(2y+15)/128$.
The last three periodic terms give respectively $a_r/16$,
$b_r/27$ and $c_r/25$: multiplication by $1-z$ reduces
the last two numerators to $2-z-z^2$ and
$4-z-z^2-z^3-z^4$. This proves (P-F35c) for every
nonnegative integer argument and hence identifies the whole
constituent polynomials.

Every polynomial $P_r(y-1)$ has nonnegative coefficients.
Indeed its coefficients of degrees four, three and two are
respectively
\[
              \frac1{2880},\qquad\frac{13}{1440},
                       \qquad\frac{113}{1440}.
\]
Its linear coefficient is
\[
             \frac{247}{960}+\frac{(-1)^r}{64}
                     \ \ge\ \frac{29}{120}>0.
\]
Its constant term is
\[
 \frac{20831}{86400}+\frac{13(-1)^r}{128}
                        +\frac{a_r}{16}+\frac{b_r}{27}
                        +\frac{c_r}{25}.
\]
The sum $13(-1)^r/128+a_r/16$ takes the four values
$21/128,-5/128,5/128,-21/128$ on residues modulo four.
Since $b_r\ge-1$ and $c_r\ge-1$, the constant term is at least
\[
 \frac{20831}{86400}-\frac{21}{128}-\frac1{27}-\frac1{25}=0.
 \tag{P-F35e}
\]
All coefficient inequalities have therefore been verified
for every constituent without an interpolation assumption.

Equation (P-F35b) says that the even constituent polynomials
of $f_4$ are of the form $P_r(s/2)$ and the odd ones are
of the form $P_r((s-25)/2)$. To justify the latter globally,
for each odd residue modulo $120$ the equality holds for
all sufficiently large nonnegative $s$; thus it is the
polynomial identity on that residue. The series is a global
quasipolynomial, either by the canonical-ring theorem already
used or directly because its numerator over $(1-t^{120})^5$
has degree $595<600$.

After shifting by $23$, the odd polynomials become
$P_r(s/2-1)$ and therefore have nonnegative coefficients
by (P-F35e) and its preceding calculations. The even
polynomials become $P_r((s+23)/2)$, obtained from the
coefficientwise nonnegative polynomial $P_r(y-1)$ by
substituting $y=(s+25)/2$. A polynomial with nonnegative
coefficients remains such after substitution of an affine
polynomial with nonnegative coefficients. Thus $p(f_4)\le23$.

Every constituent has positive degree-four coefficient
$1/46080$, so none is the zero polynomial. On the other
hand (P-F35b) gives the actual value $f_4(23)=0$:
the denominator contains only even degrees, and the odd
numerator translate first starts at $25$. For any shift
$u<23$, the corresponding nonzero constituent still
vanishes at the positive integer $23-u$. Hence strict
saturation fails for all such shifts, giving
$s(f_4)\ge23$. Coefficientwise positivity implies strict
saturation for every nonzero constituent on positive
integers, so
\[
                   23\le s(f_4)\le p(f_4)\le23.
\]
This proves (P-F35a).

Finally $120$ is a period by the denominator in (P-F35b).
Its actual primitive eighth-root, third-root and fifth-root
poles force every quasipolynomial period to be divisible
by $8,3,5$, respectively. At a primitive third or fifth root,
only the factors of the corresponding divisible degrees
can vanish and the numerator is nonzero ($1+t^{25}$ equals
$1+t$ at a third root and equals two at a fifth root).
Thus their least common multiple $120$ is the minimal period.
\end{proof}

The exploratory exact rational computation is preserved in
\nolinkurl{quaternary_index_calibration.py} and its output file.
The proof above supplies explicit constituent formulas and
uniform inequalities, independently of the finite search.
The exact source URL is
\url{https://www.math.uni-bonn.de/people/huybrech/Notes.pdf}.



% Source: branches/plethysm/erratum.tex
\section{A companion erratum to the printed positive form}
\label{sec:positive-form-erratum}

\paragraph{Inherited finding P-F12 / PL-Q22, rechecked.}
In GCT VI, Theorem 1.6.1(d) and equation (4.4), the numerator
index cannot literally be the degree $d$ of the stretching
quasipolynomial, with numerator degree bounded by $d$.
The Kahle--Micha\l ek ray already contradicts that reading.
The correct Cohen--Macaulay positive form has unrestricted
numerator degree; the denominator has $d+1$ factors when
multiplicities are counted.  This correction does not rescue
PH3, whose two clauses fail by the cubic theorem proved above.

For clarity, all data needed for the contradiction are
\[
 \phi(s)=\frac{s+r_s}{3},\qquad
 (r_0,\ldots,r_5)=(3,-1,1,0,2,-2),\qquad
 \sum_{s\ge0}\phi(s)t^s=\frac{1+t^4}{(1-t^2)(1-t^3)}.
 \tag{PNEW-KM}
\]
This is Kahle and Micha\l ek (2018), Theorem 5 and Remark 6,
printed p.~4, for $a_{(4s),(3)}^{(7s,5s)}$.
Here is also a direct coefficient verification.  Put
$d(j)=\#\{(a,b)\in\mathbb Z_{\ge0}^2:2a+3b=j\}$,
with $d(j)=0$ for $j<0$.  Weight counting in
$\operatorname{Sym}^{4s}(\operatorname{Sym}^3\mathbb C^2)$
gives the Gaussian polynomial
\[
 \prod_{i=1}^3\frac{1-q^{4s+i}}{1-q^i}.
\]
Indeed its coefficient counts partitions in the $3$ by $4s$
rectangle, equivalently the monomial weight multiplicities.
For an $\mathfrak{sl}_2$ module, the multiplicity of highest
weight $2s$ is the difference between the weight-$2s$ and
weight-$(2s+2)$ multiplicities, since each irreducible has
one copy of every weight from its highest weight down by twos.
Multiplication of the Gaussian polynomial by $1-q$ therefore
gives
\[
 \phi(s)=d(5s)-d(s-1)-d(s-2)-d(s-3)\quad(s\ge1).
\]
Terms with two numerator shifts have degree greater than $5s$.
The elementary formula
$d(6a+r)=a+1-\mathbf1_{r=1}$ for $a\ge0$, $0\le r<6$,
follows by choosing $b$ of the same parity as the argument.
Substitution proves (PNEW-KM) for $s\ge3$; $s=0,1,2$
are checked directly.  Summing
$\phi(s)=\lfloor s/3\rfloor+\mathbf1_{2\mid s}$ gives
the rational series in (PNEW-KM).

The degree is one.  A positive numerator of degree at most one
and a denominator with two factors would have the form
\[
                \frac{1+ct}{(1-t^a)(1-t^b)},\qquad c\ge0.
\]
Since $\phi(1)=0$, one must have $c=0$ and $a,b>1$.
Then $\phi(2)=\phi(3)=1$ forces $\{a,b\}=\{2,3\}$.
The resulting coefficient of $t^4$ is one, whereas
$\phi(4)=2$, a contradiction.  The unrestricted positive
numerator $1+t^4$ above is consistent with
Cohen--Macaulayness.

For later use, (PNEW-KM) also gives exactly
$p(\phi)=s(\phi)=2$.  Shifting by two makes every constituent
coefficient nonnegative.  Before that shift the residue-five
constituent $(X-2)/3$ is nonpositive at a positive integer:
its values at 1 and 2 are negative and zero, respectively.
This proves both lower bounds and the matching upper bounds.

\section{General-form index failures}

% Source: branches/plethysm/sections/positivity_torus.tex
% Reproduced inherited proof; finding P-F14,P-F15.
\subsection{Fixed-rank torus embeddings refute the unrestricted index hypotheses}

The next two findings concern the full subgroup restriction problem,
GCT VI, Problem 1.1.3, with the specifications of Section 3.4.
They do not concern its irreducible-inner plethysm subproblem,
Problem 1.1.2. The distinction is material: the broad problem permits
direct sums of polynomial torus characters.

\begin{proposition}[P-F14: fixed-rank modular indices are unbounded]
\label{prop:torusmodular}
For each integer $N\ge2$, let $H=\mathbb G_m$, $G=GL_2$ and
\[
 \rho_N(t)=\operatorname{diag}(t^{N+1},t),\qquad
 \lambda=(1),\qquad \pi=2.
\]
This is a faithful explicit polynomial embedding allowed by
GCT VI, Section 3.4.1. Its stretching multiplicity is
\[
 f_N(n)=\dim\operatorname{Hom}_{\mathbb G_m}
       (\mathbb C_{2n},\operatorname{Sym}^n(\mathbb C^2)|_{\rho_N})
       =\mathbf1_{\{N\mid n\}},\qquad n\ge0,
\]
and its generating function has minimum possible modular index $N$:
\[
 F_N(T)=\frac1{1-T^N}.
 \tag{P-F14}
\]
The input has $\operatorname{rank}(x)=5$ in the exact definition of
GCT VI, Section 3.4. Consequently Hypothesis 3.3.4 (PH3) is false for
the full class of explicit polynomial subgroup restrictions as stated.
Here $p(f_N)=s(f_N)=0$, so this example does not refute PH2 or SH(a).
\end{proposition}

\begin{proof}
The representation defining $\rho_N$ is the direct sum of the
polynomial characters of weights $N+1$ and $1$. Section 3.4.1,
``Basic explicit homomorphisms'', explicitly permits constructions
from irreducible polynomial representations by $\oplus$ and $\otimes$;
it specifies that polynomial representations of a torus are sums
of polynomial one-dimensional characters. Both exponents here are
positive. The map is faithful because its second diagonal entry is
$t$, and the inverse on its image is that entry. Its image is closed:
inside the diagonal torus it is given by $z_1=z_2^{N+1}$.

A monomial $x^k y^{n-k}$, $0\le k\le n$, in the symmetric power
has $H$-weight $(N+1)k+(n-k)=n+Nk$. It has weight $2n$ if and only
if $Nk=n$. There is exactly one such monomial when $N\mid n$ and
none otherwise; when it exists, $k=n/N$ belongs to $[0,n]$.
This proves the multiplicity and generating-function formulas.
Its natural graded multiplicity algebra is
$\mathbb C[xy^{N-1}]$, a polynomial ring with generator of degree $N$:
the weight equation forces every admissible monomial to be a unique
power of $xy^{N-1}$.

Let $\zeta$ be a primitive $N$th root of unity. The rational function
$F_N$ has a pole at $\zeta$. In any expression
\[
 F_N(T)=\frac{B(T)}{\prod_j(1-T^{a_j})},
 \qquad B\in\mathbb C[T],\quad a_j>0,
\]
the denominator must vanish at $\zeta$; otherwise the right side is
regular there. Thus $\zeta^{a_j}=1$ for some $j$, so $N\mid a_j$
and $\max_j a_j\ge N$. This argument does not need positivity of
$B$ or a bound on the number of denominator factors. The displayed
form (P-F14) attains the lower bound, proving minimality. Its
constituent polynomials are the constants zero and one, so both
indices are zero by GCT VI, Definitions 1.2.2 and 1.2.4.

Section 3.4 defines
$\operatorname{rank}(x)=\operatorname{rank}H+\operatorname{rank}G+
\|\lambda\|+\|\pi\|$.
Section 3.4.2 specifies $\|\lambda\|$ as the height for a
$GL_m$ partition, and specifies the combinatorial length of a torus
character as $\dim H$. Therefore the four terms here are
$1,2,1,1$, giving five independently of $N$. A polynomial in this
fixed parameter cannot bound the unbounded modular indices $N$.
This directly contradicts Hypothesis 3.3.4 in the stated broad scope.
\end{proof}

\begin{proposition}[P-F15: fixed-rank positivity and saturation indices are unbounded]
\label{prop:torusindices}
For odd $N\ge3$, let $H=\mathbb G_m$, $G=GL_3$ and
\[
 \sigma_N(t)=\operatorname{diag}(t,t^3,t^{N+1}),
 \qquad \lambda=(1),\qquad\pi=2.
\]
This is again a faithful explicit polynomial embedding. Its stretching
multiplicity $g_N$ satisfies
\[
 g_N(n)=\#\{(v,w)\in\mathbb Z_{\ge0}^2:2v+Nw=n\},
 \qquad
 \sum_{n\ge0}g_N(n)T^n=\frac1{(1-T^2)(1-T^N)}.
\]
In the definitions of GCT VI,
\[
 p(g_N)=s(g_N)=N-2,
 \qquad \operatorname{rank}(x)=6.
 \tag{P-F15}
\]
Its natural graded multiplicity algebra is a polynomial ring on two
generators of degrees $2,N$. Thus the ungraded affine scheme is the
same smooth normal affine plane for every $N$; over finite fields it
has pure Tate ordinary and intersection cohomology, as in P-F10.
Thus Hypothesis 3.3.2(a) (the uniform polynomial saturation-index
bound) and Hypothesis 3.3.3 (PH2) are false for the full subgroup
restriction class with the stated parameter. This finding does not
assert a refutation of the density assertion in Hypothesis 3.3.2(b).
\end{proposition}

\begin{proof}
The three characters have positive weights, so the same exact source
definition as in P-F14 applies. The first diagonal entry is $t$,
giving faithfulness and the inverse on the image. Inside the diagonal
torus, that image is closed by the equations
$z_2=z_1^3$ and $z_3=z_1^{N+1}$.
A monomial $x^u y^v z^w$ of degree $n$ has $u+v+w=n$ and weight
\[
 u+3v+(N+1)w=n+2v+Nw.
\]
It has the required weight $2n$ exactly when $2v+Nw=n$.
Conversely every nonnegative solution of that equation satisfies
$v+w\le n$, because both coefficients are at least two.
It therefore gives the unique admissible value $u=n-v-w$.
This is a bijection with the monomial weight basis and proves the
counting formula and its generating series.
In fact, $0\le w\le\lfloor n/N\rfloor$ and
$w\equiv n\pmod2$, because $N$ is odd. Counting those values gives
the closed formula
\[
 g_N(n)=\left\lfloor
 \frac{\lfloor n/N\rfloor+2-(n\bmod2)}2
 \right\rfloor.
\]
It holds also at $n=0,1$. Integer division and parity computation
evaluate it in time polynomial in the binary lengths of $n,N$.

More precisely, let $S=\mathbb C[x,y,z]$, with ordinary polynomial
degree and torus weights $1,3,N+1$. Its diagonal multiplicity algebra
\[
 T_N=\bigoplus_{n\ge0}(S_n)_{2n}
\]
is closed under multiplication, since degrees and weights both add.
The monomial condition is $u=v+(N-1)w$, so its monomial basis is
\[
 x^{v+(N-1)w}y^v z^w=(xy)^v(x^{N-1}z)^w,
 \qquad v,w\ge0.
\]
The exponent triples for distinct pairs $(v,w)$ are distinct.
Hence these two generators are algebraically independent and generate
all of $T_N$, proving
$T_N=\mathbb C[xy,x^{N-1}z]\simeq\mathbb C[U,V]$ with
$\deg U=2$, $\deg V=N$. The same monomial proof works over every
field. For $k=\mathbb F_q$ and any prime $\ell\nmid q$, P-F10's
affine-space computation therefore gives
$H_c^4(\operatorname{Spec}(T_N)_{\overline k},\mathbb Q_\ell)
=\mathbb Q_\ell(-2)$ and zero in all other degrees, with intersection
complex $\mathbb Q_\ell[2]$. This is an actual graded multiplicity
algebra, rather than only a ring with an equal Hilbert function.

P-F10, already proved in full, applies with $a=2$, $b=N$ since
$\gcd(2,N)=1$, and gives
$p(g_N)=s(g_N)=2N-2-N=N-2$. To make the exact index test explicit,
the constituents can also be read from
\[
 \frac1{(1-T^2)(1-T^N)}
 =\frac{\displaystyle\sum_{i=0}^{N-1}\sum_{j=0}^{1}T^{2i+Nj}}
 {(1-T^{2N})^2}.
\]
The exponents $r=2i+Nj$ represent all residues modulo $2N$ exactly
once: reduction modulo two determines $j$, and reduction modulo $N$
then determines $i$. The constituent belonging to $r$ is
$n\mapsto (n-r)/(2N)+1$. This is valid even when $0\le n<r$,
because the only possible such index in the residue is $r-2N$,
where the formula is zero. The greatest representative is $3N-2$.
After a nonnegative integer shift $c$, the smallest constant term
is $(c-(N-2))/(2N)$. All slopes are positive, so coefficientwise
positivity is equivalent to $c\ge N-2$. For strict saturation
each nonzero constituent must be positive at every integer $n\ge1$;
the smallest value at $n=1$ is $(1+c-(N-2))/(2N)$, which is
strictly positive precisely when $c\ge N-2$. This rechecks both
indices directly with the source conventions.

The four terms in the source's rank parameter are now $1,3,1,1$,
giving six. Since $N-2$ is unbounded, no polynomial in this fixed
parameter bounds either index. These are precisely the universal
bounds asserted in Hypotheses 3.3.2(a) and 3.3.3. A sequence of
counterexamples to a universal bound does not by itself determine
the density appearing in Hypothesis 3.3.2(b).
\end{proof}

\paragraph{Exact source scope and consequences.}
The primary locators are GCT VI, Problem 1.1.3; Hypotheses 3.3.1--3.3.4,
printed pp.~50--51; the definition of $\operatorname{rank}(x)$ in
Section 3.4, printed pp.~52--53; Section 3.4.1, printed pp.~53--54;
and Section 3.4.2, the specifications of $[\rho]$, items (4)--(5),
and of $[\lambda],[\pi]$, printed pp.~55--56.
The source expressly applies Hypotheses 3.3.1--3.3.4 to this
subgroup-restriction input. The maps above use its permitted
direct sums, rather than an irreducible representation of $H$.
Every irreducible representation of the torus is one-dimensional,
so these two- and three-dimensional representations are not inputs
to the irreducible-inner plethysm Problem 1.1.2. No refutation of
Hypotheses 1.6.6--1.6.7 for that narrower family follows.
The conditional implication PH3$\Rightarrow$PH2 in P-F12 is also
unaffected; it does not assert PH3 for these broader inputs.

Both examples even have homogeneous rational-polytope models of
constant ambient dimension. For P-F14 use the singleton
$\{1/N\}\subset\mathbb R$, whose $n$th dilation contains one
integer exactly when $N\mid n$. For P-F15 use
\[
 \{(v,w)\in\mathbb R_{\ge0}^2:2v+Nw=1\}.
\]
Their descriptions have binary length $O(\log N)$, as do the
torus-character inputs. Thus PH1-style efficient descriptions are
compatible with these failures of the index bounds. The indices
in question grow exponentially in this binary length; adding only
the bitlength of $\rho$ to the rank parameter would not make a
polynomial index bound true. A valid replacement must restrict
the family or retain suitable magnitude information about the
embedding weights. These explicit counterexamples do not identify
what occurs for the narrower determinant/permanent obstruction
families, and no historical priority claim is made.
Both examples have easy exact counting formulas---divisibility in
P-F14 and the displayed floor formula in P-F15---so their large
indices do not themselves imply a computational lower bound.

In particular, the broad rank-only bounds cannot follow solely from
purity, smoothness, normality, a positive monomial basis, or an
efficient PH1 polytope: the examples above have all these properties.
Quantitative control of the grading or embedding weights is a
separate requirement even when the graded algebra already realizes
the multiplicities exactly.



% Source: branches/plethysm/density_converse.tex
\section{The density clause and a precise stabilizer converse}
\label{sec:plethysm-density-converse}

\paragraph{The exact narrow question remains open.}
Write $\ell(a)=1+\lfloor\log_2a\rfloor$ for each positive
partition part and add these lengths over the three partitions;
zero parts are omitted.  Let $\mathcal D_N$ be the ordinary
plethysm triples of total length exactly $N$ with nonzero initial
coefficient, and let $\mathcal E_N$ be those whose literal
saturation index is positive.  The conditional reading of
GCT VI, Hypothesis 1.6.5(b), printed p.~25, is
\[
 \text{for every fixed }c>0,\qquad
          |\mathcal E_N|/|\mathcal D_N|\le N^{-c}
                  \quad\text{for all sufficiently large }N.
 \tag{PNEW-density-target}
\]
The new positive initial example proves that this exceptional
set is nonempty; it does not settle its asymptotic density.
The full inherited proof P-F48, reproduced below and rechecked
here, bounds all the cubic rescalings and their inner and outer
determinant twists by an exponentially small proportion.
No map with logarithmic raw-length overhead from arbitrary
ordinary plethysm inputs to such exceptions has been constructed.
Consequently this work neither proves nor refutes
(PNEW-density-target).  This is an explicit unresolved part of
the user request, not an inference from the refutation of SH(a).
PNEW-009 below does prove this exact-shell conclusion after
fixing the inner representation and rank, with constants
depending on that fixed embedding.  Summing over changing
embeddings remains a separate uniformity problem.

\begin{proposition}[PNEW-008: no input-growth converse for the intrinsic graded quotient]
For each integer $b\ge0$ take the ordinary plethysm input
\[
 \mu=(3),\quad
 \lambda_b=(35+b,b^{679}),\quad
 \pi_b=((7+136b)^{15}),
 \tag{PNEW-twisted-positive}
\]
with zero parts omitted.  These inputs have unbounded raw
partition length, initial coefficient one, and strict saturation
index at least seven.  All their canonical graded multiplicity
algebras are isomorphic, compatibly with the grading.
In particular the orders of all stabilizers for the intrinsic
grading action on their nonzero affine quotients are bounded by
one fixed finite constant, independent of $b$.

Thus a positive saturation index, even with positive initial
coefficient, does not force grading stabilizer orders to grow
with the encoded input.  This is a statement about the fixed
effective graded quotient, not about every possible affine
presentation and its raw general-linear determinant character.
\end{proposition}
\begin{proof}
Let $V=\mathbb C^{15}$ and $W=\operatorname{Sym}^3V$.
Then $\dim W=680$ and $\det W=(\det V)^{136}$.
Indeed each variable has equal total exponent in the monomial
basis, and the sum of those exponents is $3\cdot680$.
For every $s\ge0$ the determinant-twist identity gives
\[
 S^{s\lambda_b}W
     =\operatorname{Sym}^{35s}W\otimes(\det W)^{bs}
     =\operatorname{Sym}^{35s}W\otimes(\det V)^{136bs}.
\]
Cancelling the same character from the target
$(\det V)^{(7+136b)s}$ identifies the multiplicity space
with that of $(\det V)^{7s}$ in $\operatorname{Sym}^{35s}W$.
These identifications respect Cartan multiplication: the chosen
one-dimensional character vectors multiply by tensor product
and their common factors cancel degree by degree.  The whole
graded algebra is therefore the single ring with Hilbert
function $h(s)=f_{15}(7s)$ of PNEW-007.  Its initial coefficient
and index bounds follow from that proposition.  The lengths
are unbounded, since the positive parts tend to infinity with $b$.

For the stabilizer assertion, choose homogeneous algebra
generators $u_1,\ldots,u_t$ of this fixed finitely generated
graded ring, and put $D=\max_i\deg u_i$.  At any nonzero
point $y$ of its affine cone, some positive-degree generator
does not vanish: otherwise every positive-degree function
would vanish and $y$ would be the graded origin.  A scalar
$z\in\mathbb C^*$ fixes $y$ precisely when
$z^{\deg u_i}=1$ for all the nonzero values $u_i(y)$.
Its stabilizer is therefore
\[
                  \boldsymbol\mu_{\gcd\{\deg u_i:u_i(y)\ne0\}},
\]
and has order at most $D$.  The bound is uniform over all
points and all $b$, because the ring and its grading are the
same.  It applies in particular to every stable point for
this grading action on the nonzero locus.
\end{proof}

\paragraph{The simpler fixed Kahle--Micha\l ek hole.}
The same construction applied to the fixed ray
$\phi(s)=a_{(4s),(3)}^{(7s,5s)}$ gives
\[
       \lambda_b=(4+b,b,b,b),\qquad
       \pi_b=(7+6b,5+6b),\qquad\mu=(3).
\]
Its stretching algebra is unchanged, since
$\det(\operatorname{Sym}^3\mathbb C^2)=(\det\mathbb C^2)^6$.
The checked inherited identification PLE-F8 realizes it as
$\mathbb C[U,V,Z]/(Z^2-UV^2-4U^4)$ with degrees $(2,3,4)$.
Hence all its nonzero grading stabilizers have order at most
4, while $p(\phi)=s(\phi)=2$ and the input lengths are unbounded.
This simpler example has $\phi(1)=0$, so it is excluded from
the conditioned population in (PNEW-density-target).  PNEW-008
removes that exclusion using an explicitly proved positive initial
coefficient.

\paragraph{Why the determinant-character convention matters.}
The intrinsic grading action is the one detected by Hilbert poles
and by the fixed-locus argument in the plethysm refutation.
It is invariant under the graded-algebra isomorphisms just proved.
A raw $GL$ action on a newly twisted Cartan cone can acquire a
different central kernel, even when its effective graded quotient
is unchanged.  On a preinvariant Cartan tensor product the whole
center can instead act trivially.  These are different stabilizer
questions.  We therefore make no assertion that PNEW-008 bounds
the determinant image in every such raw presentation.  A converse
formulated using that presentation-dependent determinant image
requires an explicit normalization before it has the intrinsic
meaning suggested by Hilbert-series poles.

\paragraph{The broad density clause.}
GCT VI, Hypothesis 3.3.2(b), printed p.~51, concerns specified
subgroup-restriction inputs of length at most $N$.  It allows
ordered products of groups and explicit polynomial maps, with
additive costs as in Sections 3.4.1--3.4.2, printed pp.~53--57.
The inherited P-F63 proof reproduced below, including its
finite-coding assumption, has been reread in full and checked
against those source rules.  It refutes this broad clause under
every finite source-compatible coding that admits constant-cost
adjoining of one fixed ordered factor and has
$\exp(O(N\log(N+2)))$ inputs of source cost at most $N$.
It does not supply an ordinary plethysm padding map, and therefore
does not change the unresolved narrow conclusion above.


% Source: branches/plethysm/sections/positivity_density_scope.tex
% Reproduced inherited proof; finding P-F48.
\subsection{The density clause and the size of the cubic families}
\label{sec:cubicdensityscope}

\begin{proposition}[P-F48: the cubic constructions are exponentially sparse in the stated input population]
\label{prop:cubicdensityscope}
For a positive integer $a$, put
$\ell(a)=1+\lfloor\log_2a\rfloor$. For a partition, sum this
quantity over its positive parts; zero parts are omitted.
Write
\[
 L(\lambda,\mu,\pi)=\ell(\lambda)+\ell(\mu)+\ell(\pi).
\]
This is exactly the input-size convention in GCT VI \cite{PGCTVI},
Section~1.1, printed p.~9. The narrow SH(b), Hypothesis~1.6.5(b),
printed p.~25, concerns triples of total length \emph{exactly}
$N$ with $a_{\lambda,\mu}^{\pi}\ne0$. The unified SH(b),
Hypothesis~3.3.2(b), printed p.~51, instead says length
\emph{at most} $N$ and uses the specified input $x$.
For its canonical ordinary-plethysm specification,
Section~3.4, printed p.~53, that length is
$n+L(\lambda,\mu,\pi)$: the group $GL_n$ contributes $n$.
Thus these are different input populations, not just
different cutoffs. Neither clause imposes a fixed inner
shape, fixed outer shape, fixed target shape, or a fixed rank.

Let $\mathcal D_N$ be the set of partition triples of length
exactly $N$ with nonzero ordinary plethysm coefficient, and
put $\mathcal D_{\le N}=\bigcup_{j\le N}\mathcal D_j$.
The counting fractions for narrow triples, using either cutoff,
and an exceptional set
$\mathcal E$ are
\[
 \frac{|\mathcal E\cap\mathcal D_N|}{|\mathcal D_N|},
 \qquad
 \frac{|\mathcal E\cap\mathcal D_{\le N}|}
      {|\mathcal D_{\le N}|},
 \tag{P-F48a}
\]
The first is the narrow clause's natural
conditional reading; the second is its at-most-$N$ variant.
The source states the density in words rather than displaying
a fraction. If instead that wording is read
as density among all partition triples, the estimates below
continue to hold, since that denominator is larger. No
distribution on partitions of a fixed weight or on a chosen
one-parameter family is substituted for the source's population.

For $n\ge4$, set
\[
 g=\gcd(n,3),\quad m=\binom{n+2}{3},\quad
 d=\frac{rn}{g},\quad k=\frac{3r}{g},
 \qquad r\ge1.
\]
The rescaled cubic triples
\[
                  ((d),(3),(k^n))
 \tag{P-F48b}
\]
form a set $\mathcal C$. Their stretching functions are
$f_n(rs)$, in the notation of P-F31--P-F32.
More generally, for arbitrary integers $q,b\ge0$ put
\[
 \begin{split}
  \mu_{n,q}&=(q+3,q^{\,n-1}),\\
  \lambda_{n,r,b}&=(d+b,b^{\,m-1}),\\
  e_{n,q}&=\frac{m(3+nq)}n
          =\frac{(n+1)(n+2)}2+qm,\\
  \pi_{n,r,q,b}&=((k+qd+be_{n,q})^n),
 \end{split}
 \tag{P-F48c}
\]
again omitting all zero parts. These natural inner and outer
determinant twists form a set $\mathcal C^{\det}$ containing
$\mathcal C$. Every one of them has the same stretching
function $f_n(rs)$ as its untwisted counterpart.

For every integer $N\ge12$,
\[
 \begin{split}
 |\mathcal D_N|&\ge 2^{N/2-4},\\
 |\{x\in\mathcal C:L(x)\le N\}|
       &\le N\,2^{N/5},\\
 |\{x\in\mathcal C^{\det}:L(x)\le N\}|
       &\le N(N+1)^3\,2^{N/4}.
 \end{split}
 \tag{P-F48d}
\]
Consequently the conditional density of \emph{any} subset
of $\mathcal C$ is at most
$16N\,2^{-3N/10}$, and that of any subset of
$\mathcal C^{\det}$ is at most
$16N(N+1)^3\,2^{-N/4}$, for both populations in
(P-F48a). In particular these densities are smaller than
$N^{-c}$ for every fixed $c>0$ and all sufficiently large
$N$.

The same conclusion holds separately in the canonical
ordinary-plethysm input population of the unified clause.
Let $\widehat{\mathcal D}_{\le N}$ count admissible
$(h,\lambda,\mu,\pi)$ with nonzero coefficient and
$h+L(\lambda,\mu,\pi)\le N$. Lift the constructed
families using all their admissible ambient ranks
$h\ge n$, including padding above the displayed rank $n$.
For $N\ge13$ their conditional densities in this population
are at most
\[
 32N^2\,2^{-3N/10},\qquad
 32N^2(N+1)^3\,2^{-N/4},
 \tag{P-F48h}
\]
respectively. This statement uses the canonical
ordinary-plethysm specification; it is not a claim
about counting arbitrarily many syntactic descriptions
of a composite subgroup embedding.

The unscaled family $r=1$, including all its twists in
(P-F48c), has zero initial multiplicity and is excluded
from $\mathcal D_N$ altogether. These facts do not prove
SH(b), do not assert that the rescaled functions have zero
strict saturation index, and do not preclude a different
family of positive density.
\end{proposition}

\begin{proof}
\emph{1. A nonzero population of exponential size in every
large length shell.}
For odd $N=2u+1$, choose any integer
$a\in[2^{u-1},2^u-1]$ and take
\[
                         (\lambda,\mu,\pi)=((1),(a),(a)).
\]
Its coefficient is one, since the outer Schur functor is
the identity. Its length is exactly $1+u+u=N$.
There are $2^{u-1}$ such triples.

For even $N=2u+4$, choose
\[
                    \left\lceil\frac{2^{u+1}}3\right\rceil
                              \le a\le 2^u-1
\]
and take $((3),(a),(3a))$. Here $\ell(a)=u$ and
$\ell(3a)=u+2$, so the total length is $2+u+(u+2)=N$.
The coefficient is one: the highest weight $(3a)$ occurs
once in $\operatorname{Sym}^3(\operatorname{Sym}^a V)$.
Indeed its weight space is spanned by the cube of the
highest weight vector of $\operatorname{Sym}^a V$.
This also gives a direct verification in any sufficient
ambient rank, or just in rank one.
For $u\ge4$, the number of choices is at least
$2^u/3-1\ge2^{u-2}$.
Both parity cases therefore give the first inequality in
(P-F48d) for $N\ge12$. In particular the denominator is
nonzero in every such shell.

\emph{2. The determinant twists preserve the actual ray.}
Let $V=\mathbb C^n$, $W=\operatorname{Sym}^3 V$, and
$W_q=W\otimes(\det V)^q$. Then
\[
              W_q=V_{\mu_{n,q}}(GL_n),\qquad \dim W_q=m.
\]
The determinant character of $W$ is
$(\det V)^{3m/n}$. To see this, each diagonal coordinate
occurs with the same total exponent in its monomial
basis, and the sum of those $n$ exponents is $3m$.
The integer $3m/n=(n+1)(n+2)/2$ proves both the formula
and the asserted integrality of $e_{n,q}$.
Consequently
\[
                    \det W_q=(\det V)^{e_{n,q}}.
\]
For every integer $s\ge0$, the usual determinant twist
identity for a highest weight gives
\[
 \begin{split}
 V_{s\lambda_{n,r,b}}(GL(W_q))
   &\simeq \operatorname{Sym}^{sd}(W_q)
                          \otimes(\det W_q)^{sb}\\
   &\simeq \operatorname{Sym}^{sd}(W)
                    \otimes(\det V)^{s(qd+be_{n,q})}.
 \end{split}
 \tag{P-F48e}
\]
The target $V_{s\pi_{n,r,q,b}}(GL_n)$ is the character
$(\det V)^{s(k+qd+be_{n,q})}$. Cancelling the common
determinant twist reduces its multiplicity to that of
$(\det V)^{sk}$ in $\operatorname{Sym}^{sd}(W)$.
The central degree condition is $3sd=nsk$, so that
this is precisely the $SL_n$-invariant multiplicity
$f_n(rs)$. This proves the ray identity for all $s$,
including the boundary cases $q=0$ and $b=0$.

\emph{3. The bit budgets are literal partition budgets.}
For (P-F48b), $d\ge r$ and $k\ge r$, whence
\[
             L((d),(3),(k^n))
                      \ge 2+(n+1)\ell(r).
 \tag{P-F48f}
\]
The target has $n$ positive parts, so $n\le N$ whenever
the total length is at most $N$. For each such $n$,
the number of possible positive $r$ is at most
$2^{N/(n+1)}\le2^{N/5}$.
Summing this bound over at most $N$ values of $n$
proves the second inequality in (P-F48d).

For the larger family define $\ell(0)=0$ only for the
two parameters $q,b$. This convention counts their
omitted zero parts; it does not add zeros to the source's
partition encoding. Put
$u=\ell(r)$, $v=\ell(q)$ and $w=\ell(b)$.
The first part of $\lambda_{n,r,b}$ is at least $r$;
its other $m-1$ positive parts, when $b>0$, equal $b$.
The $n$ parts of $\mu_{n,q}$ are each at least $q$
when $q>0$. The target has $n$ parts each at least $r$.
These are disjoint entries in the three encoded
partitions, and therefore
\[
              L(\lambda_{n,r,b},\mu_{n,q},\pi_{n,r,q,b})
                   \ge (n+1)u+nv+(m-1)w
                   \ge4(u+v+w).
 \tag{P-F48g}
\]
When $q=0$ the omitted contribution on the right is
zero, while the actual inner partition still has its
positive part $3$; when $b=0$ only the first outer
part remains. Thus (P-F48g) is valid also in those
cases without double counting any entry.

There are at most $(N+1)^3$ possible bit-length triples
$(u,v,w)$. For each one, the number of parameter triples
$(r,q,b)$ is at most $2^{u+v+w}\le2^{N/4}$, using the
single zero option when $v=0$ or $w=0$.
Again $n\le N$ because the target has $n$ positive
parts. Summing gives the third inequality in
(P-F48d). Distinct parameter choices that happen to
give the same input only reduce the count.

\emph{4. The density and nonzero-input conclusions.}
For the exact-length density, bound its numerator by
the corresponding at-most-$N$ count in (P-F48d) and
its denominator by the first line there.
For the at-most-$N$ narrow density use the same numerator
bound and $|\mathcal D_{\le N}|\ge|\mathcal D_N|$.
This gives exactly the two displayed exponential
estimates, even if every member of the constructed
family were an SH(b) exception. Exponential decay
times a fixed polynomial is smaller than every
inverse power of $N$.

For the canonical unified ordinary-plethysm inputs,
each underlying triple has at most $N$ admissible
ambient ranks of cost at most $N$.
The stable plethysm coefficient, including its stretched
versions, is unchanged by such padding: specialization
of the symmetric-function character to $h$ variables
preserves the coefficients of all Schur functions of
height at most $h$, including the target.
Multiplying the previous numerator bounds by $N$
therefore covers every padded realization of the family.
The one-row examples in the first part of the proof
are all admissible already at rank one.
Use those of narrow length $N-1$, then add the rank
cost one. They show
\[
        |\widehat{\mathcal D}_{\le N}|
                  \ge 2^{(N-1)/2-4}
                  \ge 2^{N/2-5}
                         \qquad(N\ge13).
\]
Combining this with the enlarged numerator bounds
proves (P-F48h). Any larger population with these
canonical inputs and the same encoding costs has a
larger denominator; no inference about a different
encoding's numerator is made.

Finally P-F32 applied at $s=1$ gives
$f_n(1)=0$ for every $n\ge4$: its exterior parameter
is $k=3/\gcd(n,3)\le3$, for which
$(k-1)(k-2)/2\le1<n$.
Equation (P-F48e) preserves this zero under all the
determinant twists.
If instead $f_n(r)>0$, powers of a nonzero invariant
show $f_n(rs)>0$ for every positive integer $s$.
This last statement concerns realized values only.
Definition~1.2.4 of GCT VI tests every nonzero
constituent polynomial at every positive integer,
including arguments outside that constituent's
residue class. Realized nonvanishing alone thus
does not establish the strict saturation assertion,
and none was used in the counting proof.
\end{proof}

\paragraph{Scope of the remaining density problem.}
The two source versions use different length cutoffs
and, in their canonical ordinary-plethysm specifications,
different rank costs. The proof explicitly treats
each, so their equivalence is not assumed.
Neither changing the conditioning to
rectangular targets nor counting only the parameter
$r$ would test the stated clause. Nor does padding an
ambient rank create new narrow inputs: the source
counts the three partitions and omits their zero parts.
It does create distinct canonical unified inputs,
and the extra factor $N$ in (P-F48h) explicitly covers
them.
The proved conclusion closes this particular
amplification attempt, including its natural determinant
twists. Establishing or refuting SH(b) for the entire
source population remains open here.


% Source: branches/plethysm/sections/positivity_subgroup_density_padding.tex
% Reproduced inherited proof; finding P-F63.
\paragraph{A constant-size product factor obstructs the broad subgroup density clause.}
The narrow ordinary-plethysm population in P-F48 must be distinguished
from the product-closed subgroup-restriction population in GCT VI,
Problem~1.1.3.  The following argument concerns the latter population.
It uses its cumulative convention $\langle x\rangle\le N$, as stated
in Hypothesis~3.3.2(b), printed p.~51, and the additive product
specifications in \S\S3.4.1--3.4.2, printed pp.~53--57:
\url{https://arxiv.org/pdf/0704.0229}.

\emph{The finite-coding convention.}
Inputs are the specified tuples
$x=([H],[\rho],[\lambda],[\pi])$, rather than isomorphism classes
of representations with all their specifications forgotten.
Products have ordered factor specifications; an ordered expression
tree, or a list of its factors, records the functorial constructions.
These are the ordinary finite specifications referred to by the
source's phrases ``composed'' and ``defined additively.''  They
can be serialized in $O(N\log(N+2))$ binary bits when the source's
additive bit cost is at most $N$.  We prove the required bound below.
The result holds for any such finite source-compatible convention
with constant-cost adjoining of a fixed product factor.  It does
not depend on choosing a particular parenthesis code or a particular
constant in that serialization bound.  The source does not specify
the exact delimiters for composite maps or the compact labels for
every reductive target; the argument requires only these stated
additivity and finite-description properties.  If redundant
zero-cost syntax were instead counted as infinitely many distinct
inputs of a fixed cost, the quoted density would not define a
finite counting problem at all.

Let $A(N)$ count the nonzero subgroup-restriction inputs with source
bit cost at most $N$, and let $E(N)$ count those among them whose
stretching quasipolynomial has nonzero strict saturation index.
There is a fixed integer $C>0$ such that
\[
                    E(N+C)\ge A(N).                 \tag{P-F63a}
\]
Consequently this population cannot satisfy
$E(N)/A(N)\le N^{-c}$ eventually for every $c>0$.
Thus the broad, product-closed version of GCT VI
Hypothesis~3.3.2(b) fails under the finite input-specification
convention just stated.  This conclusion makes no assertion about
the narrow triple population in Hypothesis~1.6.5(b), about a
population restricted to a fixed group or embedding, or about
ordinary Kronecker and ordinary plethysm inputs considered on
their own.

\emph{Proof of the fixed factor.}
Take $H_0=(\mathbb C^*)^2$ and the faithful polynomial map
\[
 \rho_0(u,v)=\operatorname{diag}(uv,u,1,v^4)\in GL_4.
\]
It is faithful because its character weights $(1,1)$ and $(1,0)$
generate the character lattice $\mathbb Z^2$.  All character
coordinates are nonnegative, so this is one of the basic explicit
polynomial representations allowed in \S3.4.1.  Set
\[
          \lambda_0=(4),\qquad \pi_0=(2,5).
                                                        \tag{P-F63b}
\]
The multiplicity $h(s)$ of the character $(u^2v^5)^s$ in
$\operatorname{Sym}^{4s}(\mathbb C^4)$ counts the nonnegative
integer solutions
\[
 x_1+x_2+x_3+x_4=4s,\qquad
 x_1+x_2=2s,\qquad x_1+4x_4=5s.
                                                        \tag{P-F63c}
\]
Every solution has, uniquely, the form
\[
 (x_1,x_2,x_3,x_4)=(s+4j,s-4j,s+j,s-j),
 \qquad j\in\mathbb Z,\quad |j|\le s/4.
                                                        \tag{P-F63d}
\]
Indeed, write $j=s-x_4$ and solve the three displayed equations;
the first two nonnegativity inequalities are $|j|\le s/4$,
and these imply the other two when $s\ge0$.
Monomials are a weight basis, so no further multiplicity factor
occurs.  Hence
\[
 h(s)=2\lfloor s/4\rfloor+1,
 \qquad h_r(X)=\frac X2+1-\frac r2\quad(0\le r<4).
                                                        \tag{P-F63e}
\]
In particular $h(1)=1$ and $h_3(X)=(X-1)/2$ is nonzero but
$h_3(1)=0$.  Under the source's literal definition, a nonzero
constituent must be positive at every positive integer for strict
saturation.  Thus this is already a nonzero input with a strict
saturation exception.  Its saturation and positivity indices
are both exactly one: zero fails at $h_3(1)$, whereas after
shifting by one all constituent coefficients are nonnegative
and every nonzero constituent is positive at positive integers.

\emph{Padding an arbitrary positive input.}
For an input $x=(H,\rho:H\to G,\lambda,\pi)$, form
\[
 \begin{split}
 H^+&=H_0\times H,\qquad G^+=GL_4\times G,\qquad
 \rho^+=\rho_0\times\rho,\\
 \lambda^+&=(\lambda_0,\lambda),\qquad
 \pi^+=(\pi_0,\pi).
 \end{split}                                           \tag{P-F63f}
\]
Problem~1.1.3 permits an arbitrary connected reductive target
$G$, not only a single general linear group.  Thus $G^+$ is
allowed directly; no block embedding, enlarged partition or
extra restriction calculation is being used.  The second
composition rule of \S3.4.1 explicitly permits the product of
maps with separate source groups.  The irreducible modules for
these product groups are external tensor products.  Their Hom
spaces factor, and therefore the stretching function of $x^+$ is
\[
                        f_{x^+}(s)=h(s)f_x(s).          \tag{P-F63g}
\]
For the source's disconnected extension the same statement holds:
the finite-component factor $\operatorname{Sym}^s(W^*)$ in
Theorem~4.1.2 is unchanged, while the new connected torus factor
is an external tensor factor.  This is the extension actually
defined in \S4.1, printed pp.~65--66; no multiplication of a
finite-group highest-weight label is assumed.

We recall why $f_x(1)>0$ implies that every constituent of
$f_x$ is nonzero.  In the connected case the canonical graded
algebra is
\[
 T_x=\bigoplus_{s\ge0}
       (V_{s\lambda}(G)\otimes V_{s\pi}(H)^*)^H
                                                        \tag{P-F63h}
\]
with Cartan multiplication, exactly GCT VI \S4.1.2,
equation~(4.5), printed p.~69.  It is a domain: the two Cartan
algebras are homogeneous coordinate domains of their flag cones,
their tensor product is a domain over $\mathbb C$, and taking
a diagonal graded subalgebra and then invariants preserves this
property.  For the disconnected extension, tensor also with the
polynomial domain $\operatorname{Sym}(W^*)$ and take its matching
degree before invariants; the same argument proves the claim.
A nonzero element of $(T_x)_1$ has nonzero powers in every
positive degree.  Thus $f_x(s)>0$ for every positive integer $s$.
Each residue class contains arbitrarily large such integers,
so none of its constituent polynomials is identically zero.

Choose a period $L$ for $f_x$, and use the common period
$L'=\operatorname{lcm}(L,4)$ for (P-F63g).
For any residue $a$ with $a\equiv3\pmod4$, its constituent is
\[
                \frac{X-1}{2}\,q_{a\bmod L}(X).
                                                        \tag{P-F63i}
\]
It is a nonzero polynomial and vanishes at $X=1$.
It remains a constituent polynomial after passage to the least
period: equality on the infinite corresponding progression
identifies the two polynomials.  Thus $s(f_{x^+})\ge1$,
while $f_{x^+}(1)=f_x(1)>0$.  No sign assumption about the
coefficients or off-residue values of $q_a$ was needed.

\emph{The input count and the density contradiction.}
Keep the original input specification as the second ordered
factor in (P-F63f).  Deleting the first factor recovers it, so
$x\mapsto x^+$ is injective on specifications.  In \S3.4.2
the group, composite-map and product-weight costs are all
additive.  The original specifications are used once, without
copying or expanding their representation dimensions.
Adjoining the fixed specifications (P-F63b) therefore increases
the source cost by a fixed constant $C$.  This proves (P-F63a).
Here a basic constituent retains its existing source specification
and cost when included in a composite.  In particular the special
ordinary-plethysm specification $[\rho]=\mu$ on printed p.~53
is not replaced by a newly expanded generic specification on
adjoining the product factor.  This is precisely the composite
rule on printed p.~56 that adds the costs of its basic
constituents.  A convention that changed the old constituent's
cost by an unbounded amount after adjoining a fixed factor would
not satisfy the constant-additive product convention of this
statement.

Here is a sufficiently generous finite-coding bound; no exact
asymptotic for $A(N)$ is required.  Write all integer/weight
payloads in binary, retaining the source's unary rank data.
Their total charged length is at most $N$.  There are at most
$O(N)$ nonempty payload blocks and positive-cost basic factors.
The source uses a fixed list of group types and the fixed
operations direct sum, tensor product and product of maps.
The composition tree has $O(N)$ nodes after binary bracketing;
listing their types, arities, endpoints and payload lengths
uses $O(N\log(N+2))$ bits.  The same bound covers a list or a
directed acyclic representation with references, provided the
source's additive cost continues to count repeated constituents
by occurrence; references do not introduce free extra operations.
Root-system
and classifying-label information is already part of the
charged payload.  Trivial empty factors do not supply
additional mathematical data; counting arbitrarily many
zero-cost repetitions would instead destroy finiteness of the
source population, as noted above.  Consequently, for some
constant $B>0$ depending only on the finite coding convention,
\[
            A(N)\le\exp\bigl(BN\log(N+2)\bigr).
                                                        \tag{P-F63j}
\]
In a fully delimited bitstring convention the sharper
$2^{O(N)}$ bound is available, but is not needed.

Suppose the claimed density property held.  Fix one real
number $c>BC$, and let $N_0$ be sufficiently large that
$A(N_0)>0$ and
$E(m)/A(m)\le m^{-c}$ for every $m\ge N_0$.
Combining this with (P-F63a) gives
\[
       A(N+C)\ge (N+C)^c A(N)\qquad(N\ge N_0).
                                                        \tag{P-F63k}
\]
Iterate at $N=N_0,N_0+C,\ldots,N_0+(k-1)C$.
Since $N_0+iC\ge iC$, integration of the increasing function
$\log x$ gives
\[
 \begin{split}
 \log A(N_0+kC)
 &\ge\log A(N_0)+c\sum_{i=1}^k\log(N_0+iC)\\
 &\ge c\bigl(k\log C+k\log k-k+1\bigr)+\log A(N_0).
 \end{split}                                           \tag{P-F63l}
\]
Dividing by $(N_0+kC)\log(N_0+kC+2)$ and letting $k$ tend
to infinity gives lower limit at least $c/C>B$,
contradicting (P-F63j).  Equivalently, some fixed polynomial
lower bound for the exceptional proportion holds along an
unbounded sequence of cutoffs.  This suffices to negate the
source's faster-than-every-inverse-polynomial assertion.
\hfill$\square$

The restriction to the broad product-closed population is
essential.  The padding changes the ambient group and embedding,
and does not turn a single ordinary plethysm triple into another
single ordinary plethysm triple.  In particular it does not
overrule the exponential sparsity calculation in P-F48 or the
remaining narrow SH(b) question.  The argument also does not
construct a counting algorithm or an arithmetic-circuit lower
bound; it is a literal constituent and input-population argument.



% Source: branches/plethysm/fixed_rank_density.tex
\section{An exact-shell density theorem at fixed rank and inner partition}
\label{sec:fixed-rank-density}

The density population of SH(b) allows the rank and all three
partitions to vary.  The following restricted theorem uses its
literal sum-of-binary-lengths encoding and its nonzero initial
coefficient conditioning, but fixes the rank and inner partition.
It does not substitute a Euclidean-size population for that encoding.

\begin{theorem}[PNEW-009: fixed proper embeddings have exponentially sparse strict saturation exceptions]
Fix $n\ge2$ and a polynomial highest weight $\mu$ of $GL_n$.
Put $d=|\mu|$ and $m=\dim S_\mu\mathbb C^n$, and suppose
$m>n$.  Among all ordinary plethysm inputs $(\lambda,\mu,\pi)$
with $\operatorname{ht}(\lambda)\le m$, $\operatorname{ht}(\pi)\le n$,
$a_{\lambda,\mu}^{\pi}>0$ and
\[
              \ell(\lambda)+\ell(\mu)+\ell(\pi)=N,
\]
the proportion with $s(a_{s\lambda,\mu}^{s\pi})>0$ is
$O_{n,\mu}(N^{m+n-1}2^{-\eta_{n,\mu}N})$ for some
$\eta_{n,\mu}>0$.  The denominator is nonzero for every
sufficiently large $N$.

In particular the narrow SH(b) conclusion holds on every fixed
rank, fixed cubic-inner-partition population with $n\ge2$.
No uniform estimate on the constants as $n$ or $\mu$ varies is
asserted.
\end{theorem}

We isolate the elementary estimates and the classical branching
input, so the scope of the conclusion can be checked separately.

\begin{lemma}[Polynomial sublevel estimate]
Let $P$ be a nonzero homogeneous polynomial of degree $e\ge1$
on $\mathbb R^r$, and fix a full lattice $\Lambda$ there.
For each fixed $A,C>0$, the number of $u\in\Lambda$ with
\[
       \|u\|\le CR,\qquad |P(u)|\le AR^{e-1}
\]
is $O(R^{r-1/e})$ as $R\to\infty$.
\end{lemma}
\begin{proof}
Choose a primitive lattice vector $v$ with $P(v)\ne0$;
such a vector exists because a nonzero polynomial cannot vanish
on every lattice point.  Complete $v$ to a lattice basis and
write its coordinates as $(y_1,\ldots,y_r)$.  For fixed
$y_2,\ldots,y_r$, the polynomial in $y_1$ has degree $e$
and constant leading coefficient $P(v)$.  Factor it over
$\mathbb C$.  If its absolute value is at most $AR^{e-1}$,
at least one of its $e$ linear factors has absolute value at
most $(A/|P(v)|)^{1/e}R^{1-1/e}$.  Its real arguments then
lie in at most $e$ intervals of that radius, irrespective
of the imaginary parts of the roots.  There are
$O(R^{1-1/e})$ integral choices for $y_1$ and $O(R^{r-1})$
choices for the remaining coordinates.  A fixed lattice basis
change changes only the constants in the norm bound.
\end{proof}

\begin{lemma}[Piecewise quasipolynomials and strict saturation]
Let $C\subset E\simeq\mathbb R^r$ be a full-dimensional
rational polyhedral cone, with lattice $\Lambda$.
Suppose $M:C\cap\Lambda\to\mathbb Z_{\ge0}$ is given by
finitely many quasipolynomials on a finite conic polyhedral
subdivision.  Assume that on every ray the resulting stretching
function is interpreted by these global constituent polynomials.
Let $e_*$ be the maximum of 1 and the degrees of the nonzero
constituent polynomials on full-dimensional chambers.  Then
\[
 \#\{u\in C\cap\Lambda:\|u\|\le R,
                      \ s(M(su))>0\}
                 =O(R^{r-1/e_*}).
 \tag{PNEW-density-Euclidean}
\]
The same estimate holds when one conditions on $M(u)>0$.
\end{lemma}
\begin{proof}
Discard the finitely many chamber walls at a cost of
$O(R^{r-1})$ lattice points.  Fix a full-dimensional chamber
and a common finite-index period lattice for all its
quasipolynomials.  Write its constituent polynomials as
$P_\rho=\sum_{j=0}^{e_\rho}P_{\rho,j}$ with homogeneous
pieces.  If $P_\rho\not\equiv0$, its top piece is nonnegative
on the chamber.  Indeed, for any point $v$ in its interior,
choose lattice points in the residue class $\rho$ within
bounded distance of $tv$.  They remain in the chamber for
large $t$, and
$t^{-e_\rho}P_\rho(tv+O(1))\to P_{\rho,e_\rho}(v)$;
the terms before the limit are nonnegative.  Continuity gives
the assertion on the closure as well.

For $\|u\|\le R$ and $e_\rho\ge1$, the sum of the absolute
values of all lower pieces at $u$ is at most
$A_\rho R^{e_\rho-1}$.  If
\[
        P_{\rho,e_\rho}(u)>A_\rho R^{e_\rho-1},
 \tag{PNEW-leading-dominates}
\]
then $P_\rho(Xu)>0$ for every real $X\ge1$: multiply the
top term by $X^{e_\rho}$, and bound every lower power by
that same power.  A nonzero degree-zero constituent is a
positive constant; a zero polynomial is ignored by the
definition of the strict saturation index.

The constituents of $s\mapsto M(su)$ are among the
polynomials $X\mapsto P_\rho(Xu)$: choose a common scalar
period of the finite quotient lattice and restrict to each
residue of $s$.  Thus all nonzero constituents are positive
at every positive integer whenever
(PNEW-leading-dominates) holds for every applicable nonzero
$P_\rho$.  Each failure lies in a polynomial sublevel set
covered by the preceding lemma.  There are only finitely
many chambers and residue polynomials.  Summing their
bounds and the wall contribution proves the estimate.
Conditioning on $M(u)>0$ can only decrease this numerator.
\end{proof}

\begin{lemma}[From size bounds to literal binary-length shells]
Let $D\ge3$ and
\[
 E=\{(x_1,\ldots,x_D):c_1x_1+\cdots+c_Dx_D=0\},
 \qquad c_i\in\mathbb Z\setminus\{0\},
\]
where coefficients of both signs occur.  Let $\Lambda=E\cap
\mathbb Z^D$, so $r=\dim E=D-1$.  Suppose a nonnegative
piecewise quasipolynomial $M$ on a finite union of rational
cones in $E\cap\mathbb R_{\ge0}^D$ has a full-dimensional
chamber with a nonzero residue polynomial and with strictly
positive coordinates.  Suppose its exceptional lattice
points of norm at most $R$ number $O(R^{r-\epsilon})$,
where $0<\epsilon\le1$.

Allow zero coordinates and omit them when summing binary
lengths.  Then among the points with $M(u)>0$ and
$\sum_i\ell(u_i)=L$, the exceptional proportion is
\[
        O\!\left(L^{D-1}2^{-\epsilon L/(2D(D-1))}\right),
 \tag{PNEW-raw-shell-bound}
\]
provided exceptions on proper coordinate faces are counted
arbitrarily.  The denominator is nonzero for every sufficiently
large $L$.
\end{lemma}
\begin{proof}
\emph{A denominator in every large shell.}
Choose a positive vector $v$ in the interior of a chamber
where a constituent has strictly positive top homogeneous
piece.  Avoid all equalities $v_i=2^k v_j$.
Also avoid dyadic values of its individual coordinates, by
a sufficiently small scalar adjustment if necessary.
Only finitely many of these hyperplanes meet a fixed small
compact positive neighborhood.  None contains $E$ when
$D\ge3$ and every $c_i$ is nonzero, so the choice is possible.
As $t$ increases from 1 to 2, each coordinate of $tv$ crosses
one binary threshold, and these $D$ crossings are distinct.
There are consequently positive open subintervals on which
the sum of its coordinate binary lengths takes $D$
consecutive integer values $C_0,\ldots,C_0+D-1$.
Choose a small open ball $B_j\subset E$ around one point
of each such interval, with the same length pattern and
positive top piece throughout its closure.

On $2^bB_j$ the total length is $Db+C_0+j$.
In a fixed nonzero residue class the constituent is positive
there for all large $b$, uniformly on the ball; its top
piece dominates its lower pieces.  The number of points of
that full-rank lattice coset in $2^bB_j$ is
$\Omega(2^{br})$ (inscribe a fixed parallelepiped and count
its scaled lattice translates).  The $D$ consecutive values
of $C_0+j$ cover every residue modulo $D$.  Therefore the
number of nonzero inputs in every sufficiently large shell
is at least
\[
                         c\,2^{rL/D}
 \tag{PNEW-shell-denominator}
\]
for one fixed $c>0$.

\emph{Unbalanced sizes and omitted coordinates.}
First suppose all $D$ coordinates are positive.  Fix their
lengths $b_1,\ldots,b_D$ with sum $L$, and let
$B=\max b_i$.  The equation for $E$ determines a coordinate
having length $B$ uniquely from all the others.  Thus the
number of choices is at most $2^{L-B}$, whether or not the
determined value is an integer.  There are at most
$D(L+1)^{D-1}$ possible length patterns and choices of a
maximum coordinate.  Put
\[
                      \delta=\frac{\epsilon}{2D(D-1)}.
\]
Those with $B>L/D+\delta L$ therefore number
$O(L^{D-1}2^{rL/D-\delta L})$.

If precisely $D'<D$ coordinates are positive, the same
argument gives $O(L^{D'-1}2^{(1-1/D')L})$, unless all
remaining coefficients have one sign, in which case there
are no such points.  The possible empty or single-coordinate
cases contribute no large shell.  Since
\[
 \frac{D-1}{D}-\frac{D'-1}{D'}\ge\frac1{D(D-1)}\ge\delta,
\]
all proper coordinate faces have the required exponentially
smaller count, even if every such point is exceptional.

\emph{Balanced exceptional points.}
For the remaining full-height inputs,
$\|u\|\le C_D2^{L/D+\delta L}$, so the assumed size bound
gives at most
\[
 O\!\left(2^{(r-\epsilon)(L/D+\delta L)}\right)
       \le O\!\left(2^{rL/D-\epsilon L/(2D)}\right).
\]
Combine the three numerator bounds and divide by
(PNEW-shell-denominator).  This proves
(PNEW-raw-shell-bound) without replacing its length shell
by a norm ball.
\end{proof}

\paragraph{The classical branching input and the full cone.}
The restriction associated to the fixed inner representation
is a fixed homomorphism $SL_n\to SL_m$ with finite kernel.
For the compact branching theorem take $K$ to be the image
of $SU(n)$ in $SU(m)$.  Its weight lattice identifies with
a finite-index sublattice of the $SU(n)$ weight lattice.
Apply the theorem on that image lattice, and extend the
multiplicity function by zero on the other congruence
classes.  This remains a finite piecewise quasipolynomial
description on the full lattice.
The map of Lie algebras is injective: $\mathfrak{sl}_n$ is
simple and the representation is nontrivial.  Its image is
proper since $m>n$.  No nonzero ideal of this image is an
ideal of $\mathfrak{sl}_m$, because the latter is simple.
The standard full-cone lemma for branching therefore gives
a solid cone in the product of the two semisimple weight
spaces, of dimension $(m-1)+(n-1)$.
A precise accessible statement is Baldoni and Vergne (2015),
Section 5.1, Lemma 61.
Their Theorem 63 and Propositions 65--67 supply a finite
conic piecewise quasipolynomial formula for branching
multiplicities, including global stretching constituents.
These are classical branching theorems used as cited inputs;
we do not claim to reprove the quantization theorem from
which the chamber description follows.

For clarity, passing to polynomial $GL$ weights introduces
exactly the degree equation
\[
                      d|\lambda|=|\pi|.
 \tag{PNEW-degree-hyperplane}
\]
Every semisimple branching pair lifts to this hyperplane on
its appropriate finite-index congruence lattice.  Indeed the
central $n$-th roots acting on $S_\mu\mathbb C^n$ have
character $z^d$, which gives the necessary congruence for
the determinant exponent of the target.  Adding one to all
$m$ parts of $\lambda$ adds $e=dm/n$ to all $n$ parts of
$\pi$.  The number $e$ is an integer because
$\det(S_\mu\mathbb C^n)=(\det\mathbb C^n)^{dm/n}$:
the total exponent $dm$ is equally distributed among the
$n$ variables.  This determinant direction is independent
of the semisimple weight coordinates.  Adding a sufficiently
large multiple makes every part positive.  Thus the
polynomial branching cone has dimension $m+n-1$, is full
inside (PNEW-degree-hyperplane), and has an open set of
positive, strictly dominant weights.

\begin{proof}[Proof of PNEW-009]
Apply the piecewise quasipolynomial lemma to the branching
function in the degree hyperplane, with $D=m+n$,
$r=D-1$ and $\epsilon=1/e_*$.  A full-dimensional chamber
has some nonzero residue polynomial because its cone is
the asymptotic support of branching.  The preceding
paragraph ensures a chamber of positive polynomial weights.
Partition dominance imposes linear inequalities only and
does not affect either lattice estimate.

Apply the binary-length lemma with $m$ coefficients equal
to $d$ and $n$ coefficients equal to $-1$.
It uses $L=N-\ell(\mu)$, where $\ell(\mu)$ is a fixed
constant.  Thus one may take
\[
                \eta_{n,\mu}
                    =\frac1{2(m+n)(m+n-1)e_*}
\]
in the displayed polynomial-times-exponential bound,
absorbing the fixed length offset in the implicit constant.
The numerator is precisely the positive-initial inputs
with a nonzero strict saturation index; the denominator is
the positive-initial population of the stated shell.
For $\mu=(3)$ one has $m=\binom{n+2}{3}>n$ for every
$n\ge2$, so the final assertion follows.
\end{proof}

\paragraph{What remains for the original SH(b).}
PNEW-009 eliminates a fixed-embedding density obstruction,
but its degree bound, periods, chamber constants and onset
are allowed to depend on $n$ and $\mu$.  The original
population sums over changing embeddings.  A countable
union of exponentially sparse fixed-parameter families need
not be sparse in that growing population.  The remaining
gate is a uniform estimate strong enough to sum their
exceptional counts, or an explicit family violating such
an estimate.  No such uniform estimate is proved here.



\begin{thebibliography}{9}
\bibitem[Mulmuley(2009)]{PGCTVI} Mulmuley, K. D. (2009).
\emph{Geometric complexity theory VI: The flip via positivity}
(arXiv:0704.0229v4). \url{https://arxiv.org/abs/0704.0229v4}
\bibitem[Kahle and Micha\l ek(2018)]{KM}
Kahle, T., \& Micha\l ek, M. (2018).
Obstructions to combinatorial formulas for plethysm.
\emph{Electronic Journal of Combinatorics, 25}(1), P1.41.
Theorem 5 and Remark 6, p.~4.
\url{https://arxiv.org/abs/1507.07131}
\bibitem[Luna(1973)]{Luna} Luna, D. (1973).
Slices \'etales. \emph{M\'emoires de la Soci\'et\'e Math\'ematique
de France, 33}, 81--105. Section I.1, p.~84; Section III.1,
slice theorem and Remark 1, pp.~97--98.
\url{https://www.numdam.org/item/10.24033/msmf.110.pdf}
\bibitem[Huybrechts(2023)]{Huybrechts} Huybrechts, D. (2023).
\emph{The geometry of cubic hypersurfaces}. Cambridge University Press.
Author manuscript: Chapter 3, Corollary 1.6, pp.~149--150;
Chapter 4, Theorem 4.1 and proof of Corollary 4.2, p.~196.
\url{https://www.math.uni-bonn.de/people/huybrech/Notes.pdf}
\bibitem[McCrory et al.(1994)]{MSV} McCrory, C., Shifrin, T.,
\& Varley, R. (1994). \emph{Siegel modular forms generated by
invariants of cubic hypersurfaces} (arXiv:alg-geom/9403005v3).
Section 3.1.2(i), p.~15; Section 3.2, p.~16.
\url{https://arxiv.org/abs/alg-geom/9403005}
\bibitem[Zheng(2021)]{Zheng} Zheng, Z. (2021).
\emph{On abelian automorphism groups of hypersurfaces}
(arXiv:2004.09008v3). Theorems 4.3 and 4.8.
\url{https://arxiv.org/abs/2004.09008v3}
\bibitem[Baldoni and Vergne(2015)]{BV}
Baldoni, V., \& Vergne, M. (2015).
\emph{Multiplicity of compact group representations and applications
to Kronecker coefficients} (arXiv:1506.02472v1).
Section 5.1, Lemma 61, Theorem 63 and Propositions 65--67.
\url{https://arxiv.org/abs/1506.02472}
\end{thebibliography}
